\documentclass[11pt]{amsart}
\usepackage[margin=1.1in]{geometry}
\usepackage{amsmath,amssymb,amsthm,mathrsfs,mathtools}
\usepackage{enumitem}
\usepackage{microtype}
\usepackage{hyperref}
\usepackage[nameinlink,capitalize]{cleveref}
\hypersetup{colorlinks=true,linkcolor=blue,citecolor=blue,urlcolor=blue}
\allowdisplaybreaks
\pdfstringdefDisableCommands{%
  \def\({}%
  \def\){}%
  \def\tilde#1{#1}%
  \def\mathfrak#1{#1}%
  \def\mathrm#1{#1}%
  \def\mathbb#1{#1}%
  \def\tau{tau}%
  \def\lambda{lambda}%
  \def\infty{infty}%
}

\theoremstyle{plain}
\newtheorem{theorem}{Theorem}[section]
\newtheorem{lemma}[theorem]{Lemma}
\newtheorem{proposition}[theorem]{Proposition}
\newtheorem{corollary}[theorem]{Corollary}
\theoremstyle{definition}
\newtheorem{definition}[theorem]{Definition}
\newtheorem{assumption}[theorem]{Assumption}
\theoremstyle{remark}
\newtheorem{remark}[theorem]{Remark}

\newcommand{\dd}{\mathrm d}
\newcommand{\RR}{\mathbb R}
\newcommand{\C}{\mathbb C}
\newcommand{\supp}{\mathrm{supp}}
\newcommand{\nn}{\nabla}
\newcommand{\tr}{\mathrm{tr}}
\newcommand{\fint}{\mathop{\ooalign{\(\int\)\cr\hidewidth\(-\)\hidewidth\cr}}}

\title[Type II Regularity]{Type II Regularity for Navier--Stokes}
\author{Guillem Duran Ballester}
\email{guillem@fragile.tech}
\date{}

\begin{document}

\begin{abstract}
We establish a local Type II exclusion criterion for the three-dimensional
Navier--Stokes equations near a possible singular point.  From a positive
Caffarelli--Kohn--Nirenberg concentration sequence we select renormalized local
windows around a retained core and analyze the resulting alternatives by
compactness, pressure reconstruction, and local energy estimates.  The argument
combines a compact single-core exclusion criterion, a repaired-gauge
representation for nondegenerate concentration cores, local Calderon--Zygmund
pressure control, a Caccioppoli estimate on compact cylinders, multibubble and
cascade reductions, and scale-collapse cost estimates.  The conclusion is
conditional on the analytic inputs and rigidity hypotheses stated in the paper;
no additional stationary or self-similar classification theorem is used
implicitly.
\end{abstract}

\maketitle

\setcounter{tocdepth}{2}
\tableofcontents


\clearpage
\section{Local Single-Core Criterion}\label{sec:single-core-criterion}
\subsection{Introduction}

We first isolate a local compactness mechanism for Type II concentration in
the three-dimensional Navier--Stokes equations, in the standard class of
suitable weak solutions introduced by Leray and the partial-regularity theory of Scheffer and
Caffarelli--Kohn--Nirenberg
\cite{Leray1934,Scheffer1976,CaffarelliKohnNirenberg1982,Lin1998}.  The
setting is a sequence of
renormalized solutions on fixed parabolic cylinders, centered at a selected
concentration core and written in scale-translation variables.  The result
excludes the case in which one compact core carries a fixed positive
Caffarelli--Kohn--Nirenberg density while the localized dissipation along
selected windows vanishes.  The theorem is conditional on the local
representation and pressure decomposition/stability inputs stated below; the
compact-cylinder \(A,E,C,D\) bounds are part of the branch hypothesis.

The notation and estimates follow the usual local energy, critical-space, and
monograph conventions for Navier--Stokes
\cite{Ladyzhenskaya1969,Temam1977,ConstantinFoias1988,Kato1984,EscauriazaSereginSverak2003,LemarieRieusset2002}.
The argument is local.  Uniform Caffarelli--Kohn--Nirenberg bounds and bounded
modulation coefficients give compactness on smaller cylinders.  Vanishing
localized dissipation forces the compactness limit to be spatially constant.
If the constant is zero, strong velocity convergence and pressure stability
contradict the positive core density.  If the constant is nonzero, the sequence
has a nonzero bounded selected-window limit, which is incompatible with the
local Type II condition.

\subsection{Setup and definitions}

For \(R>0\), write
\[
  Q_R=B_R\times(-R^2,0).
\]
For a suitable pair \((v,q)\) on \(Q_R\), define
\[
  A_v(R)=R^{-1}\operatorname*{ess\,sup}_{-R^2<s<0}
  \int_{B_R}|v(y,s)|^2\,dy,
\]
\[
  E_v(R)=R^{-1}\int_{Q_R}|\nabla v|^2\,dy\,ds,
\]
\[
  C_v(R)=R^{-2}\int_{Q_R}|v|^3\,dy\,ds,
\]
\[
  D_q(R)=R^{-2}\int_{-R^2}^0\int_{B_R}
  |q(y,s)-(q)_{B_R}(s)|^{3/2}\,dy\,ds.
\]
All pressure quantities are understood modulo functions of time.

\begin{definition}[Bounded selected-window limit]\label{paper1:def:bounded-window-limit}
Let \((V_n,P_n)\) be a sequence of represented local rescalings on a compact
cylinder \(Q_r\).  A selected-window subsequence has a bounded selected-window limit
on \(Q_r\) if, after passing to a subsequence,
\[
  V_n\to V_*
  \qquad\text{strongly in }L^3(Q_r),
\]
where \(V_*\in L^\infty(Q_r)\).  The limit is nonzero if
\(V_*\not\equiv0\) on \(Q_r\).
\end{definition}

\begin{definition}[Local Type II sequence]\label{paper1:def:typeII}
A represented sequence is called a local Type II sequence if it has positive
local Caffarelli--Kohn--Nirenberg concentration on some fixed compact cylinder
and no selected-window subsequence has a nonzero bounded selected-window limit on
any compact cylinder in the sense of \cref{paper1:def:bounded-window-limit}.
\end{definition}

\begin{definition}[Compact single-core vanishing-dissipation branch]\label{paper1:def:single-core-branch}
Fix radii \(0<r_*<R_*\), constants \(M<\infty\), \(\eta_0>0\), and a sequence
of represented local rescalings \((V_n,P_n,a_n,b_n)\) on \(Q_{R_*}\).  The
sequence is a compact single-core vanishing-dissipation branch if:
\begin{enumerate}[label=\textup{(\roman*)}]
\item each \((V_n,P_n)\) is suitable on \(Q_{R_*}\) and solves the represented
      Navier--Stokes equation
      \begin{equation}\label{paper1:eq:represented-ns}
      \partial_s V_n+(V_n\cdot\nabla)V_n+\nabla P_n
      =\nu\Delta V_n+a_n(s)(V_n+y\cdot\nabla V_n)+b_n(s)\cdot\nabla V_n,
      \qquad \nabla\cdot V_n=0;
      \end{equation}
\item the compact-cylinder suitable estimates are uniformly bounded:
      \[
      A_{V_n}(R_*)+E_{V_n}(R_*)+C_{V_n}(R_*)+D_{P_n}(R_*)\le M;
      \]
\item the selected core has positive CKN concentration on \(Q_{r_*}\):
      \[
      C_{V_n}(r_*)+D_{P_n}(r_*)\ge\eta_0
      \qquad\text{for all }n;
      \]
\item the localized scale and center are nondegenerate for the selected core,
      so the representation theorem gives the variables in
      \eqref{paper1:eq:represented-ns};
\item the modulation coefficients are uniformly bounded on the selected cylinder:
      \[
      \|a_n\|_{L^\infty(-R_*^2,0)}+\|b_n\|_{L^\infty(-R_*^2,0)}\le M;
      \]
\item the branch has vanishing localized dissipation in the sense of
      \cref{paper1:def:localized-dissipation} below.
\end{enumerate}
\end{definition}

\begin{definition}[Vanishing localized dissipation]\label{paper1:def:localized-dissipation}
Let \(r_*<\rho<R_*\).  Selected windows are first translated and rescaled to the
fixed cylinders \(Q_{r_*}\Subset Q_\rho\Subset Q_{R_*}\).  The localized
dissipation on \(Q_\rho\) is
\[
  \mathcal D_n(\rho)=\int_{Q_\rho}|\nabla V_n|^2\,dy\,ds.
\]
The branch has vanishing localized dissipation if, after passing to a
selected-window subsequence, there is \(\rho\in(r_*,R_*)\) such that
\[
  \mathcal D_n(\rho)\to0.
\]
\end{definition}

\subsection{Local analytic inputs}

The compact single-core argument is based on the following local analytic statements.

\begin{theorem}[Local concentration dichotomy]
\label{paper1:thm:paper0-dichotomy}
Let \((u,p)\) be a suitable weak solution on
\(\mathbb R^3\times(T-\delta,T)\), and suppose \(\Sigma(T)\ne\emptyset\).
Then, using the pressure-normalized quantities \(C,D\), there exist
\(x_0\in\Sigma(T)\), \(\eta>0\), and \(r_n\downarrow0\) such that
\[
  C((x_0,T),r_n)+D((x_0,T),r_n)\ge\eta
  \qquad\text{for all }n.
\]
Moreover, at the concentration point either a local Type I scale-invariant
bound holds, or the point lies in the local Type II alternative, meaning that
no such local Type I bound holds.
\end{theorem}

\begin{proof}
The statement is the standing local concentration dichotomy assumed in the sequel.
\end{proof}

\begin{proposition}[Bridge to local Type II sequences]
\label{paper1:prop:paper0-typeII-bridge}
Let \((x_0,T)\) be a concentration point in the local Type II alternative of
\cref{paper1:thm:paper0-dichotomy}.  After choosing the positive concentration scales
given there and applying the local scale-translation selection used in the
present section to a nondegenerate compact single core, the represented selected
windows form a sequence with positive local CKN concentration on a fixed compact
cylinder.  In the selected-window Type II branch, no subsequence has a nonzero
bounded selected-window limit on any compact cylinder.  Hence the represented
selected windows are a local Type II sequence in the sense of \cref{paper1:def:typeII}.
\end{proposition}

\begin{proof}
The local concentration dichotomy supplies the positive CKN concentration sequence and the exclusion of
the local Type I alternative at the chosen concentration point.  The
scale-translation selection rewrites those shrinking cylinders as fixed
selected windows.  The last condition is the selected-window
formulation of the Type II branch considered here: any branch with a nonzero
bounded selected-window limit belongs to the bounded-window alternative,
not to \cref{paper1:def:typeII}.
\end{proof}

\begin{proposition}[Local representation input]\label{paper1:prop:representation-input}
For a nondegenerate selected single core, the local scale-translation
construction gives represented variables \((V_n,P_n,a_n,b_n)\) satisfying
\eqref{paper1:eq:represented-ns} on compact cylinders.  Failure of the localized
nondegeneracy condition is a multibubble, cascade, or gauge-degenerate
alternative, not a compact single-core branch.
\end{proposition}

\begin{proposition}[Local pressure decomposition and stability]
\label{paper1:prop:pressure-reconstruction}
For every \(B_r\Subset B_R\), the represented pressure decomposes as
\[
  P_n=P_{{\rm loc},n}+H_n,
  \qquad
  -\Delta P_{{\rm loc},n}=\partial_i\partial_j(\zeta V_{n,i}V_{n,j}),
\]
where \(\zeta\in C_c^\infty(B_R)\) equals one on \(B_r\), and \(H_n\) is
harmonic in \(B_r\).  The local part satisfies the Calderon--Zygmund bound
\cite{CalderonZygmund1952,Stein1970}
\[
  \|P_{{\rm loc},n}\|_{L^{3/2}(B_R)}
  \le C\|V_n\|_{L^3(B_R)}^2
\]
for a.e. time, and the harmonic part is controlled modulo spatial means by the
local pressure norm.

Moreover, if \(V_n\to0\) strongly in \(L^3(Q_R)\), then for every \(r<R\),
\[
  P_n-(P_n)_{B_r}\to0
  \quad\text{strongly in }L^{3/2}(Q_r).
\]
\end{proposition}

\begin{lemma}[Localization of pressure oscillation bounds]\label{paper1:lem:D-localization}
Let \(0<r<R\), and let \(q\in L^{3/2}(Q_R)\).  Then
\[
  D_q(r)\le C\left(\frac{R}{r}\right)^2D_q(R),
\]
where \(C\) is universal.
\end{lemma}

\begin{proof}
For a.e. \(s\in(-r^2,0)\), Jensen's inequality gives
\[
  |(q)_{B_r}(s)-(q)_{B_R}(s)|^{3/2}
  \le \frac1{|B_r|}\int_{B_r}|q(y,s)-(q)_{B_R}(s)|^{3/2}\,dy .
\]
Hence
\[
  \int_{B_r}|q-(q)_{B_r}(s)|^{3/2}\,dy
  \le
  C\int_{B_r}|q-(q)_{B_R}(s)|^{3/2}\,dy .
\]
Integrating over \((-r^2,0)\) and using \(Q_r\subset Q_R\),
\[
\begin{aligned}
  D_q(r)
  &\le
  Cr^{-2}\int_{-r^2}^0\int_{B_r}
  |q-(q)_{B_R}(s)|^{3/2}\,dy\,ds \\
  &\le
  Cr^{-2}\int_{-R^2}^0\int_{B_R}
  |q-(q)_{B_R}(s)|^{3/2}\,dy\,ds
  =
  C\left(\frac{R}{r}\right)^2D_q(R).
\end{aligned}
\]
\end{proof}

\subsection{Compactness and pressure stability}

\begin{proposition}[Compactness on selected windows]\label{paper1:prop:compactness}
Let \((V_n,P_n)\) be suitable on \(Q_{R_*}\), solve \eqref{paper1:eq:represented-ns}, and satisfy
\[
  A_{V_n}(R_*)+E_{V_n}(R_*)+C_{V_n}(R_*)+D_{P_n}(R_*)\le M,
\]
with
\[
  \|a_n\|_{L^\infty(-R_*^2,0)}+\|b_n\|_{L^\infty(-R_*^2,0)}\le M.
\]
Then for every \(r<R_*\), after passing to a subsequence there are \((V,P)\)
such that
\[
  V_n\stackrel{*}{\rightharpoonup}V
  \quad\text{in }L^\infty(-r^2,0;L^2(B_r)),
\]
\[
  V_n\rightharpoonup V
  \quad\text{in }L^2(-r^2,0;H^1(B_r)),
\]
\[
  V_n\to V
  \quad\text{strongly in }L^3(Q_r),
\]
and
\[
  P_n-(P_n)_{B_r}\rightharpoonup P
  \quad\text{in }L^{3/2}(Q_r).
\]
After passing to a further subsequence, \(a_n\) and \(b_n\) converge weak-* in
\(L^\infty\) to bounded coefficients, and the limit satisfies the corresponding
represented equation distributionally on smaller cylinders.
\end{proposition}

\begin{proof}
The bounds on \(A\) and \(E\) give weak-* compactness in \(L^\infty L^2\) and
weak compactness in \(L^2H^1\) on smaller cylinders.  Fix
\(r<r_1<R_*\), and choose \(m\) so large that
\[
  H^m_0(B_{r_1})\hookrightarrow W^{1,\infty}_0(B_{r_1}).
\]
We claim that \(\partial_sV_n\) is uniformly bounded in
\[
  L^1(-r_1^2,0;H^{-m}(B_{r_1})).
\]
Let \(\varphi\in H^m_0(B_{r_1})\).  Pairing \eqref{paper1:eq:represented-ns} with
\(\varphi\), integrating by parts in space where needed, and using the
embedding above gives
\[
  |\langle (V_n\cdot\nabla)V_n,\varphi\rangle|
  =
  \left|\int_{B_{r_1}}V_{n,i}V_{n,j}\partial_j\varphi_i\,dy\right|
  \le
  C\|V_n\|_{L^3(B_{r_1})}^2\|\varphi\|_{H^m},
\]
\[
  |\langle \nu\Delta V_n,\varphi\rangle|
  \le
  C\|\nabla V_n\|_{L^2(B_{r_1})}\|\varphi\|_{H^m},
\]
\[
  |\langle \nabla P_n,\varphi\rangle|
  =
  \left|\int_{B_{r_1}}(P_n-(P_n)_{B_{r_1}})\nabla\cdot\varphi\,dy\right|
  \le
  C\|P_n-(P_n)_{B_{r_1}}\|_{L^{3/2}(B_{r_1})}\|\varphi\|_{H^m}.
\]
The modulation terms obey
\[
  |\langle a_nV_n,\varphi\rangle|
  \le
  C\|V_n\|_{L^2(B_{r_1})}\|\varphi\|_{H^m},
\]
\[
  |\langle a_n\,y\cdot\nabla V_n,\varphi\rangle|
  =
  \left|\int_{B_{r_1}}a_n V_{n,i}\partial_j(y_j\varphi_i)\,dy\right|
  \le
  C\|V_n\|_{L^2(B_{r_1})}\|\varphi\|_{H^m},
\]
and
\[
  |\langle b_n\cdot\nabla V_n,\varphi\rangle|
  =
  \left|\int_{B_{r_1}}V_{n,i}b_{n,j}\partial_j\varphi_i\,dy\right|
  \le
  C\|V_n\|_{L^2(B_{r_1})}\|\varphi\|_{H^m},
\]
where the constant uses \(r_1\) and the uniform \(L^\infty\) bounds for
\(a_n,b_n\).  Integrating these estimates in time gives a uniform bound because
\[
  \int_{-r_1^2}^0\|V_n(s)\|_{L^3(B_{r_1})}^2\,ds
  \le
  r_1^{2/3}\|V_n\|_{L^3(Q_{r_1})}^2,
\]
\[
  \int_{-r_1^2}^0\|\nabla V_n(s)\|_{L^2(B_{r_1})}\,ds
  \le
  r_1\|\nabla V_n\|_{L^2(Q_{r_1})},
\]
\[
  \int_{-r_1^2}^0
  \|P_n(s)-(P_n)_{B_{r_1}}(s)\|_{L^{3/2}(B_{r_1})}\,ds
  \le
  r_1^{2/3}
  \|P_n-(P_n)_{B_{r_1}}\|_{L^{3/2}(Q_{r_1})},
\]
and
\[
  \int_{-r_1^2}^0\|V_n(s)\|_{L^2(B_{r_1})}\,ds
  \le
  r_1^2
  \|V_n\|_{L^\infty(-r_1^2,0;L^2(B_{r_1}))}.
\]
The uniform \(C\), \(E\), \(D\), and \(A\) bounds therefore prove the claimed
time-derivative bound; the \(D\)-bound on \(Q_{r_1}\) follows from the bound on
\(Q_{R_*}\) by \cref{paper1:lem:D-localization}.

The compact embedding \(H^1(B_{r_1})\Subset L^2(B_r)\), together with
Aubin--Lions compactness \cite{Aubin1963,LionsMagenes1972,Simon1987}, gives
\[
  V_n\to V\qquad\text{strongly in }L^2(Q_r)
\]
after passing to a subsequence.  The same \(L^\infty L^2\cap L^2H^1\) bounds
give the parabolic interpolation estimate
\[
  \|V_n\|_{L^{10/3}(Q_r)}\le C(r,M).
\]
Indeed, for a.e. \(s\), the Sobolev and Gagliardo--Nirenberg inequalities give
\[
  \|V_n(s)\|_{L^{10/3}(B_r)}
  \le
  C\|V_n(s)\|_{L^2(B_r)}^{2/5}
  \|\nabla V_n(s)\|_{L^2(B_r)}^{3/5}.
\]
Raising to the power \(10/3\), integrating in time, and using the \(A\) and
\(E\) bounds yields
\[
  \int_{-r^2}^0\|V_n(s)\|_{L^{10/3}(B_r)}^{10/3}\,ds
  \le
  C
  \left(\operatorname*{ess\,sup}_{-r^2<s<0}
  \|V_n(s)\|_{L^2(B_r)}^2\right)^{2/3}
  \int_{Q_r}|\nabla V_n|^2\,dy\,ds
  \le C(r,M).
\]
The same estimate applies to \(V\) by weak lower semicontinuity, and therefore
to \(V_n-V\) by the triangle inequality.  Interpolating between \(L^2\) and
\(L^{10/3}\), one obtains
\[
  \|V_n-V\|_{L^3(Q_r)}
  \le
  \|V_n-V\|_{L^2(Q_r)}^{1/6}
  \|V_n-V\|_{L^{10/3}(Q_r)}^{5/6}.
\]
The second factor is uniformly bounded and the first tends to zero, so
\[
  V_n\to V\qquad\text{strongly in }L^3(Q_r).
\]

The pressure compactness follows from the uniform \(D\)-bound:
after subtracting spatial means,
\[
  P_n-(P_n)_{B_r}
  \rightharpoonup P
  \qquad\text{in }L^{3/2}(Q_r).
\]
We pass to the limit distributionally using the preceding convergences.  It
remains to justify the modulation terms.  For every compactly supported
test field \(\varphi\),
\[
  \int a_n(s)\, y\cdot\nabla V_n\cdot\varphi
  =
  -\int a_n(s)\,V_{n,i}\,\partial_j(y_j\varphi_i),
\]
where summation over repeated spatial indices is understood.  Equivalently,
\[
  \partial_j(y_j\varphi_i)=3\varphi_i+y_j\partial_j\varphi_i .
\]
Hence this term passes to the limit by strong local \(L^2\) convergence of
\(V_n\) and weak-* convergence of \(a_n\): if \(F_n\to F\) in \(L^1\) and
\(a_n\rightharpoonup^* a\) in \(L^\infty\), then
\[
  \int a_nF_n-\int aF
  =
  \int a_n(F_n-F)+\int (a_n-a)F\to0.
\]
The term \(a_nV_n\) is immediate from the same argument, while
\[
  \int b_n(s)\cdot\nabla V_n\cdot\varphi
  =
  -\int V_{n,i}\,b_{n,j}(s)\,\partial_j\varphi_i
\]
passes to the limit by strong local \(L^2\) convergence of \(V_n\) and weak-*
convergence of \(b_n\).  These passages give the limiting represented equation
in the sense of distributions.
\end{proof}

\begin{remark}[Use of bounded modulation]\label{paper1:rem:bounded-modulation}
The uniform \(L^\infty\) bounds on \(a_n\) and \(b_n\) enter the compactness
argument in exactly two places.  First, they give the negative Sobolev
time-derivative bound for the three lower-order modulation terms in
\eqref{paper1:eq:represented-ns}.  Second, after subsequential weak-* convergence of
\(a_n,b_n\), they allow the modulation terms to pass to the distributional
limit once \(V_n\to V\) strongly in \(L^2_{\rm loc}\).
\end{remark}

\begin{lemma}[Stability of local pressure decomposition]\label{paper1:lem:pressure-stability}
Let \(V_n\to V\) strongly in \(L^3(Q_R)\), and use the pressure decomposition
of \cref{paper1:prop:pressure-reconstruction}.  Then for every \(r<R\),
\[
  P_{{\rm loc},n}\to P_{\rm loc}
  \quad\text{strongly in }L^{3/2}(Q_r),
\]
where \(-\Delta P_{\rm loc}=\partial_i\partial_j(\zeta V_iV_j)\).  If
\(V\equiv0\) on \(Q_R\), then
\[
  P_n-(P_n)_{B_r}\to0
  \quad\text{strongly in }L^{3/2}(Q_r).
\]
More locally in time, if \(V_n\to0\) strongly in
\[
  L^3(B_R\times(-r^2,0)),
\]
then
\[
  P_n-(P_n)_{B_r}\to0
  \quad\text{strongly in }L^{3/2}(Q_r).
\]
\end{lemma}

\begin{proof}
Strong convergence in \(L^3\) gives
\[
  V_{n,i}V_{n,j}\to V_iV_j
  \quad\text{strongly in }L^{3/2}(Q_R).
\]
Indeed,
\[
  \|V_{n,i}V_{n,j}-V_iV_j\|_{L^{3/2}(Q_R)}
  \le
  \bigl(\|V_n\|_{L^3(Q_R)}+\|V\|_{L^3(Q_R)}\bigr)
  \|V_n-V\|_{L^3(Q_R)}.
\]
Calderon--Zygmund estimates for the compactly supported localized source yield
strong convergence of the localized pressures in \(L^{3/2}\) on smaller balls:
\[
  \|P_{{\rm loc},n}-P_{\rm loc}\|_{L^{3/2}(Q_r)}
  \le
  C_{r,R}
  \|V_{n,i}V_{n,j}-V_iV_j\|_{L^{3/2}(Q_R)}.
\]

Assume now that \(V\equiv0\) on \(Q_R\).  Then \(P_{\rm loc}=0\).  By the
zero-limit stability part of \cref{paper1:prop:pressure-reconstruction},
\[
  P_n-(P_n)_{B_r}\to0
  \quad\text{strongly in }L^{3/2}(Q_r).
\]
The same proof is local in time.  If \(V_n\to0\) strongly in
\(L^3(B_R\times(-r^2,0))\), the Calderon--Zygmund estimate and harmonic
pressure control are applied only for times \(s\in(-r^2,0)\), giving the
strong convergence on \(Q_r\).
\end{proof}

\begin{lemma}[Zero-dissipation limit]\label{paper1:lem:zero-dissipation}
If
\[
  V_n\rightharpoonup V\quad\text{in }L^2(I;H^1(B_\rho))
\]
and
\[
  \int_I\int_{B_\rho}|\nabla V_n|^2\,dy\,ds\to0,
\]
then \(\nabla V=0\) a.e. on \(B_\rho\times I\).  Hence for a.e. time,
\(V(\cdot,s)\) is spatially constant on \(B_\rho\).
\end{lemma}

\begin{proof}
Weak lower semicontinuity gives
\[
  \int_I\int_{B_\rho}|\nabla V|^2\,dy\,ds
  \le \liminf_{n\to\infty}\int_I\int_{B_\rho}|\nabla V_n|^2\,dy\,ds=0.
\]
Thus \(\nabla V=0\) a.e. on \(B_\rho\times I\).  By Fubini, for a.e.
\(s\in I\) one has \(\nabla V(\cdot,s)=0\) in \(L^2(B_\rho)\).  Since
\(B_\rho\) is connected, every \(H^1(B_\rho)\) function with zero weak
gradient is a.e. equal to a constant.  Hence \(V(\cdot,s)\) is spatially
constant for a.e. \(s\).
\end{proof}

\begin{lemma}[Spatially constant limits]\label{paper1:lem:constant-limit}
Let \(V\) be a limit on \(Q_\rho\) and suppose \(\nabla V=0\) a.e. there.
Then \(V(y,s)=c(s)\) on \(Q_\rho\), with
\(c\in L^\infty(-\rho^2,0)\).  If \(c\not\equiv0\) on \(Q_\rho\), then
the approximating selected-window subsequence has a nonzero bounded
selected-window limit.
\end{lemma}

\begin{proof}
The spatial constancy follows from \cref{paper1:lem:zero-dissipation}.  Since
\(V(\cdot,s)=c(s)\) on \(B_\rho\), the local \(L^\infty_sL^2_y\) bound gives
\[
  |c(s)|^2|B_\rho|\le \operatorname*{ess\,sup}_{-\rho^2<s<0}
  \int_{B_\rho}|V(y,s)|^2\,dy,
\]
so \(c\in L^\infty(-\rho^2,0)\).  If \(c\not\equiv0\) on \(Q_\rho\), then the
strong \(L^3\) convergence of \cref{paper1:prop:compactness} gives a subsequential
limit to a nonzero bounded function on \(Q_\rho\), which is precisely a nonzero
bounded selected-window limit in the sense of \cref{paper1:def:bounded-window-limit}.
\end{proof}

\subsection{Contradiction lemmas}

\begin{lemma}[Persistence of positive core concentration]\label{paper1:lem:persistence}
For a compact single-core branch in the sense of \cref{paper1:def:single-core-branch},
there are fixed \(r_*>0\) and \(\eta_0>0\) such that along the selected-window
subsequence,
\[
  C_{V_n}(r_*)+D_{P_n}(r_*)\ge\eta_0
  \qquad\text{for all }n.
\]
\end{lemma}

\begin{proof}
The lower bound is part of the definition of the selected compact core and is preserved by
passing to subsequences of the selected windows.
\end{proof}

\begin{lemma}[Zero constant contradicts positive CKN concentration]\label{paper1:lem:zero-contradiction}
Let \(r_*<\rho\).  Suppose the compactness limit in \cref{paper1:prop:compactness} is
spatially constant on \(Q_\rho\) and is identically zero on \(Q_{r_*}\).  Then
\[
  C_{V_n}(r_*)+D_{P_n}(r_*)\to0.
\]
Consequently this case contradicts \cref{paper1:lem:persistence}.
\end{lemma}

\begin{proof}
Strong \(L^3(Q_{r_*})\) convergence to zero gives
\[
  C_{V_n}(r_*)
  =
  r_*^{-2}\int_{Q_{r_*}}|V_n|^3\,dy\,ds
  =
  r_*^{-2}\|V_n\|_{L^3(Q_{r_*})}^3
  \to0.
\]
It remains to control the pressure term.  Since \(V\) is spatially constant on
\(Q_\rho\), the identity \(V=0\) on \(Q_{r_*}\) implies
\[
  V=0
  \qquad\text{on }B_\rho\times(-r_*^2,0).
\]
Indeed, for a.e. \(s\in(-r_*^2,0)\), the spatially constant value on
\(B_\rho\) vanishes on the nonempty subball \(B_{r_*}\), hence vanishes on all
of \(B_\rho\).  The strong \(L^3_{\rm loc}\) convergence therefore gives
\[
  V_n\to0
  \qquad\text{strongly in }L^3(B_\rho\times(-r_*^2,0)).
\]
Applying \cref{paper1:lem:pressure-stability} with \(R=\rho\) and \(r=r_*\), in its
local-in-time form, gives
\[
  P_n-(P_n)_{B_{r_*}}\to0
  \qquad\text{strongly in }L^{3/2}(Q_{r_*}).
\]
Therefore
\[
  D_{P_n}(r_*)
  =
  r_*^{-2}
  \|P_n-(P_n)_{B_{r_*}}\|_{L^{3/2}(Q_{r_*})}^{3/2}
  \to0.
\]
Thus
\[
  C_{V_n}(r_*)+D_{P_n}(r_*)\to0,
\]
contradicting the uniform lower bound
\[
  C_{V_n}(r_*)+D_{P_n}(r_*)\ge\eta_0.
\]
\end{proof}

\begin{lemma}[Nonzero constant contradicts the Type II regime]\label{paper1:lem:nonzero-contradiction}
If the compactness limit is spatially constant and nonzero on a compact
subcylinder, then the original sequence is not a local Type II sequence.
\end{lemma}

\begin{proof}
By \cref{paper1:lem:constant-limit}, the selected-window subsequence has a nonzero
bounded selected-window limit.  This contradicts \cref{paper1:def:typeII}.
\end{proof}

\subsection{Main theorem}

\begin{theorem}[Local single-core Type II criterion]\label{paper1:thm:single-core-criterion}
Under the analytic inputs
\cref{paper1:prop:paper0-typeII-bridge,paper1:prop:representation-input,paper1:prop:pressure-reconstruction},
a compact single-core vanishing-dissipation branch in the sense of
\cref{paper1:def:single-core-branch} cannot be a local Type II sequence.
\end{theorem}

\begin{proof}
By the definition of vanishing localized dissipation, after passing to a
selected-window subsequence there is a fixed radius \(\rho\in(r_*,R_*)\) such
that
\[
  \int_{Q_\rho}|\nabla V_n|^2\,dy\,ds\to0
\]
and the positive core bound on \(Q_{r_*}\) remains valid along this subsequence.
Apply compactness on the fixed compact
cylinder \(Q_\rho\) and pass to a subsequence.  The compactness proposition
gives
\[
  V_n\rightharpoonup V\quad\text{in }L^2(-\rho^2,0;H^1(B_\rho))
\]
and strong convergence in \(L^3\) on every smaller cylinder.  Since the
localized dissipation tends to zero, the zero-dissipation lemma gives
\[
  \nabla V=0\quad\text{a.e. on }Q_\rho.
\]
Thus \(V(y,s)=c(s)\) on \(Q_\rho\), and in particular on the selected core
cylinder \(Q_{r_*}\).

If this spatially constant limit is zero on the core cylinder, the zero-limit
contradiction lemma gives
\[
  C_{V_n}(r_*)+D_{P_n}(r_*)\to0,
\]
contradicting persistence of positive core concentration.

If the spatially constant limit is not identically zero on \(Q_{r_*}\), the
nonzero-limit lemma produces a nonzero bounded selected-window limit, contradicting
the definition of a local Type II sequence.  The two cases exhaust all
spatially constant limits, so the branch cannot occur.
\end{proof}

\begin{remark}[Minimal local inputs]
The proof uses two local analytic inputs from the Type II argument: the
localized representation for a nondegenerate selected core, the compact-ball
pressure reconstruction with stability on smaller balls.  The
compact-cylinder \(A,E,C,D\) bounds and vanishing localized dissipation are
part of the branch hypothesis, not separate analytic inputs.  Multibubble,
cascade, gauge-degenerate, and compactness-failure alternatives are outside
this single-core theorem; they are treated by the local multibubble and
cascade exclusion theorem.
\end{remark}

\clearpage
\section{Repaired-Gauge Construction}\label{sec:repaired-gauge-construction}
\subsection{Assumptions and notation}\label{paper2:sec:assumptions}

Let \(u:\RR^{3}\times (T-\delta,T)\to \RR^{3}\), \(p:\RR^{3}\times (T-\delta,T)\to\RR\)
be a suitable weak Navier--Stokes solution with viscosity \(\nu>0\).
For \(z=(x_0,t_0)\) and \(R>0\),
\[
Q_R(z):=B_R(x_0)\times (t_0-R^2,t_0),\qquad
Q_R^*:=B_R\times I_*.
\]
Fix a reference physical center-time \((x_\bullet,t_\star)\) around which \(B_{R_*}(x_\bullet)\) is taken.
For a fixed spatial center \(x_0\) we also use
\[
Q_R(t):=Q_R(x_0,t).
\]
Fix a compact window \(I_*:=[\tau_0-\ell,\tau_0+\ell]\), \(\ell>0\), and write
\(I_{\rm ph}:=[t_\star-\Lambda_{\max}^2\ell,\; t_\star+\Lambda_{\max}^2\ell]\).

For a cylinder \(Q_R(x_0,t_0)=B_R(x_0)\times (t_0-R^2,t_0)\), set
\[
\begin{aligned}
A(u;Q_R(x_0,t_0))&:=R^{-1}\sup_{s\in(t_0-R^2,t_0)}
\int_{B_R(x_0)}|u(x,s)|^2\,dx,\\
E(u;Q_R(x_0,t_0))&:=R^{-1}\int_{Q_R(x_0,t_0)}|\nabla u|^2\,dx\,ds,\\
C(u;Q_R(x_0,t_0))&:=R^{-2}\int_{Q_R(x_0,t_0)}|u|^3\,dx\,ds,\\
D(p;Q_R(x_0,t_0))&:=R^{-2}\int_{t_0-R^2}^{t_0}\int_{B_R(x_0)}
\left|p(x,s)-\langle p(s)\rangle_{B_R(x_0)}\right|^{3/2}\,dx\,ds.
\end{aligned}
\]
Here \(\langle p(s)\rangle_{B_R(x_0)}:=|B_R|^{-1}\int_{B_R(x_0)}p(x,s)\,dx\);
pressure is always understood modulo time-dependent constants.

\begin{definition}[Single-core hypotheses]\label{paper2:def:SC}
Let \((R_*,r_*,M,\eta,\kappa,\Lambda_{\min},\Lambda_{\max})\) satisfy
\(\Lambda_{\min}\le \Lambda_{\max}<\infty\).
We say \((u,p)\) satisfies \(\mathrm{SC}(R_*,r_*,M,\eta,\kappa)\) on \(Q_{R_*}(x_\bullet,t_\star)\)
if the following hold:
\begin{enumerate}
\item[(\textup{SC1})] \(A(u;Q_{R_*}(x_\bullet,t)) + E(u;Q_{R_*}(x_\bullet,t))
  + C(u;Q_{R_*}(x_\bullet,t)) + D(p;Q_{R_*}(x_\bullet,t))\le M\) for every
  \(t\in[t_\star-R_*^2,t_\star)\).
\item[(\textup{SC2})] There exists a physically selected core map
  \((x_{\rm c}(t),\rho_{\rm c}(t))\) on \(I_{\rm ph}\) such that
  \(\rho_{\rm c}(t)\in(0,R_*]\), \(x_{\rm c}(t)\in B_{R_*/2}(x_\bullet)\), and for some
  \(\Theta_0>0\),
  \[
  \int_{B_{r_*}(0)}|\rho_{\rm c}(t)^{-1}u(x_{\rm c}(t)+\rho_{\rm c}(t)y,t)|^3\dd y\ge \eta
  \quad \text{for a.e. }t\in I_{\rm ph}.
  \]
\item[(\textup{SC3})] There exists \(0<r_{\rm sep}\le r_*/2\) such that for every
  \(0<r\le r_{\rm sep}\), every competitor pair
  \[
  \mathcal C_{\rm core}(t):=\left\{(\bar x,\bar\rho):|\bar x-x_\bullet|\le \frac{R_*}{2},\
  0<\bar\rho\le R_*\right\},
  \]
  \((\bar x,\bar\rho)\in \mathcal C_{\rm core}(t)\setminus\{(x_{\rm c}(t),\rho_{\rm c}(t))\}\)
  with
  \[
  \rho_{\rm c}(t)^{-1}|\bar x-x_{\rm c}(t)|
  +\left|\frac{\bar\rho}{\rho_{\rm c}(t)}-1\right|\ge r_{\rm sep}
  \]
  has separated local weighted mass in the same core-scale variables:
  \[
  \left|\int_{B_r}|\rho_{\rm c}(t)^{-1}u(x_{\rm c}(t)+\rho_{\rm c}(t)y,t)|^3\,\dd y
  - \int_{B_r}|\bar\rho^{-1}u(\bar x+\bar\rho y,t)|^3\,\dd y\right|
  \ge \eta.
  \]
  This quantitative non-competing core condition is imposed on the
  selected physical core directly; it is not a consequence of the repaired-gauge
  construction.
\item[(\textup{SC4})] For the selected core,
  \[
  \sup_{t\in I_{\rm ph}}\int_{B_{2R_*}}|y|^{-p}|\rho_{\rm c}(t)^{-1}u(x_{\rm c}(t)+\rho_{\rm c}(t) y,t)|^3\,dy\le M,
  \]
  for some \(0<p<3\).
\item[(\textup{SC5})] The preliminary frame associated with the selected core is
  the absolutely continuous representative
  \[
  0<\Lambda_{\min}\le \Lambda_{\rm pre}(\tau)\le \Lambda_{\max},\quad
  \Lambda_{\rm pre}\in AC(I_*),\quad
  X_{\rm pre}\in AC(I_*;\RR^3).
  \]
  It is tied to the selected physical core by the time change
  \(dt_{\rm c}/d\tau=\Lambda_{\rm pre}(\tau)^2\), \(t_{\rm c}(\tau_0)=t_\star\):
  one has
  \[
  X_{\rm pre}(\tau)=x_{\rm c}(t_{\rm c}(\tau)),\qquad
  \Lambda_{\rm pre}(\tau)=\rho_{\rm c}(t_{\rm c}(\tau))
  \]
  for a.e. \(\tau\in I_*\).
\end{enumerate}
\end{definition}
\noindent
The parameter \(\kappa\) denotes the quantitative constants obtained in the local gauge estimates and is not used in items \textup{(SC1)--(SC5)}.

\begin{remark}
The selected core \((x_{\rm c},\rho_{\rm c})\) and its separation/nondegeneracy
conditions are specified in physical variables; no repaired-gauge constraint is used in
their definition.  This avoids the circularity of defining the core via the chart
to be constructed.
\end{remark}

\subsection{Preliminary frame map from the selected core}\label{paper2:sec:preliminary-frame}

The transport map is constructed from the selected core; the core selection
itself is part of the hypotheses.

Define
\[
\frac{dt_{\rm c}}{d\tau}=\Lambda_{\rm pre}(\tau)^2,\qquad t_{\rm c}(\tau_0)=t_\star,
\]
and use the compatibility in \textup{(SC5)} to identify
\(\Lambda_{\rm pre}(\tau)=\rho_{\rm c}(t_{\rm c}(\tau))\) and
\(X_{\rm pre}(\tau)=x_{\rm c}(t_{\rm c}(\tau))\) a.e.

\begin{proposition}[Preliminary frame map regularity]\label{paper2:prop:prelim-frame}
Under \(\mathrm{SC}(R_*,r_*,M,\eta,\kappa)\), the map
\[
\Phi_{\rm pre}(\tau,Y):=(x,t)=(X_{\rm pre}(\tau)+\Lambda_{\rm pre}(\tau)Y,\ t_{\rm c}(\tau))
\]
is well-defined and absolutely continuous in \(\tau\).
Moreover, \(1/\Lambda_{\rm pre}\), \(\Lambda_{\rm pre}\in L^\infty(I_*)\),
\(X_{\rm pre}\in W^{1,1}_{\rm loc}(I_*;\RR^3)\), and for every \(R\le R_*\),
\[
\Phi_{\rm pre}:Q_R^*\to \Phi_{\rm pre}(Q_R^*)\subset Q_{2R_*}(x_\bullet,t_\star)
\]
is bi-Lipschitz with \(\det \nabla_Y\Phi_{\rm pre}=\Lambda_{\rm pre}^3\), hence
Jacobian and inverse-Jacobian bounds depend only on \(\Lambda_{\min},\Lambda_{\max}\).
\end{proposition}
\begin{proof}
From \textup{(SC5)}, \(\Lambda_{\rm pre}\in AC(I_*)\),
\(0<\Lambda_{\min}\le \Lambda_{\rm pre}\le\Lambda_{\max}<\infty\), and
\(X_{\rm pre}\in AC(I_*;\mathbb R^3)\). The scalar ODE
\[
\frac{dt_{\rm c}}{d\tau}=\Lambda_{\rm pre}(\tau)^2,\qquad t_{\rm c}(\tau_0)=t_\star
\]
therefore has an absolutely continuous strictly increasing solution on \(I_*\), and
\[
\frac{dt_{\rm c}}{d\tau}=\Lambda_{\rm pre}(\tau)^2
\quad\text{a.e. on }I_*.
\]
The compatibility condition \textup{(SC5)} identifies this absolutely continuous frame
with the selected physical core a.e.; no differentiability of \(x_{\rm c}\) or
\(\rho_{\rm c}\) beyond the frame regularity in \textup{(SC5)} is used.
For fixed \(\tau\), \(\nabla_Y\Phi_{\rm pre}(\tau,\cdot)=\Lambda_{\rm pre}(\tau)I\), so
\[
\det\nabla_Y\Phi_{\rm pre}(\tau,\cdot)=\Lambda_{\rm pre}(\tau)^3,
\qquad
\Lambda_{\min}^3\le \det\nabla_Y\Phi_{\rm pre}\le \Lambda_{\max}^3.
\]
Hence, if \(Y_1,Y_2\in B_R\),
\[
\Lambda_{\min}|Y_1-Y_2|
\le |\Phi_{\rm pre}(\tau,Y_1)-\Phi_{\rm pre}(\tau,Y_2)|
\le \Lambda_{\max}|Y_1-Y_2|,
\]
and therefore \(\Phi_{\rm pre}(\tau,\cdot)\) is bi-Lipschitz on \(B_R\). Its inverse has
Lipschitz constant \(\Lambda_{\min}^{-1}\), so Jacobian and inverse-Jacobian bounds
are controlled by \(\Lambda_{\min},\Lambda_{\max}\).

For image inclusions, if \(x=X_{\rm pre}(\tau)+\Lambda_{\rm pre}(\tau)Y\) with
\(|Y|\le R\), then
\[
|x-X_{\rm pre}(\tau)|\le \Lambda_{\max}R,\qquad
X_{\rm pre}(\tau)\in B_{R_*/2}(x_\bullet),
\]
so \(x\in B_{R_*/2+\Lambda_{\max}R}(x_\bullet)\). As \(\tau\in I_*\), we also have
\(t=t_{\rm c}(\tau)\in I_{\rm ph}\). Therefore
\[
\Phi_{\rm pre}(Q_R^*)\subset B_{R_*/2+\Lambda_{\max}R}(x_\bullet)\times I_{\rm ph}.
\]
\end{proof}
\begin{definition}[Preliminary profile]\label{paper2:def:prelim-profile}
For \((x,t)=\Phi_{\rm pre}(\tau,Y)\), define
\[
V_{\rm pre}(Y,\tau):=\Lambda_{\rm pre}(\tau)\,u(x,t),\qquad
P_{\rm pre}(Y,\tau):=\Lambda_{\rm pre}(\tau)^2\,p(x,t).
\]
Hence \(u(x,t)=\Lambda_{\rm pre}(\tau)^{-1}V_{\rm pre}(Y,\tau)\) and
\(p(x,t)=\Lambda_{\rm pre}(\tau)^{-2}P_{\rm pre}(Y,\tau)\) under \(\Phi_{\rm pre}\).
\end{definition}

\subsection{Exact renormalized change of variables}\label{paper2:sec:change}

Define renormalized fields on \(Q_{2R}^*\) by
\[
Y=\frac{x-X(\tau)}{\Lambda(\tau)},\qquad
u(x,t)=\Lambda(\tau)^{-1}V(Y,\tau),\qquad
p(x,t)=\Lambda(\tau)^{-2}P(Y,\tau),
\]
where \(dt/d\tau=\Lambda(\tau)^2\), \(X:\,I_*\to\RR^3\), \(\Lambda:\,I_*\to(0,\infty)\),
and \(\Lambda,X\in AC(I_*)\).
Write
\[
a(\tau):=-\Lambda'(\tau),\qquad
b(\tau):=-X'(\tau).
\]

\begin{lemma}[Exact differential identities]\label{paper2:lem:change-id}
For smooth data and fixed \((x,t)\), the chain rule gives
\[
\partial_t u
=\Lambda^{-3}\left(\partial_\tau V-(X'+\Lambda' Y)\cdot\nabla_Y V-\Lambda'V\right),
\]
\[
\nn_x u=\Lambda^{-2}\nn_Y V,\qquad \Delta_x u=\Lambda^{-3}\Delta_Y V,\qquad
\nn_x p=\Lambda^{-3}\nn_Y P.
\]
In particular
\[
\nn_x\cdot u=\Lambda^{-2}\nn_Y\cdot V.
\]
\end{lemma}
\begin{proof}
For fixed \((x,t)\), \(\tau=\tau(t)\) and \(Y=(x-X(\tau))/\Lambda(\tau)\).
Differentiate
\[
u(x,t)=\Lambda(\tau)^{-1}V\bigl(Y,\tau\bigr)
\]
in \(t\):
\[
\partial_t u
=-\frac{\Lambda'}{\Lambda^2}V
+\frac1\Lambda\partial_\tau V\,\frac{d\tau}{dt}
+\frac1\Lambda\nabla_YV\cdot\frac{dY}{dt}.
\]
Since \(d\tau/dt=\Lambda^{-2}\) and
\[
Y_\tau=-\Lambda^{-1}X'-\frac{\Lambda'}{\Lambda}Y,\qquad
\frac{dY}{dt}=Y_\tau\,\frac{d\tau}{dt},
\]
we have
\[
\partial_tu
=\Lambda^{-3}\left(\partial_\tau V-(X'+\Lambda'Y)\cdot\nabla_YV-\Lambda'V\right).
\]
For spatial derivatives, \(Y=(x-X)/\Lambda\) gives \(\nabla_xY=\Lambda^{-1}I\), hence
\[
\nabla_xu=\Lambda^{-2}\nabla_YV,\qquad
\Delta_xu=\Lambda^{-3}\Delta_YV.
\]
Similarly \(p(x,t)=\Lambda^{-2}P(Y,\tau)\) yields
\[
\nabla_x p=\Lambda^{-3}\nabla_YP.
\]
Taking divergence of \(u=\Lambda^{-1}V\) gives
\[
\nabla_x\cdot u=\Lambda^{-2}\nabla_Y\cdot V.
\]
\end{proof}

\begin{theorem}[Renormalized Navier--Stokes equation]\label{paper2:thm:ren-NS}
Assume \(u,p\) solve
\[
\partial_t u+(u\cdot\nn_x)u+\nn_x p-\nu\Delta_xu=0,\qquad \nn_x\cdot u=0
\]
in distributions. Under the transformation above, \(V,P\) satisfy, on the image
renormalized cylinder, the exact equation
\[
\partial_\tau V+(V\cdot\nn)V+\nn P-\nu\Delta V
+a(\tau)\bigl(V+Y\cdot\nn V\bigr)+b(\tau)\cdot\nn V=0,\qquad \nn\cdot V=0,
\]
with the sign convention fixed by \(a=-\Lambda'\), \(b=-X'\).
Conversely, any \(V,P\) satisfying this equation pull back to a distributional
solution \(u,p\) of Navier--Stokes on the physical image.
\end{theorem}
\begin{proof}
To avoid hidden regularity loss, first write the physical equation tested against
\(\Phi\in C_c^\infty\) in the physical cylinder and then use the smooth pull-back formula
(an approximation argument gives distributions). By Lemma~\ref{paper2:lem:change-id},
\[
\partial_tu+\bigl(u\cdot\nabla_x\bigr)u+\nabla_xp-\nu\Delta_xu=0
\]
becomes
\[
\Lambda^{-3}\!\left(\partial_\tau V-(X'+\Lambda'Y)\cdot\nabla V-\Lambda'V\right)
+\Lambda^{-3}(V\cdot\nabla)V
+\Lambda^{-3}\nabla P
-\Lambda^{-3}\nu\Delta V=0.
\]
Multiplying by \(\Lambda^3\) and using \(a=-\Lambda'\), \(b=-X'\) gives
\[
\partial_\tau V+(V\cdot\nabla)V+\nabla P-\nu\Delta V
+a\bigl(V+Y\cdot\nabla V\bigr)+b\cdot\nabla V=0.
\]
The divergence condition is
\[
\nabla_x\cdot u=\Lambda^{-2}\nabla\cdot V=0
\quad\Longrightarrow\quad
\nabla\cdot V=0.
\]
Conversely, assume \((V,P)\) solve this equation on the renormalized cylinder.
Set \(u,p\) by the same change-of-variables map.
Then the preceding identities run in reverse and give
 \[
\partial_t u+(u\cdot\nabla_x)u+\nabla_xp-\nu\Delta_xu=0,\qquad
\nabla_x\cdot u=0
\]
in \(\mathcal D'\) on the physical image.
\end{proof}

\subsection{Domain localization and support}\label{paper2:sec:domains}

Let \(R>0\) be fixed and \(0<\ell\le \ell_0\). Set \(Q_R^*=B_R\times I_*\).
With \(\Phi\) as above define physical support maps
\[
\mathcal Q^{\rm ph}_{R}:=\Phi(Q_R^*),\qquad
\mathcal Q^{\rm ph}_{2R}:=\Phi(Q_{2R}^*).
\]

\begin{proposition}[Compact support transport]\label{paper2:prop:domain}
Assume \(\Lambda_{\min}\le \Lambda\le \Lambda_{\max}\) on \(I_*\).
Then
\[
\begin{gathered}
X(\tau)\in B_{R_*/2}(x_\bullet),\qquad |x-X(\tau)|\le \Lambda_{\max}R,\\
\mathcal Q^{\rm ph}_{R}\subset B_{R_*/2+ \Lambda_{\max}R}(x_\bullet)\times I_{\rm ph},\\
\mathcal Q^{\rm ph}_{2R}\subset B_{R_*/2+2\Lambda_{\max}R}(x_\bullet)\times I_{\rm ph},
\end{gathered}
\]
for every \(\tau\in I_*\) and \(x=X(\tau)+\Lambda(\tau)Y\) with \(Y\in B_{2R}\).
In particular, if \(|c(\tau)|\le \delta\) from Proposition~\ref{paper2:prop:preliminary-correction},
then \(X(\tau)\in B_{R_*/2+\Lambda_{\max}\delta}(x_\bullet)\) on \(I_*\).
Hence, \(\chi_R\in C_c^\infty(B_{2R})\) pulls back to
\(\chi_R^\tau:=\chi_R\circ \Phi^{-1}\in C_c^\infty(\mathcal Q^{\rm ph}_{2R})\)
with support depending only on \(R,\Lambda_{\min},\Lambda_{\max}\).
\end{proposition}
\begin{proof}
Fix \(\tau\in I_*\) and \(Y\in B_R\). Since \(x=X(\tau)+\Lambda(\tau)Y\),
\[
|x-X(\tau)|\le |Y|\,\Lambda_{\max}\le \Lambda_{\max}R,\qquad
\Lambda_{\max}^{-1}|x-X(\tau)|\le |Y|.
\]
Because \(X(\tau)\in B_{R_*/2}(x_\bullet)\), this gives
\[
x\in B_{R_*/2+\Lambda_{\max}R}(x_\bullet).
\]
Similarly, if \(Y\in B_{2R}\), then \(x\in B_{R_*/2+2\Lambda_{\max}R}(x_\bullet)\).
Together with \(t=t_{\rm c}(\tau)\in I_{\rm ph}\), these give the claimed
inclusions for \(\mathcal Q_{R}^{\rm ph}\) and \(\mathcal Q_{2R}^{\rm ph}\).

For pull-backs of cutoffs, \(\Phi\) is bi-Lipschitz at every \(\tau\), so composition
\[
\chi_R^\tau(x,t):=\chi_R\!\left(\frac{x-X(\tau)}{\Lambda(\tau)}\right)
\]
is smooth on \(\mathcal Q^{\rm ph}_{2R}\), supported where
\(Y\in B_{2R}\), hence supported where \((x,t)\in\mathcal Q^{\rm ph}_{2R}\). The support size depends
only on \(R,\Lambda_{\min},\Lambda_{\max}\) and the fixed \(C_c^\infty\) norm of
\(\chi_R\), not on the particular branch.
\end{proof}

\subsection{Transport of local quantities under the chart}\label{paper2:sec:local-transfer}

Let \(t=t(\tau)=t_{\rm c}(\tau)\) and, for \(R>0\), define
\[
J_\tau^{-}(R):=\bigl(t(\tau)-\Lambda_{\max}^{2}R^2,\ t(\tau)\bigr].
\]
Define renormalized local quantities on \(Q_R^*(\tau)\) by
\[
\begin{aligned}
A_V(\tau,R)&:=\sup_{s\in(\tau-R^2,\tau)}\int_{B_R}|V(\cdot,s)|^2\,dY,\qquad
E_V(\tau,R):=\int_{\tau-R^2}^\tau\!\!\int_{B_R}|\nabla V|^2\,dY\,d s,\\
C_V(\tau,R)&:=\int_{\tau-R^2}^\tau\!\!\int_{B_R}|V|^3\,dY\,d s,\qquad
D_V(\tau,R):=\int_{\tau-R^2}^\tau\!\!\int_{B_R}\Bigl|P-\langle P\rangle_{B_R}\Bigr|^{3/2}\,dY\,d s.
\end{aligned}
\]
For interval notation and a center \(X\), define
\[
A_u(J,R;X):=\sup_{s\in J}\int_{B_R(X)}|u(\cdot,s)|^2\,dx,\quad
E_u(J,R;X):=\int_J\int_{B_R(X)}|\nabla u|^2\,dx\,ds,
\]
\[
C_u(J,R;X):=\int_J\int_{B_R(X)}|u|^3\,dx\,ds,\quad
D_u(J,R;X):=\int_J\int_{B_R(X)}|p-\langle p\rangle_{B_R(X)}|^{3/2}\,dx\,ds.
\]
In the inequalities below we abbreviate
\[
\begin{aligned}
A_u(J,R)&:=A_u(J,R;X(\tau)),& E_u(J,R)&:=E_u(J,R;X(\tau)),\\
C_u(J,R)&:=C_u(J,R;X(\tau)),& D_u(J,R)&:=D_u(J,R;X(\tau)).
\end{aligned}
\]

\begin{proposition}[Cylinder comparison on compact windows]\label{paper2:prop:local-transfer}
Assume \(\Lambda_{\min}\le \Lambda\le \Lambda_{\max}\) and formulas of
Lemma~\ref{paper2:lem:change-id}.
Then for every \(\tau\in I_*\), \(R>0\),
\[
A_V(\tau,R)\le \Lambda_{\min}^{-1}A_u(J_\tau^{-}(R),\Lambda_{\max}R),
\]
\[
E_V(\tau,R)\le \Lambda_{\min}^{-3}E_u(J_\tau^{-}(R),\Lambda_{\max}R),
\]
\[
C_V(\tau,R)\le \Lambda_{\min}^{-2}C_u(J_\tau^{-}(R),\Lambda_{\max}R),
\]
\[
D_V(\tau,R)\le \Lambda_{\min}^{-2}D_u(J_\tau^{-}(R),\Lambda_{\max}R).
\]
In particular, all four renormalized quantities are bounded whenever the physical
quantities are bounded on the corresponding compact physical cylinders.
\begin{proof}
These are direct changes of variables under \(\Phi_{\rm pre}\):
\[
V(\cdot,\tau)=\Lambda(\tau)u(\cdot,t(\tau)),\qquad
P(\cdot,\tau)=\Lambda(\tau)^2 p(\cdot,t(\tau)),
\]
\[
\nabla_YV=\Lambda\nabla_xu,\qquad dY\,d\tau=\Lambda^{-5}\,dx\,dt,\qquad
d\tau=\Lambda^{2}\,dt.
\]
Insert these formulas into each definition and use
\(\Lambda_{\min}\le\Lambda\le\Lambda_{\max}\), plus
\[
\tau-s\in(0,R^2)\implies t(\tau)-\Lambda_{\max}^{2}R^2\le t(s)\le t(\tau)-\Lambda_{\min}^{2}R^2.
\]
\end{proof}
\end{proposition}

\subsection{Local repaired gauge map}\label{paper2:sec:gauge-map}

Fix \(R>0\), \(0<p<3\), and \(\chi_R\in C_c^\infty(B_{2R})\), \(\chi_R\equiv1\) on
\(B_R\). Let the local energy profile space be
\[
\mathcal V_R:=\left\{V\in L^3(B_{2R})\cap W^{1,3/2}(B_{2R};\RR^3):
|Y|^{-p/3}V\in L^3(B_{2R})\right\},
\]
with norm
\[
\|V\|_{\mathcal V_R}:=\|V\|_{L^3(B_{2R})}
+\||Y|^{-p/3}V\|_{L^3(B_{2R})}
+\|\nn V\|_{L^{3/2}(B_{2R})}.
\]
For differentiating the repaired scale and centering gauges we use the
admissible gauge class
\[
\mathcal A_{p,R}:=\left\{V:
\begin{array}{l}
V\in L^3(B_{4R})\cap W^{1,3}_{\rm loc}(B_{4R};\RR^3),\\[2mm]
\displaystyle \int_{B_{4R}}|Y|^{-p}|V|^3\,dY<\infty,\quad
\displaystyle \int_{B_{4R}}|Y|^{-p}|V|\,|Y\cdot\nabla V|\,dY<\infty,\\[2mm]
\displaystyle \int_{B_{4R}}|Y|^{-p-1}|V|^3\,dY<\infty
\end{array}\right\}.
\]
The larger ball is only a domain buffer: all gauge integrals are supported in
\(B_{2R}\), while \((\lambda,c)\) is restricted near \((1,0)\) so that
\(\lambda^{-1}(B_{2R}-c)\subset B_{4R}\).

\begin{definition}[Gauge actions]\label{paper2:def:gauge-action}
For \((\lambda,c)\in (0,\infty)\times\RR^3\), define
\[
V_{\lambda,c}(Y):=\lambda^{-1}V(\lambda^{-1}(Y-c)).
\]
For pressure profiles use
\[
P_{\lambda,c}(Y):=\lambda^{-2}P(\lambda^{-1}(Y-c)).
\]
We also write \(V^{\lambda,c}:=V_{\lambda,c}\) and \(P^{\lambda,c}:=P_{\lambda,c}\).
\end{definition}

\begin{definition}[Finite-dimensional gauge map]\label{paper2:def:gauge-map}
For \(V\in\mathcal V_R\) and \( (\lambda,c) \in (0,\infty)\times\RR^3\), set
\[
\mathcal F(\lambda,c;V)
:=\bigl(\mathcal F_0,\mathcal F_1,\mathcal F_2,\mathcal F_3\bigr),
\]
with
\[
\mathcal F_0(\lambda,c;V):=
\int_{\RR^3}\chi_R(Y)\,|Y|^{-p}|V_{\lambda,c}(Y)|^3\dd Y-\Theta_0,\qquad
\mathcal F_j(\lambda,c;V):=\int_{\RR^3}Y_j\,\chi_R(Y)\,|V_{\lambda,c}(Y)|^2\dd Y
\]
for \(j=1,2,3\). The gauge conditions are \(\mathcal F(\lambda,c;V)=0\).
Codomain is \(\RR^4\), domain \((0,\infty)\times\RR^3\times\mathcal V_R\).
\end{definition}

\begin{definition}[Localized modulation matrix]\label{paper2:def:gauge-nondeg}
For \(J\subset I_*\) and a profile path \(W:J\to\mathcal V_R\), define
\[
J_W(\tau):=
D_{(\lambda,c)}\mathcal F\bigl(1,0;W(\tau)\bigr)
\in\mathbb R^{4\times4}.
\]
The rows are ordered as the scale constraint followed by the three centering
constraints.
\end{definition}

\begin{lemma}[Weighted integrability and differentiation under integrals]\label{paper2:lem:weighted-finite}
Assume \(V,W\in\mathcal A_{p,R}\). Then \(G_0(V):=\int\chi_R|Y|^{-p}|V|^3\)
is finite and
\[
\frac{\dd}{\dd\epsilon}\Big|_{\epsilon=0}
\int \chi_R|Y|^{-p}|V+\epsilon W|^3\,\dd Y
=\int 3\chi_R|Y|^{-p}|V|\,V\cdot W\,\dd Y.
\]
\end{lemma}
\begin{proof}
For \(W\in\mathcal A_{p,R}\), define
\[
F_\varepsilon(Y):=\chi_R(Y)|Y|^{-p}|V(Y)+\varepsilon W(Y)|^3.
\]
For each \(Y\), \(\varepsilon\mapsto F_\varepsilon(Y)\) is \(C^1\) and
\[
\partial_\varepsilon F_\varepsilon
=3\chi_R|Y|^{-p}|V+\varepsilon W|(V+\varepsilon W)\cdot W.
\]
For \(|\varepsilon|\le1\),
\[
|F_\varepsilon(Y)|+|\partial_\varepsilon F_\varepsilon|
\le C\chi_R|Y|^{-p}\bigl(|V|^3+|W|^3\bigr)\in L^1(B_{2R})
\]
by the defining weighted \(L^3\) condition of \(\mathcal A_{p,R}\).
Hence dominated convergence allows differentiation under the integral, and
\[
\frac{\dd}{\dd\varepsilon}\Big|_{\varepsilon=0}\int_{B_{2R}}\chi_R|Y|^{-p}|V+\epsilon W|^3\,\dd Y
=\int_{B_{2R}}3\chi_R|Y|^{-p}|V|\,V\cdot W\,\dd Y.
\]
\end{proof}

\begin{lemma}[Gauge covariance]\label{paper2:lem:gauge-covariance}
For \(V\in\mathcal V_R\), \((\lambda,c)\in(0,\infty)\times\RR^3\), and
\(\alpha=0,\dots,3\),
\[
\mathcal F_\alpha(1,0;V_{\lambda,c})=\mathcal F_\alpha(\lambda,c;V).
\]
\begin{proof}
For each component this is a direct substitution:
\[
\mathcal F_0(1,0;V_{\lambda,c})
=\int\chi_R(Y)|Y|^{-p}|V_{\lambda,c}(Y)|^3-\Theta_0\,\dd Y
=\mathcal F_0(\lambda,c;V),
\]
\[
\mathcal F_j(1,0;V_{\lambda,c})
=\int Y_j\chi_R(Y)|V_{\lambda,c}(Y)|^2\,\dd Y
=\mathcal F_j(\lambda,c;V),\quad j=1,2,3.
\]
\end{proof}
\end{lemma}

\begin{proposition}[Differentiability of \(\mathcal F\)]\label{paper2:prop:gauge-C1}
For each fixed \(V\in\mathcal A_{p,R}\), the map
\((\lambda,c)\mapsto \mathcal F(\lambda,c;V)\) is \(C^1\) in a neighborhood of
\((1,0)\). In each component \( \mathcal F_\alpha\) all derivatives under the
integral sign are legitimate.
\begin{proof}
Fix \((\lambda,c)\) near \((1,0)\), and abbreviate
\[
V_{\lambda,c}(Y):=\lambda^{-1}V\!\left(\lambda^{-1}(Y-c)\right).
\]
By direct differentiation,
\[
\partial_\lambda V_{\lambda,c}
=-\lambda^{-1}V_{\lambda,c}-\lambda^{-1}(Y-c)\cdot\nabla V_{\lambda,c},\qquad
\partial_{c_k}V_{\lambda,c}=-\lambda^{-2}\partial_kV(\lambda^{-1}(Y-c))=-\partial_kV_{\lambda,c}.
\]
The neighborhood of \((1,0)\) is chosen so that the arguments of \(V\) remain
inside \(B_{4R}\) for \(Y\in\supp\chi_R\). The scale derivative is controlled by
the two weighted integrability conditions
\[
\int |Y|^{-p}|V|^3<\infty,\qquad
\int |Y|^{-p}|V|\,|Y\cdot\nabla V|<\infty,
\]
transported by the smooth change of variables. The translation derivative in
the scale row is treated by the source identity
\[
3\int |Y|^{-p}|V|\,V\cdot\partial_kV
=p\int Y_k|Y|^{-p-2}|V|^3,
\]
obtained first for smooth compactly supported profiles and then by approximation
in \(\mathcal A_{p,R}\); the right-hand side is finite because
\(\int |Y|^{-p-1}|V|^3<\infty\). The centering rows are supported in \(B_{2R}\)
and are controlled by \(V\in L^2_{\rm loc}\) and \(\nabla V\in L^3_{\rm loc}\).
These bounds give an integrable majorant for the difference quotients, so
dominated convergence applies.
\[
\mathcal F_0(\lambda,c;V)
=\int_{\RR^3}\chi_R(Y)\,|Y|^{-p}|V_{\lambda,c}|^3\,\dd Y-\Theta_0,
\]
for which dominated convergence gives
\[
\partial_{\lambda}\mathcal F_0
=\int_{\RR^3}\chi_R|Y|^{-p}\,3|V_{\lambda,c}|\,V_{\lambda,c}\cdot\partial_\lambda V_{\lambda,c}\,dY,
\]
\[
\partial_{c_k}\mathcal F_0
=\int_{\RR^3}\chi_R|Y|^{-p}\,3|V_{\lambda,c}|\,V_{\lambda,c}\cdot\partial_{c_k}V_{\lambda,c}\,dY.
\]
Similarly
\[
\mathcal F_j(\lambda,c;V)=\int_{\RR^3}Y_j\chi_R(Y)\,|V_{\lambda,c}|^2\,\dd Y,\quad j=1,2,3,
\]
hence
\[
\partial_\lambda \mathcal F_j
\ =\int_{\RR^3}2Y_j\chi_R\,V_{\lambda,c}\cdot\partial_\lambda V_{\lambda,c}\,dY,
\]
\[
\partial_{c_k}\mathcal F_j
=\int_{\RR^3}2Y_j\chi_R\,V_{\lambda,c}\cdot\partial_{c_k}V_{\lambda,c}\,dY.
\]
Thus each component admits pointwise \(C^1\) formulas and these are continuous in
\((\lambda,c)\) near \((1,0)\). Therefore \((\lambda,c)\mapsto\mathcal F(\lambda,c;V)\in C^1\)
near \((1,0)\).
\end{proof}
\end{proposition}

\begin{proposition}[Full local Jacobian at \((1,0)\)]\label{paper2:prop:full-jacobian}
Let \(V\in\mathcal A_{p,R}\).
Set \(V_0(Y)=V(Y)\).
At \((\lambda,c)=(1,0)\), the matrix
\[
J(V):=D_{(\lambda,c)}\mathcal F(1,0;V)\in \RR^{4\times4}
\]
has rows
\[
J_\alpha=\bigl(\partial_\lambda \mathcal F_\alpha,\ \partial_{c_1}\mathcal F_\alpha,\ 
\partial_{c_2}\mathcal F_\alpha,\ \partial_{c_3}\mathcal F_\alpha\bigr),\quad \alpha=0,\dots,3.
\]
Writing \(W_{\rm scl}=V_0+Y\cdot\nn V_0\), one has
\begin{align}
\partial_\lambda \mathcal F_0(1,0;V)
&=
-3\int\chi_R|Y|^{-p}|V_0|\,V_0\cdot W_{\rm scl}\,\dd Y,\label{paper2:eq:jac00}\\
\partial_{c_k}\mathcal F_0(1,0;V)
&=-3\int_{\RR^3}\chi_R |Y|^{-p}|V_0|\,V_0\cdot\partial_kV_0\,\dd Y,\label{paper2:eq:jac0k}\\
\partial_\lambda \mathcal F_j(1,0;V)
&=-2\int_{\RR^3}Y_j\chi_R\,V_0\cdot W_{\rm scl}\,\dd Y,\label{paper2:eq:jacjlambda}\\
\partial_{c_k}\mathcal F_j(1,0;V)
&=-2\int_{\RR^3}\chi_R Y_j V_0\cdot\partial_kV_0\,\dd Y,\label{paper2:eq:jacjk}
\end{align}
for \(j,k=1,2,3\).
\begin{proof}
Differentiate the formulas in Definition~\ref{paper2:def:gauge-map} at \((\lambda,c)=(1,0)\).

Since \(V_{\lambda,c}(Y)=\lambda^{-1}V(\lambda^{-1}(Y-c))\),
\[
\partial_\lambda V_{\lambda,c}
 =-\lambda^{-1}V_{\lambda,c}-\lambda^{-1}Y\cdot\nabla V_{\lambda,c},\qquad
\partial_{c_k}V_{\lambda,c}=-\lambda^{-2}\partial_kV(\lambda^{-1}(Y-c)).
\]
At \((1,0)\), \(\partial_\lambda V_{\lambda,c}|_{(1,0)}=-(V+Y\cdot\nabla V)\),
\(\partial_{c_k}V_{\lambda,c}|_{(1,0)}=-\partial_kV\).

For \(\mathcal F_0\), write
\[
\mathcal F_0(\lambda,c;V)=\int \chi_R(Y)|Y|^{-p}|V_{\lambda,c}(Y)|^3\,dY-\Theta_0.
\]
Derivative in \(\lambda\):
\[
\partial_\lambda \mathcal F_0
\ =\int\chi_R|Y|^{-p}3|V_{\lambda,c}|\,V_{\lambda,c}\cdot\partial_\lambda V_{\lambda,c}\,dY
\]
Evaluating at \((1,0)\) yields \eqref{paper2:eq:jac00}.

For \(k=1,2,3\), \(c_k\)-derivative of \(\mathcal F_0\) has no cutoff term:
\[
\partial_{c_k}\mathcal F_0(1,0;V)
=\int \chi_R|Y|^{-p}\,3|V|\,V\cdot(-\partial_kV)\,dY,
\]
which is \eqref{paper2:eq:jac0k}.

For \(\mathcal F_j\),
\[
\mathcal F_j(\lambda,c;V)=\int Y_j\chi_R(Y)\,|V_{\lambda,c}|^2\,dY.
\]
Hence
\[
\partial_\lambda \mathcal F_j
=\int 2Y_j\chi_R\,V_{\lambda,c}\cdot\partial_\lambda V_{\lambda,c}\,dY,
\]
and at \((1,0)\):
\[
\partial_\lambda \mathcal F_j(1,0;V)
=-2\int Y_j\chi_R\,V\cdot W_{\rm scl}\,dY,
\]
which is \eqref{paper2:eq:jacjlambda}. Similarly,
\[
\partial_{c_k}\mathcal F_j(1,0;V)
=\int 2Y_j\chi_R\,V\cdot\partial_{c_k}V_{\lambda,c}\big|_{(1,0)}\,dY
=-2\int \chi_RY_jV\cdot\partial_kV\,dY,
\]
giving \eqref{paper2:eq:jacjk}.
\end{proof}
\end{proposition}

\begin{proposition}[Translation block and Schur-complement invertibility]\label{paper2:thm:quant-jac}
Let \(V:J\to\mathcal A_{p,R}\) satisfy the repaired gauge constraints
\[
\mathcal F(1,0;V(\tau))=0\qquad\text{for a.e. }\tau\in J.
\]
Write the modulation matrix in block form
\[
J(V(\tau))=
\begin{pmatrix}
\alpha(\tau)&u(\tau)^T\\
v(\tau)&A(\tau)
\end{pmatrix},
\]
where \(\alpha=\partial_\lambda\mathcal F_0\), \(u_\ell=\partial_{c_\ell}\mathcal F_0\),
\(v_j=\partial_\lambda\mathcal F_j\), and \(A_{j\ell}=\partial_{c_\ell}\mathcal F_j\).
Assume that for a.e. \(\tau\in J\)
\[
m_\chi(\tau):=\int |V(Y,\tau)|^2\chi_R(Y)\,dY\ge m_0>0,
\]
\[
\int_{A_R}|V(Y,\tau)|^2\,dY\le\varepsilon_0,\qquad
3C_\chi\varepsilon_0\le \frac{m_0}{2},
\]
where \(A_R=\{R\le |Y|\le2R\}\) and
\[
C_\chi:=\max_{j,\ell}\|Y_j\partial_\ell\chi_R\|_{L^\infty}.
\]
Assume also
\[
C_u:=\operatorname*{ess\,sup}_{\tau\in J}\int |Y|^{-p-1}|V(Y,\tau)|^3\,dY<\infty
\]
and, with \(C_\nabla:=\||Y|\nabla\chi_R\|_{L^\infty}\),
\[
6p\left(\frac{2}{m_0}\right)C_u\,R C_\nabla\varepsilon_0
\le \frac{p\Theta_0}{2}.
\]
Then \(J(V(\tau))\) is invertible for a.e. \(\tau\in J\), and
\[
\operatorname*{ess\,sup}_{\tau\in J}\|J(V(\tau))^{-1}\|_{\infty\to\infty}
\le K(R,p,\Theta_0,m_0,C_u,\varepsilon_0,\chi_R).
\]
\end{proposition}
\begin{proof}
The entries are finite by \(V(\tau)\in\mathcal A_{p,R}\) and the compact support
of \(\chi_R\). With the present parameter convention,
\(\partial_\lambda V_{\lambda,c}|_{(1,0)}=-(V+Y\cdot\nabla V)\). Hence the
scale entry satisfies the signed transversality identity
\[
\alpha(\tau)=-p\int |Y|^{-p}|V(Y,\tau)|^3\,dY=-p\Theta_0
\]
on the gauge surface. Thus \(|\alpha|\ge p\Theta_0\).

For the centering block, integration by parts gives
\[
A_{j\ell}
=-\delta_{j\ell}\int |V|^2\chi_R
-\int |V|^2Y_j\partial_\ell\chi_R.
\]
Writing \(A=-m_\chi I_3+E\), the annular support of \(\nabla\chi_R\) gives
\[
|E_{j\ell}|\le C_\chi\int_{A_R}|V|^2\le C_\chi\varepsilon_0,\qquad
\|E\|_{\infty\to\infty}\le3C_\chi\varepsilon_0\le m_0/2.
\]
Thus \(A=-m_\chi(I_3-m_\chi^{-1}E)\) and
\(\|m_\chi^{-1}E\|_{\infty\to\infty}\le1/2\), so \(A^{-1}\) exists and
\(\|A^{-1}\|_{\infty\to\infty}\le2/m_0\).

For the scale-to-translation block, the repaired scale-row formula gives
\[
\left|u_\ell\right|
\le p\int |Y|^{-p-1}|V|^3\,dY
\le p\int |Y|^{-p-1}|V|^3\,dY,
\]
so \(\|u\|_\infty\le pC_u\). For the translation-to-scale block,
\[
v_j=-\int |V|^2Y_j(Y\cdot\nabla\chi_R)\,dY
\]
because the centering constraint cancels the interior term. Therefore
\[
\|v\|_\infty\le2R\,C_\nabla\varepsilon_0.
\]
The Schur complement \(S=\alpha-u^TA^{-1}v\) satisfies
\[
|u^TA^{-1}v|
\le 3\|u\|_\infty\|A^{-1}\|_{\infty\to\infty}\|v\|_\infty
\le 6p\left(\frac{2}{m_0}\right)C_uRC_\nabla\varepsilon_0
\le \frac{p\Theta_0}{2}.
\]
Since \(|\alpha|=p\Theta_0\), this gives \(|S|\ge p\Theta_0/2\).
The block inverse formula
\[
J^{-1}=
\begin{pmatrix}
S^{-1} & -S^{-1}u^TA^{-1}\\
-A^{-1}vS^{-1}&A^{-1}+A^{-1}vS^{-1}u^TA^{-1}
\end{pmatrix}
\]
then gives the displayed bound.
\end{proof}

\subsection{Repaired gauge by implicit inversion in time}\label{paper2:sec:gauge-solve}
\begin{lemma}[Time chain rule for preliminary profile]\label{paper2:lem:diff-F-pre}
For each
\(\alpha=0,\dots,3\), set \(G_\alpha(\tau):=\mathcal F_\alpha(1,0;V_{\rm pre}(\tau))\).
Then \(G_\alpha\in W^{1,1}_{\rm loc}(I_*)\) and for a.e. \(\tau\),
\[
G_\alpha'(\tau)=D_V\mathcal F_\alpha(1,0;V_{\rm pre}(\tau))[\partial_\tau V_{\rm pre}(\tau)].
\]
The right-hand side is obtained from the profile formulas of
Proposition~\ref{paper2:prop:gauge-C1} and is in \(L^1_{\rm loc}(I_*)\).
\begin{proof}
Let \(V_{\rm pre}^\varepsilon\) be a spatial and temporal Steklov mollification
inside a slightly larger compact cylinder.  For the smoothed profiles the
ordinary chain rule gives
\[
\frac d{d\tau}\mathcal F_\alpha(1,0;V_{\rm pre}^\varepsilon(\tau))
=D_V\mathcal F_\alpha(1,0;V_{\rm pre}^\varepsilon(\tau))
[\partial_\tau V_{\rm pre}^\varepsilon(\tau)] .
\]
The transformed equation expresses \(\partial_\tau V_{\rm pre}^\varepsilon\)
as the mollification of
\(\nu\Delta V_{\rm pre}-(V_{\rm pre}\cdot\nabla)V_{\rm pre}-\nabla P_{\rm pre}\)
plus the chart transport terms.  Pairing this distribution with the compactly
supported Gateaux gradients below is legitimate by the local suitable bounds,
the pressure bound, and the weighted \(L^3\) control in \textup{(SC4)}.  Passing
\(\varepsilon\downarrow0\) gives the asserted \(W^{1,1}_{\rm loc}\) identity.
For \(W\in\mathcal V_R\), these derivative formulas hold:
\[
\begin{aligned}
D_V\mathcal F_0(1,0;V)[W]&=\int_{\mathbb R^3}3\chi_R|Y|^{-p}|V|\,V\cdot W\,\dd Y,\\
D_V\mathcal F_j(1,0;V)[W]&=\int_{\mathbb R^3}2Y_j\chi_R\,V\cdot W\,\dd Y,\qquad j=1,2,3,
\end{aligned}
\]
where both right-hand sides are finite on \(\mathcal V_R\) by compact support, \(p<3\), and
\(V,W\in L^3\cap W^{1,3/2}_{\mathrm{loc}}\).
The right-hand side is finite by the explicit formulas above and by the distributional
equation paired with compactly supported test gradients; the local energy,
pressure, and weighted \(L^3\) bounds give an \(L^1_{\rm loc}(I_*)\) majorant.
\end{proof}
\end{lemma}
\begin{lemma}[AC implicit inversion with parameters]\label{paper2:lem:AC-IFT}
Let \(U\subset\RR^4\) be open and \(G:I_*\times U\to\RR^4\). Assume:
\begin{enumerate}
\item[(i)] for each \(s\in U\), \(\tau\mapsto G(\tau,s)\in W^{1,1}_{\rm loc}(I_*;\RR^4)\), and
  \(D_\tau G\in L^1_{\rm loc}(I_*;L^\infty(U;\RR^4))\);
\item[(ii)] for a.e. \(\tau\in I_*\), \(s\mapsto G(\tau,s)\) is \(C^1\) on \(U\), with \(D_sG\) continuous in \((\tau,s)\);
\item[(iii)] for some \((\tau_0,s_0)\in I_*\times U\), \(G(\tau_0,s_0)=0\), and there is \(\rho,C>0\) such that
\[
\sup_{(\tau,s)\in I_*\times B_\rho(s_0)}\|D_sG(\tau,s)^{-1}\|_{\mathcal L(\RR^4,\RR^4)}\le C.
\]
\end{enumerate}
Then after possibly shrinking \(I_*\) to a subinterval \(J\ni\tau_0\), there is a unique
map \(s\in AC(J;U)\), \(s(\tau_0)=s_0\), with
\[
G(\tau,s(\tau))=0\quad\text{a.e. in }J.
\]
Moreover
\[
s'(\tau)=-D_sG(\tau,s(\tau))^{-1}\,D_\tau G(\tau,s(\tau))\quad\text{a.e. on }J.
\]
\begin{proof}
By (ii), the matrix field \(D_sG\) is continuous in \(s\), so there is \(r_0\in(0,\rho)\) such that
\(\inf_{(\tau,s)\in I_*\times B_{r_0}(s_0)}|\det D_sG(\tau,s)|>0\) after possibly
shrinking \(I_*\); this keeps \(D_sG^{-1}\) bounded and continuous on \(B_{r_0}(s_0)\).

Set
\[
F(\tau,s):=-D_sG(\tau,s)^{-1}D_\tau G(\tau,s).
\]
Then \(F(\tau,\cdot)\) is continuous on \(B_{r_0}(s_0)\) for a.e. \(\tau\) and
\(F(\cdot,s)\in L^1_{\rm loc}(I_*;\mathbb R^4)\) by (i), hence \(F\) is a Carathéodory
right-hand side.  Since the gauge map is \(C^1\) in the finite variables, after
shrinking \(r_0\) there is \(L\in L^1(J)\) such that
\[
|F(\tau,s_1)-F(\tau,s_2)|\le L(\tau)|s_1-s_2|
\quad\text{for }s_1,s_2\in B_{r_0}(s_0).
\]
Choose \(\delta>0\) so small that
\[
\int_{\tau_0-\delta}^{\tau_0+\delta}\sup_{s\in B_{r_0}(s_0)}|F(\sigma,s)|\,d\sigma<r_0.
\]
On \(J=[\tau_0-\delta,\tau_0+\delta]\cap I_*\) define the Picard operator
\[
(\mathcal Ts)(\tau):=s_0+\int_{\tau_0}^{\tau}F(\sigma,s(\sigma))\,d\sigma
\]
on \(C(J;B_{r_0}(s_0))\).  The preceding bound keeps \(\mathcal T\) in the same
closed ball, and the Lipschitz estimate gives
\[
\|\mathcal T s_1-\mathcal T s_2\|_{C(J)}
\le \left(\int_J L(\sigma)\,d\sigma\right)\|s_1-s_2\|_{C(J)}.
\]
After decreasing \(\delta\) once more, the coefficient is \(<1\), so Banach's
fixed-point theorem gives a unique solution
\[
s(\tau)=s_0+\int_{\tau_0}^{\tau}F(\sigma,s(\sigma))\,d\sigma,\qquad
\tau\in J,
\]
which stays in \(B_{r_0}(s_0)\subset U\). This solution is absolutely continuous.

Define \(H(\tau):=G(\tau,s(\tau))\). Since \(G(\cdot,s(\cdot))\in W^{1,1}_{\rm loc}\) from (i),
the chain rule gives for a.e. \(\tau\in J\)
\[
H'(\tau)=D_\tau G(\tau,s(\tau))+D_sG(\tau,s(\tau))s'(\tau)
=D_\tau G(\tau,s(\tau))+D_sG(\tau,s(\tau))F(\tau,s(\tau))=0.
\]
Hence \(H\equiv G(\tau_0,s_0)=0\) on \(J\), which is the claimed constraint.
Uniqueness follows from Grönwall on the difference of two AC solutions on \(B_{r_0}(s_0)\),
using Lipschitz continuity of \(F(\tau,\cdot)\) on that ball.
\end{proof}
\end{lemma}

\begin{definition}[Retained nondegenerate repaired-gauge regime]\label{paper2:def:admissible-gauge-path}
Let \(J\subset I_*\) be compact. The retained nondegenerate repaired-gauge
regime over the preliminary profile \(V_{\rm pre}\) consists of branches for
which there are
\[
\lambda\in AC(J;(0,\infty)),\qquad c\in AC(J;\mathbb R^3),
\]
such that, for a.e. \(\tau\in J\),
\[
V(\tau):=V_{\rm pre}^{\lambda(\tau),c(\tau)}\in\mathcal A_{p,R},\qquad
\mathcal F(1,0;V(\tau))=0.
\]
The path is quantitatively nondegenerate if the translation-block, annular,
weighted first-moment, and Schur-complement bounds in
Proposition~\ref{paper2:thm:quant-jac} hold on \(J\).
If no such local AC selection exists, or if the quantitative nondegeneracy
hypotheses fail persistently, the branch lies in the gauge-degenerate,
multibubble, or cascade alternatives rather than in the retained regime.
\end{definition}

\begin{proposition}[Gauge regularity on the retained regime]\label{paper2:prop:preliminary-correction}
Let \((\lambda,c)\) be the repaired-gauge selection of the retained
nondegenerate regime on \(J\). Then
\[
\mathcal F\bigl(\lambda(\tau),c(\tau);V_{\rm pre}(\tau)\bigr)=0
\quad\text{for a.e. }\tau\in J,
\]
and \((\lambda,c)\in AC(J;\mathbb R^4)\) by definition. If the path is
quantitatively nondegenerate, then
\[
\operatorname*{ess\,sup}_{\tau\in J}\|J(V(\tau))^{-1}\|<\infty.
\]
\end{proposition}
\begin{proof}
The constraint follows from Lemma~\ref{paper2:lem:gauge-covariance} applied to
\(V(\tau)=V_{\rm pre}^{\lambda(\tau),c(\tau)}\). The absolute continuity is part
of Definition~\ref{paper2:def:admissible-gauge-path}. The uniform inverse bound is
Proposition~\ref{paper2:thm:quant-jac}.
\end{proof}

\begin{definition}[Repaired chart]\label{paper2:def:repaired-chart}
Define
\[
\Lambda(\tau):=\Lambda_{\rm pre}(\tau)\lambda(\tau),\qquad
X(\tau):=X_{\rm pre}(\tau)+\Lambda_{\rm pre}(\tau)c(\tau),
\]
\[
V(\tau):=V_{\rm pre}^{\lambda(\tau),c(\tau)},\qquad
P(\tau):=P_{\rm pre}^{\lambda(\tau),c(\tau)}\ \text{(same physical pressure pull-back under }\Phi).
\]
\end{definition}

\begin{theorem}[AC repaired parameters]\label{paper2:thm:repaired-existence}
Assume \(\mathrm{SC}(R_*,r_*,M,\eta,\kappa)\), and suppose the branch lies in the
retained nondegenerate repaired-gauge regime on \(J\), with selection
\((\lambda,c)\). Then the corrected
chart \((\Lambda,X)\) is absolutely continuous on \(J\), satisfies
\[
a(\tau)=-\Lambda'(\tau)\in L^1_{\rm loc}(J),\qquad
b(\tau)=-X'(\tau)\in L^1_{\rm loc}(J),
\]
and the gauge constraints \(\mathcal F(1,0;V(\tau))=0\) hold a.e.
Equivalently,
\[
\mathcal F(\lambda(\tau),c(\tau);V_{\rm pre}(\tau))=0
\]
for the repaired-gauge selection \((\lambda,c)\).
\end{theorem}
\begin{proof}
By Definition~\ref{paper2:def:admissible-gauge-path}, \((\lambda,c)\in AC(J)\) and the
repaired gauge constraints hold for \(V(\tau)\) a.e.
By definition \(V(\tau)=V_{\rm pre}^{\lambda(\tau),c(\tau)}\), so Lemma~\ref{paper2:lem:gauge-covariance}
applied at \((\lambda,c)=(\lambda(\tau),c(\tau))\) gives
\[
\mathcal F(1,0;V_{\rm pre}^{\lambda,c})=\mathcal F(\lambda,c;V_{\rm pre})
\]
implies \(\mathcal F(1,0;V(\tau))=0\) a.e.

Since \(\Lambda=\Lambda_{\rm pre}\lambda\) and \(X=X_{\rm pre}+\Lambda_{\rm pre}c\),
and \(\Lambda_{\rm pre},X_{\rm pre}\) are AC by \textup{(SC5)}, the products and sums
preserve AC, hence \((\Lambda,X)\in AC(I_*)\). Differentiating,
\[
a=-\Lambda' =-(\Lambda_{\rm pre}'\lambda+\Lambda_{\rm pre}\lambda')\in L^1_{\rm loc},\qquad
b=-X' =-(X_{\rm pre}'+\Lambda_{\rm pre}'c+\Lambda_{\rm pre}c')\in L^1_{\rm loc}.
\]
This gives the claimed regularity.
\end{proof}

\begin{remark}[Gauge degeneracy]
Failure of invertibility of \(J\) is the gauge-degenerate alternative to the
retained single-core regime for this fixed renormalization mechanism.  In the
Type II decomposition of alternatives this case belongs to the
multibubble/cascade and scale-rigidity alternatives, rather than to the
representation theorem.
\end{remark}

\subsection{Pressure transport and decomposition}\label{paper2:sec:pressure}
\begin{lemma}[Renormalized pressure equation]\label{paper2:lem:pressure-eq}
If \(V,P\) satisfy the renormalized Navier--Stokes equation in
Theorem~\ref{paper2:thm:ren-NS}, then
\[
-\Delta P=\partial_i\partial_j(V_iV_j)\qquad\text{in }\mathcal D'(Q_{2R}^*).
\]
\end{lemma}
\begin{proof}
Take test \(\varphi\in C_c^\infty(Q_{2R}^*)\). Since \(V,P\in L^2_{\rm loc}\times
L^{3/2}_{\rm loc}\) on \(Q_{2R}^*\), all pairings below are legitimate.
Expand
\[
\nn\cdot\bigl((V\cdot\nabla)V\bigr)
=\sum_{i,j}\partial_i\partial_j(V_iV_j)-\sum_{i,j}\partial_iV_i\,\partial_jV_j.
\]
Because \(\nn\cdot V=0\), this reduces to \(\partial_i\partial_j(V_iV_j)\).
Taking divergence of
\[
\partial_\tau V+(V\cdot\nabla)V+\nabla P-\nu\Delta V
+a(V+Y\cdot\nabla V)+b\cdot\nabla V=0
\]
and using \(\nn\cdot V=0\), \(\nn\cdot(V+Y\cdot\nabla V)=0\), \(\nn\cdot(b\cdot\nabla V)=b\cdot\nn(\nn\cdot V)=0\),
and \(\nn\cdot\Delta V=\Delta(\nn\cdot V)=0\), we get in \(\mathcal D'\):
\[
\Delta P+\partial_i\partial_j(V_iV_j)=0.
\]
Hence \(-\Delta P=\partial_i\partial_j(V_iV_j)\) on \(Q_{2R}^*\).
\end{proof}

\begin{theorem}[Localized pressure decomposition]\label{paper2:thm:pressure-decomp}
Let \(\zeta\in C_c^\infty(B_{2R})\), \(\zeta\equiv1\) on \(B_R\), and assume
\(V\in L^1_{\mathrm{loc}}(I_*;L^3(B_{2R}))\), \(P\in L^1_{\mathrm{loc}}(I_*;L^{3/2}(B_{2R}))\).
For a.e. \(\tau\in I_*\) there exists
\[
P(\cdot,\tau)=P_{\rm loc}(\cdot,\tau)+H(\cdot,\tau)
\]
on \(B_R\) such that
\[
\begin{cases}
-\Delta P_{\rm loc}(\cdot,\tau)=\partial_i\partial_j\bigl(\zeta V_iV_j\bigr),\\
-\Delta H(\cdot,\tau)=0\ \text{on }B_R,\quad
\fint_{B_R}H(\cdot,\tau)\,dY=0.
\end{cases}
\]
Moreover
\[
\|P_{\rm loc}(\tau)\|_{L^{3/2}(B_R)}
\lesssim \|V(\tau)\|_{L^3(B_{2R})}^2,\qquad
\|H(\tau)\|_{L^q(B_r)}\lesssim_{q,R,r}
\|P(\tau)\|_{L^{3/2}(B_{2R})}+\|V(\tau)\|_{L^3(B_{2R})}^2
\]
for each \(1\le q<\infty\) and each \(0<r<R\), with constants depending only on \(R\), \(r\), and \(q\).
\begin{proof}
For fixed \(\tau\), set
\[
P_{\rm loc}= \mathcal R_i\mathcal R_j(\zeta V_iV_j),\qquad
H:=P-P_{\rm loc}.
\]
\(\zeta V_iV_j\in L^{1}_{\mathrm{loc}}(I_*;L^{3/2}(B_{2R}))\), so the Calderón--Zygmund
operator \(f\mapsto \mathcal R_i\mathcal R_j f\) gives
\[
P_{\rm loc}(\cdot,\tau)=\mathcal R_i\mathcal R_j(\zeta V_iV_j)(\cdot,\tau),
\qquad
\|P_{\rm loc}(\tau)\|_{L^{3/2}(B_R)}
\lesssim \|V(\tau)\|_{L^3(B_{2R})}^2.
\]
Because \(\mathcal R_i\mathcal R_j\) is linear bounded on \(L^{3/2}\),
\[
P_{\rm loc}\in L^{1}_{\mathrm{loc}}(I_*;L^{3/2}(B_R)).
\]
Since \(-\Delta P_{\rm loc}=\partial_i\partial_j(\zeta V_iV_j)\) and \(\zeta\equiv1\) on \(B_R\),
 \[
-\Delta H(\cdot,\tau)
=\partial_i\partial_j\bigl((1-\zeta)V_iV_j\bigr)=0 \quad\text{in }B_R,
\]
so \(H(\cdot,\tau)\) is harmonic on \(B_R\). Imposing \(\fint_{B_R}H=0\) fixes the
time-dependent additive constant uniquely in each time slice.
For \(q\in[1,\infty)\) and \(0<r<R\), harmonicity and local
\(L^p\)-to-\(L^q\) interior estimates imply
\[
\|H(\tau)\|_{L^q(B_r)}\lesssim_{q,R,r}\|H(\tau)\|_{L^{3/2}(B_R)}.
\]
Finally, \(\|H(\tau)\|_{L^{3/2}(B_R)}\le \|P(\tau)\|_{L^{3/2}(B_R)}+\|P_{\rm loc}(\tau)\|_{L^{3/2}(B_R)}\),
giving the announced bound in terms of \(P\) and \(V\).
\end{proof}
\end{theorem}
\begin{remark}[Measurability and normalization]\label{paper2:rem:pressure-time}
If \(V,P\) are measurable in \(\tau\), then \(\tau\mapsto P_{\rm loc}(\cdot,\tau)\)
is measurable in \(L^{3/2}(B_R)\) by boundedness of \(\mathcal R_i\mathcal R_j\), hence
\(\tau\mapsto H(\cdot,\tau)\) is measurable in each \(L^q(B_r)\), \(0<r<R\).
The normalization \(\fint_{B_R}H(\cdot,\tau)\,dY=0\) fixes the time-dependent additive constant
in \(H\) uniquely for each \(\tau\), and therefore the decomposition is a legitimate
renormalized pressure splitting.
\end{remark}

\subsection{Modulation from differentiated constraints}\label{paper2:sec:modulation}
\begin{lemma}[Differentiation of constraints]\label{paper2:lem:diff-f}
For the repaired profile \(V\) constructed above, with
\(|Y|^{-p}|V|^3\in L^1_{\rm loc}(I_*;L^1)\),
Then \(G_\alpha(\tau):=\mathcal F_\alpha(1,0;V(\tau))\) are absolutely
continuous and
\[
\frac{\dd}{\dd\tau}G_\alpha(\tau)
=\int_{\RR^3}\mathbf D_\alpha(V,\nn V,\nabla\chi_R)\,\partial_\tau V\,\dd Y
\]
for a.e. \(\tau\). Here \(\mathbf D_\alpha\) is the \(L^1\)-integrable Gateaux
gradient associated to \(\mathcal F_\alpha\).
\begin{proof}
\[
G_\alpha(\tau)=\mathcal F_\alpha(1,0;V(\tau))
\]
is first evaluated on spatial-temporal mollifications \(V^\varepsilon\).  For
those smooth profiles the ordinary chain rule gives
\[
\frac d{d\tau}G_\alpha^\varepsilon(\tau)
=D_V\mathcal F_\alpha(1,0;V^\varepsilon(\tau))\,[\partial_\tau V^\varepsilon(\tau)].
\]
The local suitable bounds, pressure estimates, and compact support of the
Gateaux gradients below allow passage to the limit \(\varepsilon\downarrow0\),
which gives the displayed derivative formula for \(V\).
For \(\mathcal F_0\),
\[
\mathbf D_0(V,\nabla V,\nabla\chi_R)
=3\chi_R|Y|^{-p}|V|\,V;
\]
for \(\alpha=1,2,3\),
\[
\mathbf D_\alpha(V,\nabla V,\nabla\chi_R)
=2\chi_R\,Y_\alpha V.
\]
Each coefficient is \(L^1(B_{2R})\): \(3\chi_R|Y|^{-p}|V|\,V\in
L^{3/2}(B_{2R})\subset L^1(B_{2R})\) from \(p<3\) and \(|Y|^{-p}|V|^3\in L^1\), and
\(2\chi_RY_\alpha V\in L^2(B_{2R})\subset L^1(B_{2R})\) by compact support.
\[
\frac d{d\tau}G_\alpha(\tau)
=\int_{\RR^3}\mathbf D_\alpha(V,\nn V,\nabla\chi_R)\cdot\partial_\tau V\,\dd Y,
    \]
Since \(\mathbf D_\alpha(V,\nn V,\nabla\chi_R)\in L^1(B_{2R})\) and
\(\partial_\tau V\in L^1_{\rm loc}(I_*;L^2(B_{2R}))\),
the right-hand side is in \(L^1_{\rm loc}(I_*)\), so \(G_\alpha\in W^{1,1}_{\rm loc}(I_*)\subset AC_{\rm loc}(I_*)\).
\end{proof}
\end{lemma}

\begin{lemma}[Directional derivatives in the profile variable]\label{paper2:lem:profile-direction}
For \(\mathcal A_\tau=V+Y\cdot\nn V\) and \(\alpha=0,1,2,3\),
\[
D_V\mathcal F_\alpha(1,0;V)[\mathcal A_\tau]
=-\partial_\lambda\mathcal F_\alpha(1,0;V),\qquad
D_V\mathcal F_\alpha(1,0;V)[\partial_kV]
=-\partial_{c_k}\mathcal F_\alpha(1,0;V).
\]
Indeed, from \(V_{\lambda,c}(Y)=\lambda^{-1}V(\lambda^{-1}(Y-c))\),
\[
\partial_\lambda V_{\lambda,c}\big|_{(1,0)}=-(V+Y\cdot\nabla V)=-\mathcal A_\tau,\qquad
\partial_{c_k}V_{\lambda,c}\big|_{(1,0)}=-\partial_kV.
\]
Substituting these into the Gateaux formulas from Proposition~\ref{paper2:prop:gauge-C1} gives the two identities.
\end{lemma}

\begin{theorem}[Modulation system]\label{paper2:thm:mod-sys}
Let \((V,P)\) be the repaired pair on the compact window under consideration.
Writing \(\theta:=(a,b)\), define
\[
\mathcal A_\tau(Y):=V(\tau,Y)+Y\cdot\nn V(\tau,Y).
\]
For \(\alpha=0,1,2,3\), define
\[
\mathcal R_\alpha(\tau):=
-\int_{\RR^3}\mathbf D_\alpha\bigl(V,\nn V,\nabla\chi_R\bigr)
\cdot\bigl((V\cdot\nn)V+\nabla P-\nu\Delta V\bigr)\,\dd Y.
\]
Then for a.e. \(\tau\in I_*\),
\[
\partial_\lambda\mathcal F_\alpha(1,0;V)\,a(\tau)
+\sum_{k=1}^3\partial_{c_k}\mathcal F_\alpha(1,0;V)\,b_k(\tau)
=\mathcal R_\alpha(\tau),
\]
Equivalently,
\[
J(\tau)\theta(\tau)=\mathcal R(\tau).
\]
Hence
\[
\theta(\tau)=J(\tau)^{-1}\mathcal R(\tau)\in L^1_{\rm loc}(I_*)^4,
\qquad |\theta(\tau)|\le \kappa\,|\mathcal R(\tau)|.
\]
\begin{proof}
Set \(G_\alpha(\tau):=\mathcal F_\alpha(1,0;V(\tau))\). By Theorem~\ref{paper2:thm:repaired-existence},
\(G_\alpha(\tau)=0\) a.e.; by Lemma~\ref{paper2:lem:diff-f}, \(G_\alpha\in W^{1,1}_{\rm loc}\).
Hence for a.e. \(\tau\),
\[
0=\frac{d}{d\tau}G_\alpha(\tau)
=D_V\mathcal F_\alpha(1,0;V(\tau))\,[\partial_\tau V(\tau)].
\]
Insert the equation from Theorem~\ref{paper2:thm:ren-NS},
\[
\partial_\tau V=-(V\cdot\nn)V-\nabla P+\nu\Delta V-a\mathcal A_\tau-b\cdot\nn V.
\]
This gives
\[
0=-D_V\mathcal F_\alpha\bigl[(V\cdot\nn)V+\nabla P-\nu\Delta V\bigr]
-a\,D_V\mathcal F_\alpha[\mathcal A_\tau]
-\sum_{k=1}^3b_k\,D_V\mathcal F_\alpha[\partial_kV].
\]
By Lemma~\ref{paper2:lem:diff-f}, for any profile direction \(W\),
\[
D_V\mathcal F_\alpha(1,0;V)[W]=\int_{\RR^3}\mathbf D_\alpha(V,\nn V,\nabla\chi_R)\cdot W\,\dd Y.
\]
By Lemma~\ref{paper2:lem:profile-direction},
\[
D_V\mathcal F_\alpha[\mathcal A_\tau]=-\partial_\lambda\mathcal F_\alpha(1,0;V),\qquad
D_V\mathcal F_\alpha[\partial_kV]=-\partial_{c_k}\mathcal F_\alpha(1,0;V).
\]
Substituting into the previous line yields the linear system
\[
\partial_\lambda\mathcal F_\alpha(1,0;V)\,a(\tau)
+\sum_{k=1}^3\partial_{c_k}\mathcal F_\alpha(1,0;V)\,b_k(\tau)
=\mathcal R_\alpha(\tau),
\]
which is equivalent to \(J(\tau)\theta(\tau)=\mathcal R(\tau)\).
By definition of \(\mathbf D_\alpha\), each \(\mathcal R_\alpha\in L^1_{\rm loc}(I_*)\), and
Proposition~\ref{paper2:thm:quant-jac} gives \(|J^{-1}(\tau)|\le K\) a.e., hence
\[\theta(\tau)=J^{-1}(\tau)\mathcal R(\tau),\quad 
\theta\in L^1_{\rm loc}(I_*),\quad |\theta|\le K|\mathcal R|\text{ a.e.}.\]
\end{proof}
\end{theorem}

\begin{corollary}[Regularity of \(a,b\)]
For the repaired pair of Theorem~\ref{paper2:thm:mod-sys}, \(a,b\in L^1_{\rm loc}(I_*)\).  
If in addition \(\mathcal R\in L^\infty_{\rm loc}(I_*)\), then \(a,b\in
L^\infty_{\rm loc}(I_*)\).
\begin{proof}
By Theorem~\ref{paper2:thm:mod-sys}, \(\theta=(a,b)=J^{-1}\mathcal R\) a.e., and
\(J^{-1}\in L^\infty_{\rm loc}(I_*)\) by Theorem~\ref{paper2:thm:quant-jac}. The first claim is therefore
immediate from \(\mathcal R\in L^1_{\rm loc}\), while the second follows if \(\mathcal R\in L^\infty_{\rm loc}\).
\end{proof}
\end{corollary}

\subsection{Repaired suitable structure in renormalized variables}\label{paper2:sec:suitable}


\begin{proposition}[Transport of suitable local energy inequality]\label{paper2:prop:ren-suitable}
Suppose \((u,p)\) is suitable on \(\mathcal Q^{\rm ph}_{2R}\), and
\((V,P)\) is obtained by the repaired map \(\Phi\).
Set
\[
\tilde a(\tau):=-\Lambda(\tau)^{-1}\Lambda'(\tau),\qquad
\tilde b(\tau):=-\Lambda(\tau)^{-1}X'(\tau).
\]
Then for every nonnegative \(\phi\in C_c^\infty(Q_R^*)\) and \(\tau_1<\tau_2\),
\[
\begin{aligned}
\mathscr E_\phi(V,P)
&:=\frac12|V|^2\partial_\tau(\phi^2)
+\frac{\nu}{2}|V|^2\Delta_Y\phi^2
+\frac12|V|^2\,V\cdot\nn(\phi^2)\\
&\quad +(P_{\rm loc}+H)\,V\cdot\nn(\phi^2)
+\frac{\tilde a}{2}|V|^2\phi^2
+\frac{\tilde a}{2}|V|^2\,Y\cdot\nn(\phi^2)
+\frac12|V|^2\,\tilde b\cdot\nn(\phi^2).
\end{aligned}
\]
\begin{multline*}
\frac12\int_{B_R}|V(\tau_2)|^2\phi(\cdot,\tau_2)^2
+\nu\int_{\tau_1}^{\tau_2}\!\!\int_{B_R}|\nn V|^2\phi^2\,dY\,d\tau
\le \frac12\int_{B_R}|V(\tau_1)|^2\phi(\cdot,\tau_1)^2\\
 +\int_{\tau_1}^{\tau_2}\!\!\int_{B_R}\mathscr E_\phi(V,P)\,dY\,d\tau.
\end{multline*}
The right-hand side is controlled on compact windows by
\(\|\tilde a\|_{L^1},\|\tilde b\|_{L^1},M\), and the local bounds from
Theorem~\ref{paper2:thm:pressure-decomp}, and support bounds from Section~\ref{paper2:sec:domains}.
Hence \((V,P)\) is locally suitable on \(Q_R^*\).
\begin{proof}
Apply the physical local energy inequality on \(\mathcal Q^{\rm ph}_{2R}\): for
nonnegative \(\psi\in C_c^\infty(\mathcal Q^{\rm ph}_{2R})\),
    \[
    \frac12\int |u|^2\psi\big|_{t_1}^{t_2}
    +\nu\int_{t_1}^{t_2}\!\!\int |\nabla_xu|^2\psi\,dxdt
    \le\int_{t_1}^{t_2}\!\!\int\Bigl(
    \frac12|u|^2(\partial_t\psi+\nu\Delta_x\psi)
    +\left(\frac{|u|^2}{2}u+p\,u\right)\!\cdot\!\nabla_x\psi\Bigr)\,dxdt.
    \]
Choose
    \[
    \psi(x,t):=\Lambda(t)^{-1}\phi\!\left(\frac{x-X(t)}{\Lambda(t)},\tau(t)\right)^2,\qquad
    \tau=\tau(t),\ d\tau/dt=\Lambda^{-2},
    \]
let \(t_i=t(\tau_i)\), \(i=1,2\). 
    By Proposition~\ref{paper2:prop:prelim-frame} and Section~\ref{paper2:sec:domains}, \(\psi\in
    C_c^\infty(\mathcal Q^{\rm ph}_{2R})\).
\[
\nabla_x\psi=\Lambda^{-2}\nabla_Y\phi^2,\quad
\Delta_x\psi=\Lambda^{-3}\Delta_Y\phi^2,\quad
\partial_t\psi
=-\Lambda^{-4}\Lambda'\phi^2+\Lambda^{-3}\partial_\tau(\phi^2)-\Lambda^{-4}(X'+\Lambda'Y)\cdot\nabla_Y\phi^2.
\]
Using \(u=\Lambda^{-1}V\), \(p=\Lambda^{-2}P\), \(\nabla_xu=\Lambda^{-2}\nabla_YV\),
\(dx\,dt=\Lambda^{5}\,dY\,d\tau\), each term in the physical inequality becomes
\[
\frac12\int |u|^2\psi\big|_{t_1}^{t_2}
=\frac12\int |V|^2\phi^2\big|_{\tau_1}^{\tau_2},
\]
\[
\nu\int |\nabla_xu|^2\psi\,dxdt
=\nu\int |\nabla V|^2\phi^2\,dY\,d\tau,
\]
\[
\frac12\int |u|^2\partial_t\psi\,dxdt
=\frac12\int |V|^2\left(\partial_\tau(\phi^2)+\tilde a\phi^2+\tilde a\,Y\cdot\nabla_Y\phi^2+\tilde b\cdot\nabla_Y\phi^2\right)\,dY\,d\tau,
\]
\[
\frac{\nu}{2}\int |u|^2\Delta_x\psi\,dxdt
=\frac{\nu}{2}\int |V|^2\Delta_Y\phi^2\,dY\,d\tau,
\]
\[
\int\left(\frac{|u|^2}{2}u+p\,u\right)\cdot\nabla_x\psi\,dxdt
=\int\left(\frac12|V|^2V+(P_{\rm loc}+H)V\right)\cdot\nabla_Y\phi^2\,dY\,d\tau.
\]
Substituting these identities and grouping terms gives the proposition's right-hand
side and coefficients \(\tilde a,\tilde b\).
\end{proof}
\end{proposition}
\subsection{Stability under subsequences}\label{paper2:sec:stability}
\begin{proposition}\label{paper2:prop:stability}
Let \((u^{(n)},p^{(n)})\) satisfy \(\mathrm{SC}(R_*,r_*,M,\eta,\kappa)\) with the same
data parameters and let \((\Lambda^{(n)},X^{(n)},V^{(n)},P^{(n)},a^{(n)},b^{(n)})\)
be obtained from retained quantitatively nondegenerate repaired-gauge selections on the same \(I_*\),
with the constants in Proposition~\ref{paper2:thm:quant-jac} uniform in \(n\). After extraction \(n_j\to\infty\), there exist
subsequences and limits \((\Lambda,X,V,P,a,b)\) such that
\[
\Lambda^{(n_j)}\to\Lambda\ \text{in }L^1_{\rm loc},\quad
X^{(n_j)}\to X\ \text{in }L^1_{\rm loc},
\]
\[
\Lambda^{(n_j)\prime}\rightharpoonup\Lambda'\ \text{in }L^1_{\rm loc},\quad
X^{(n_j)\prime}\rightharpoonup X'\ \text{in }L^1_{\rm loc},
\]
\[
V^{(n_j)}\to V\ \text{in }L^2_{\rm loc}(I_*;L^2_{\rm loc}),\quad
V^{(n_j)}\rightharpoonup V\ \text{in }L^2_{\rm loc}(I_*;H^1_{\rm loc}),
\]
    \[
    P^{(n_j)}\rightharpoonup P\ \text{in }L^{3/2}_{\rm loc}(Q_R^*),\quad
    (a^{(n_j)},b^{(n_j)})\rightharpoonup(a,b)\ \text{in }L^1_{\rm loc}(I_*)^4.
    \]
    Moreover \(a=-\Lambda'\), \(b=-X'\) a.e.\ on \(I_*\), and the limit obeys all conclusions below.
\begin{proof}
From \(\mathrm{SC}\) and the retained repaired-gauge selections, \((V^{(n)},P^{(n)})\) are uniformly bounded in
\(L^\infty_\tau L^2(B_{2R})\cap L^2_\tau H^1(B_{2R})\) and \(L^{3/2}_{\rm loc}(Q_R^*)\).
Theorem~\ref{paper2:thm:quant-jac} gives uniform invertibility bounds for \(J^{(n)}\), and
Theorem~\ref{paper2:thm:mod-sys} gives \((a^{(n)},b^{(n)})\) uniformly in \(L^1_{\rm loc}\) (since
\(\mathcal R^{(n)}\in L^1_{\rm loc}\)). Hence \(\Lambda^{(n)}\) and \(X^{(n)}\) are uniformly
bounded in \(W^{1,1}_{\rm loc}\), so by BV compactness and 1D subsequences,
\[
\Lambda^{(n_j)}\to\Lambda\ \text{in }L^1_{\rm loc},\qquad
X^{(n_j)}\to X\ \text{in }L^1_{\rm loc},
\]
and \(a^{(n_j)}\rightharpoonup -\Lambda'\), \(b^{(n_j)}\rightharpoonup -X'\) in \(L^1_{\rm loc}\).
From the renormalized equation,
\[
\partial_\tau V^{(n)}\in
L^1_\tau(H^{-1}_{\rm loc})+L^{4/3}_\tau(L^{12/11}_{\rm loc}),
\]
uniformly on compact windows. Aubin--Lions compactness on bounded balls gives
\[
V^{(n_j)}\to V\ \text{in }L^2_{\rm loc}(I_*;L^2_{\rm loc}),\quad
V^{(n_j)}\rightharpoonup V\ \text{in }L^2_{\rm loc}(I_*;H^1_{\rm loc}).
\]
By Calderón--Zygmund and Theorem~\ref{paper2:thm:pressure-decomp}, \(P^{(n_j)}\rightharpoonup P\) in
\(L^{3/2}_{\rm loc}(Q_R^*)\), and with normalized pressure split \(P=P_{\rm loc}+H\),
the same for \(P_{\rm loc}\) and \(H\). Passing to limits in
\(\mathcal F(1,0;V^{(n)}(\tau))=0\), Theorem~\ref{paper2:thm:ren-NS}, and the linear modulation system
yields the corresponding limit constraints and equations.

\end{proof}
\end{proposition}

\subsection{Final representation theorem}\label{paper2:sec:final}

\begin{theorem}[Chart and repaired-gauge representation]\label{paper2:thm:A}
Assume \(\mathrm{SC}(R_*,r_*,M,\eta,\kappa)\), and suppose the branch lies in the
retained repaired-gauge regime with selection \((\lambda,c)\) on
\(J=[\tau_0-\ell,\tau_0+\ell]\subset I_*\). Then there are fields \((V,P)\) on \(Q_{2R}^J\), with
\((\Lambda,X)\in AC(J),\ a=-\Lambda',\ b=-X'\in L^1_{\rm loc}(J)\), such that
\[
u(x,t)=\Lambda(\tau)^{-1}V\!\left(\frac{x-X(\tau)}{\Lambda(\tau)},\tau\right),\quad
\frac{dt}{d\tau}=\Lambda(\tau)^2,
\]
and \(\mathcal F(1,0;V(\tau))=0\) for \(\tau\in J\).
\begin{proof}
From \textup{(SC5)} and Proposition~\ref{paper2:prop:prelim-frame}, \((\Lambda_{\rm pre},X_{\rm pre})\)
are AC on \(I_*\), so \(V_{\rm pre}\) is defined on \(Q_{2R}^*\).
The retained repaired-gauge regime supplies
\((\lambda,c)\in AC(J;(0,\infty)\times\mathbb R^3)\) and
\(\mathcal F(1,0;V_{\rm pre}^{\lambda,c})=0\) a.e.
By Definition~\ref{paper2:def:repaired-chart} and Theorem~\ref{paper2:thm:repaired-existence}, setting
\[
\Lambda(\tau)=\Lambda_{\rm pre}(\tau)\lambda(\tau),\qquad
X(\tau)=X_{\rm pre}(\tau)+\Lambda_{\rm pre}(\tau)c(\tau),\qquad
V(\tau)=V_{\rm pre}^{\lambda(\tau),c(\tau)},
\]
and \(P(\tau)=P_{\rm pre}^{\lambda(\tau),c(\tau)}\),
gives \( \Lambda,X\in AC(J)\) and \(a=-\Lambda',\ b=-X'\in L^1_{\rm loc}(J)\), and
\[
\mathcal F(1,0;V(\tau))=0\ \text{for a.e. }\tau\in J.
\]
Finally \(u=\Lambda^{-1}V(\cdot,\tau)\circ\Phi^{-1}\) by definition of the repaired map, hence the
physical identity with \(\frac{dt}{d\tau}=\Lambda^2\) follows.
\end{proof}
\end{theorem}

\begin{theorem}[Pressure and modulation]\label{paper2:thm:B}
Assume, in addition to \autoref{paper2:thm:A}, that the retained repaired-gauge
selection is quantitatively nondegenerate in the sense of Proposition~\ref{paper2:thm:quant-jac}.
Then \((V,P)\) satisfies the renormalized Navier--Stokes
equation with the sign convention of Theorem~\ref{paper2:thm:ren-NS}; there is a
decomposition \(P=P_{\rm loc}+H\) on \(Q_R^*\) as in
Theorem~\ref{paper2:thm:pressure-decomp}; and \((a,b)\) is the unique solution of
\[
J(\tau)\binom{a(\tau)}{b(\tau)}=\mathcal R(\tau),\qquad
\binom{a}{b}\in L^1_{\rm loc}(J)^4,
\]
with bounds \(|(a,b)|\le K|\mathcal R|\), where \(K\) depends on
\(R,p,\Theta_0,m_0,C_u,\varepsilon_0,\chi_R\).
\begin{proof}
Equation and sign convention are inherited directly from Theorem~\ref{paper2:thm:ren-NS} applied to
\((\Lambda,X)\) from Theorem~\ref{paper2:thm:A} and the definition of \(V\).
By Theorem~\ref{paper2:thm:pressure-decomp} applied at each time slice, the decomposition on \(Q_R^J\)
is defined and unique once \(\fint_{B_R}H=0\) is imposed.
The modulation identity is Theorem~\ref{paper2:thm:mod-sys}; uniqueness follows from invertibility of
\(J\) (Theorem~\ref{paper2:thm:quant-jac}) and the pointwise identity \((a,b)=J^{-1}\mathcal R\).
Uniform bounds on \((a,b)\) are obtained by the bound on \(J^{-1}\) there and the stated bound on
\(\mathcal R\).
\end{proof}
\end{theorem}

\begin{theorem}[Compact-window representation theorem]\label{paper2:thm:C}
Under the assumptions of \cref{paper2:thm:A,paper2:thm:B}, for every compact
window \(J\subset I_*\) and each fixed \(R>0\), setting
\(Q_R^J:=B_R\times J\),
one has:
\begin{enumerate}
\item the exact renormalized equation on \(Q_{2R}^J\),
\item the four repaired-gauge constraints \(\mathcal F(1,0;V(\tau))=0\),
\item bounds on \(a,b\in L^1_{\rm loc}(J)\) (and \(L^\infty_{\rm loc}\) under
  \(\mathcal R\in L^\infty_{\rm loc}\)),
\item local pressure decomposition \(P=P_{\rm loc}+H\) with \(L^{3/2}\) local bounds and
  harmonic remainder controls,
	\item renormalized local suitable energy inequality on \(Q_R^J\),
	\item stability under subsequences in the topologies above.
\end{enumerate}
The later reductions use only these conclusions.
\begin{proof}
From Theorem~\ref{paper2:thm:A} and the choice of \(J\subset I_*\), all quantities are defined on \(Q_{2R}^J\).
Items (1) and (2) are Theorems~\ref{paper2:thm:ren-NS} and~\ref{paper2:thm:repaired-existence};
item (3) is Theorem~\ref{paper2:thm:mod-sys} together with the inverse bound from
Proposition~\ref{paper2:thm:quant-jac}; item (4) is Theorem~\ref{paper2:thm:pressure-decomp};
item (5) is Proposition~\ref{paper2:prop:ren-suitable}; item (6) is
Proposition~\ref{paper2:prop:stability}. Combining these results gives the
stated compact-window representation.
\end{proof}
\end{theorem}

\begin{remark}
The preceding results separate two alternatives:
\begin{enumerate}
\item the nondegenerate single-core repaired-gauge regime, where the representation theorems apply;
\item gauge degeneracy (\(\det J=0\)) or compact-structure failure, which belongs
to the multibubble, cascade, or scale-rigidity alternatives.
\end{enumerate}
\end{remark}

\clearpage
\section{Multibubble and Cascade Exclusion}\label{sec:multibubble-cascade-exclusion}
\medskip
\noindent\textbf{Analytic inputs.}

The local reduction theorem uses the following analytic inputs from the
preceding sections, with profile decompositions and decoupling estimates in
their standard critical-space sense \cite{Gallagher2001,BrezisLieb1983}.  The
remaining work is the multibubble, cascade, and terminal-frame decoupling.
\begin{enumerate}[label=\textup{(\roman*)}]
\item From \cref{sec:single-core-criterion}, the compact single-core criterion \cref{paper1:thm:single-core-criterion}: a
compact single-core finite-cost Type II branch with positive local CKN
concentration, nondegenerate scale-center selection, pressure reconstruction,
bounded compact-cylinder suitable estimates, bounded modulation, and vanishing
localized dissipation is excluded.
\item From \cref{sec:repaired-gauge-construction}, the renormalized equation \cref{paper2:thm:ren-NS}, pressure decomposition \cref{paper2:thm:pressure-decomp}, modulation system \cref{paper2:thm:mod-sys}, and repaired-gauge representation \cref{paper2:thm:A,paper2:thm:B,paper2:thm:C}: the renormalized equation,
pressure decomposition modulo functions of time, bounded modulation
coefficients, renormalized suitable structure, and stability of these structures
under subsequences.
\item From \cref{sec:rough-core-exclusion}, the compact-cylinder Caccioppoli estimate \cref{paper4:thm:caccioppoli} and rough-core exclusion \cref{paper4:cor:rough-core}.
\item From \cref{sec:scale-collapse-cost}, the collapse cost exclusion \cref{paper5:thm:collapse-core-floor,paper5:thm:cost-exclusion} and autonomous reduced-limit theorem \cref{paper5:thm:autonomous-limit} for scale-rigid or scale-collapsing branches.
\end{enumerate}

\subsection{Local setting and active frames}

\medskip
\noindent\textit{Normalized local Navier--Stokes setting.}
All quantities are written in rescaled local variables \((x,t)\in Q_R=B_R\times(-R^2,0)\)
with viscosity \(\nu=1\) (critical scaling has absorbed \(\nu\)).  A local pair
\((u_n,p_n)\) on such a cylinder satisfies
\[
  \partial_tu_n-\Delta u_n+(u_n\cdot\nabla)u_n+\nabla p_n=0,\qquad \nabla\cdot u_n=0,
\]
and pressure is normalised through spatial averaging
\[
  (p_n)_{B_\rho}(t):=|B_\rho|^{-1}\int_{B_\rho}p_n(x,t)\,dx
\]
to factor out time-dependent constants.

For $r>0$, $x_0\in\mathbb R^3$, and $t_0\in\mathbb R$, let
\[
  Q_r(x_0,t_0):=B_r(x_0)\times(t_0-r^2,t_0),
  \qquad
  Q_r:=Q_r(0,0)=B_r\times(-r^2,0).
\]

Let $(u_n,p_n)$ be parabolic rescalings around a potential singular point with
positive local Caffarelli--Kohn--Nirenberg concentration on compact cylinders.
Concretely, for a base triple \((x_n,t_n,\lambda_n)\),
\[
  u_n(y,s)=\lambda_n u(x_n+\lambda_n y,\ t_n+\lambda_n^2 s),\qquad
  p_n(y,s)=\lambda_n^2 p(x_n+\lambda_n y,\ t_n+\lambda_n^2 s).
\]
For each $n$ and $R>0$, define the scaling-invariant local quantities
\[
  A_n(R):=R^{-1}\operatorname*{ess\,sup}_{-R^2<t<0}\int_{B_R}|u_n(x,t)|^2\,dx,
\]
\[
  E_n(R):=R^{-1}\int_{Q_R}|\nabla u_n(x,t)|^2\,dx\,dt,
\]
\[
  C_n(R):=R^{-2}\int_{Q_R}|u_n(x,t)|^3\,dx\,dt,
  \qquad
  D_n(R):=R^{-2}\int_{-R^2}^0\int_{B_R}|p_n-(p_n)_{B_R}(t)|^{3/2}\,dx\,dt,
\]
where
\[
  (p_n)_{B_R}(t):=|B_R|^{-1}\int_{B_R}p_n(x,t)\,dx.
\]

For an active frame \((x_n,\lambda_n)\), \(\lambda_n>0\), define the critical rescaling
\[
  \mathcal T_{x_n,\lambda_n}f(y):=\lambda_n f(x_n+\lambda_n y).
\]
If \(\phi\in L^3_\sigma(\mathbb R^3)\) is observed in this frame, its contribution is
\[
  \phi_n(y)=\mathcal T_{x_n,\lambda_n}\phi(y)=\lambda_n\phi(x_n+\lambda_n y).
\]
The critical Jacobian identity is exact:
\[
  \|\mathcal T_{x_n,\lambda_n}f\|_{L^3(E)}^3
  =\int_{x_n+\lambda_nE}|f(x)|^3\,dx
\]
for every measurable \(E\subset\mathbb R^3\).  In particular
\[
  \|\mathcal T_{x_n,\lambda_n}f\|_{L^3(\mathbb R^3)}^3
  =\|f\|_{L^3(\mathbb R^3)}^3,
\]
so \(\mathcal T_{x_n,\lambda_n}\) is an isometry on \(L^3(\mathbb R^3)\).

\begin{definition}[Bounded selected-window limit]
A selected-window subsequence has a bounded selected-window limit on a compact
cylinder $Q_R$ if there exists $u_\infty\in L^\infty(Q_R)$ such that, after
extracting a subsequence,
\[
  \forall\varepsilon>0\ \exists N:\ \forall n\ge N,\ \|u_n-u_\infty\|_{L^3(Q_R)}<\varepsilon.
\]
The limit is nonzero if $u_\infty\not\equiv0$ on $Q_R$.
\end{definition}

\begin{definition}[Local Type II sequence]
A positive local concentration sequence is called a local Type II sequence if every
selected-window rescaling has zero bounded selected-window limit on every compact
cylinder \(Q_R\):
\[
  \forall R>0,\qquad \lim_{n\to\infty}\|u_n\|_{L^3(Q_R)}=0.
\]
\end{definition}

\begin{definition}[Active frame]\label{paper3:def:active-frame}
Fix a concentrating time sequence \(t_n\uparrow T\) and a compact time interval
\(I\Subset(-\infty,0]\).  An active frame is a tuple
\[
  \mathfrak F^j=\bigl(x^j_n,\lambda^j_n,I,U^j_n,P^j_n\bigr),
  \qquad \lambda^j_n\downarrow0,
\]
where \(x^j_n\in\mathbb R^3\) and
\[
  U^j_n(y,s):=\lambda^j_n u(x^j_n+\lambda^j_n y,t_n+(\lambda^j_n)^2s),
  \qquad
  P^j_n(y,s):=(\lambda^j_n)^2p(x^j_n+\lambda^j_n y,t_n+(\lambda^j_n)^2s)
\]
on \(B_{2R}\times I\) for every fixed compact cylinder under consideration.
The frame is active at threshold \(\eta_*>0\) if for some fixed \(R_0<\infty\)
\[
  \liminf_{n\to\infty}\int_{Q_{R_0}}|U^j_n|^3\,dy\,ds\ge \eta_*^3.
\]
\end{definition}

\begin{definition}[Relations between active frames]\label{paper3:def:frame-relations}
Let \(\mathfrak F^i\) and \(\mathfrak F^j\) be two active frames.  After passing
to a subsequence exactly one of the following parameter regimes holds.
\begin{enumerate}[label=\textup{(\roman*)}]
\item They are \emph{same-point comparable-scale} if
\[
  0<c\le\frac{\lambda^i_n}{\lambda^j_n}\le C<\infty,\qquad
  \frac{|x^i_n-x^j_n|}{\max(\lambda^i_n,\lambda^j_n)}\le C
\]
for all large \(n\).
\item They form a \emph{same-point cascade} if, after relabeling,
\[
  \frac{\lambda^i_n}{\lambda^j_n}\to0,\qquad
  \frac{|x^i_n-x^j_n|}{\lambda^j_n}=O(1).
\]
The frame \(i\) is then the inner frame and \(j\) is an outer frame.
\item They are \emph{separated-point} if
\[
  \frac{|x^i_n-x^j_n|}{\lambda^i_n}\to\infty
  \quad\text{and}\quad
  \frac{|x^i_n-x^j_n|}{\lambda^j_n}\to\infty .
\]
\end{enumerate}
\end{definition}

\begin{definition}[Terminal frame and perturbative remainder]\label{paper3:def:terminal-frame}
Given a finite active family, a terminal frame is chosen by selecting an active
scale with minimal order:
\[
  \lambda^{j_*}_n\le C_j\lambda^j_n
  \qquad\text{for every active }j
\]
after passing to a subsequence.  Comparable frames at the same physical point
are grouped before this choice; if several grouped frames remain tied, one is
chosen by a fixed lexicographic ordering of their labels.  A remainder \(r_n\)
in a terminal frame is perturbative on \(B_{2R}\times I\) if
\[
  \|r_n\|_{L^\infty_I L^3(B_{2R})}\to0
\]
and its equation residual and pressure satisfy
\[
  \|f_n\|_{L^1_IH^{-1}(B_R)}+\|\nabla\pi^r_n\|_{L^1_IH^{-1}(B_R)}\to0
\]
for every compact \(B_R\times I\).
\end{definition}

\begin{definition}[Local multibubble/cascade branch]\label{paper3:def:local-multibubble}
A local multibubble/cascade branch is a positive-concentration local Type II
branch for which at least one of the following alternatives occurs:
\begin{enumerate}[label=\textup{(\roman*)}]
\item the compact-cylinder suitable estimates $A_n+E_n+C_n+D_n$ are unbounded on some
  fixed rescaled cylinder;
\item positive CKN mass splits into two or more comparable active cores;
\item a strict same-point scale cascade appears inside bounded rescaled cylinders;
\item separated physical points remain active in different local rescaled frames;
\item the localized scale or centering selection degenerates (no unique core
  selected).
\end{enumerate}
\end{definition}

\begin{definition}[Active profile family]\label{paper3:def:active-family}
An active profile family is a family
\[
  \{(\phi^j,x_n^j,\lambda_n^j)\}_{j\in J},
  \qquad \phi^j\in L^3_\sigma(\mathbb R^3),
\]
with finite or countable index set $J$, extracted along a positive local
concentration sequence.  In physical variables write
\[
  \Lambda_{\lambda,x_0}\phi(x):=\lambda^{-1}\phi\left(\frac{x-x_0}{\lambda}\right).
\]
For every finite \(J_0\subset J\) the family gives an expansion
\[
  u(\cdot,t_n)=\sum_{j\in J_0}\Lambda_{\lambda^j_n,x^j_n}\phi^j+r^{J_0}_n
\]
on the compact terminal windows under consideration.  After the finite active
subfamily has been selected, the corresponding remainder satisfies the
perturbative bounds of \cref{paper3:def:terminal-frame}.  The parameters are
considered only after grouping same-point comparable-scale profiles into
compound profiles as in \cref{paper3:def:compound-profile}.  A profile is active if
\[
  \|\phi^j\|_{L^3(\mathbb R^3)}\ge\eta_*>0.
\]
The family has finite critical mass when
\[
  \sum_{j\in J}\|\phi^j\|_{L^3(\mathbb R^3)}^3\le M_*<\infty .
\]
It captures positive local concentration if every compact terminal-window
sequence with
\[
  \liminf_{n\to\infty}\int_{Q_R}|u_n|^3>0
\]
either contains an active profile observed in that window or is absorbed into
the selected terminal frame; equivalently, no positive terminal-window
concentration remains in the perturbative remainder.
\end{definition}

\begin{definition}[Compound profile]\label{paper3:def:compound-profile}
Let \(J_k\subset J\) be a same-point comparable-scale class.  Choose one
representative \((x_n^{j_k},\lambda_n^{j_k})\).  After passing to a subsequence,
\[
  \rho_j:=\lim_{n\to\infty}\frac{\lambda^j_n}{\lambda^{j_k}_n}\in(0,\infty),
  \qquad
  z_j:=\lim_{n\to\infty}\frac{x^j_n-x^{j_k}_n}{\lambda^{j_k}_n}\in\mathbb R^3 .
\]
The compound profile in the representative frame is the \(L^3_\sigma\)-sum
\[
  \Phi^k(y):=\sum_{j\in J_k}\rho_j^{-1}
  \phi^j\left(\frac{y-z_j}{\rho_j}\right),
\]
whenever \(J_k\) is finite; finiteness follows for active classes from
\cref{paper3:lem:finite-active-count}.  If \(\|\Phi^k\|_{L^3}<\varepsilon_\star\), where
\(\varepsilon_\star\) is the critical small-data threshold, the class is
perturbative.  Otherwise \(\Phi^k\) is treated as one active object in its
representative frame.
\end{definition}

\begin{lemma}[Finite active profile count]\label{paper3:lem:finite-active-count}
Every active profile family with finite critical mass contains only finitely many
active profiles.  More precisely,
\[
  \#J\le M_*\eta_*^{-3}.
\]
\end{lemma}

\begin{proof}
For each active profile $j\in J$,
\[
  \|\phi^j\|_{L^3(\mathbb R^3)}^3\ge\eta_*^3.
\]
Let $J_0\subset J$ be finite. Summing over $J_0$ gives
\[
  \#J_0\,\eta_*^3
  \le
  \sum_{j\in J_0}\|\phi^j\|_{L^3(\mathbb R^3)}^3
  \le
  \sum_{j\in J}\|\phi^j\|_{L^3(\mathbb R^3)}^3
  \le M_*.
\]
Hence $\#J_0\le M_*\eta_*^{-3}$.  Taking the supremum over all finite
subsets $J_0$ yields the same bound for the cardinality of $J$, so $J$ is finite
and
\[
  \#J\le M_*\eta_*^{-3}.
\]
\end{proof}

\begin{lemma}[Exhaustive parameter partition]\label{paper3:lem:parameter-partition}
Let \(\mathfrak F^i,\mathfrak F^j\) be two active frames.  After passing to a
subsequence, either the pair is same-point comparable-scale, same-point cascade,
or separated-point in the sense of \cref{paper3:def:frame-relations}; if none of these
relations holds in a selected compact frame, then one of the two contributions
is locally invisible there in \(L^3\).
\end{lemma}

\begin{proof}
Pass to a subsequence so that
\[
  \frac{\lambda^i_n}{\lambda^j_n}\to \ell\in[0,\infty]
\]
in the extended half-line.  If \(0<\ell<\infty\), then either
\[
  \frac{|x^i_n-x^j_n|}{\max(\lambda^i_n,\lambda^j_n)}
\]
is bounded along a further subsequence, giving the same-point comparable-scale
case, or it tends to infinity, giving separated points because the two scales
are comparable.
\par
If \(\ell=0\), then either
\[
  \frac{|x^i_n-x^j_n|}{\lambda^j_n}
\]
is bounded, giving a same-point cascade with \(i\) inner and \(j\) outer, or it
tends to infinity.  In the latter case the \(j\)-profile is locally invisible in
the \(i\)-frame by \cref{paper3:lem:separated-vanish}.  The case \(\ell=\infty\) is the
same after interchanging \(i\) and \(j\).  These alternatives exhaust all
subsequential limits of the scale ratio and normalized center distance.
\end{proof}

\begin{theorem}[Minimal bad-configuration extraction]\label{paper3:thm:minimal-bad-extraction}
Let a positive local Type II branch admit an active profile family with finite
critical mass, perturbative remainder in the sense of
\cref{paper3:def:terminal-frame}, and concentration capture in the sense of
\cref{paper3:def:active-family}.  If the branch is not already a single active
terminal frame satisfying the single-core hypotheses, then after passing to a
subsequence one of the following occurs:
\begin{enumerate}[label=\textup{(\roman*)}]
\item compact-cylinder suitable control fails, and
\cref{paper3:lem:compactness-failure} produces a bounded-cylinder concentration core;
\item there are at least two active frames in a same-point comparable-scale
class, hence a compound profile in the sense of \cref{paper3:def:compound-profile};
\item there is a same-point cascade;
\item there are separated active physical points;
\item the scale-center selection of the selected compound or terminal frame is
degenerate.
\end{enumerate}
In every case the positive local concentration is carried by an active frame or
compound profile, not by the perturbative remainder.
\end{theorem}

\begin{proof}
If compact-cylinder suitable control fails, item \((i)\) follows from
\cref{paper3:lem:compactness-failure}.  Assume therefore that the compact-cylinder
family is bounded on the windows under consideration.
\par
By concentration capture, positive local concentration cannot remain solely in
the perturbative remainder.  Hence either one selected active frame carries the
terminal-window concentration, or another active frame remains non-negligible.
If there is one nondegenerate selected active frame and no competing active
object, the branch is in the single-core regime.  Since this case is
excluded by hypothesis, either the selected scale-center map is degenerate,
giving item \((v)\), or at least two active frames remain.
\par
For two active frames, apply \cref{paper3:lem:parameter-partition}.  Same-point
comparable-scale frames give item \((ii)\), same-point scale separation gives
item \((iii)\), and separated points give item \((iv)\).  The remaining
alternative in \cref{paper3:lem:parameter-partition} is local invisibility of one
contribution in the selected compact frame; such a contribution is absorbed
into the perturbative remainder and cannot carry the retained positive local
concentration by the capture property.  This gives the extraction statement.
\end{proof}

\subsection{Elementary local decoupling estimates}

\begin{lemma}[Critical rescalings that escape are locally invisible]\label{paper3:lem:escaping-l3}
Let \(\phi\in L^3(\mathbb R^3)\), let \(R<\infty\), and define
\[
  W_n(y)=\rho_n\phi(z_n+\rho_n y),
\]
where \(\rho_n>0\).  Assume that the sets
\[
  z_n+\rho_nB_R
\]
either have measure tending to zero, or escape to spatial infinity in the sense
that for every compact set \(K\subset\mathbb R^3\),
\[
  (z_n+\rho_nB_R)\cap K=\varnothing
\]
for all sufficiently large \(n\).  Then
\[
  \|W_n\|_{L^3(B_R)}\to0.
\]
\end{lemma}

\begin{proof}
By the critical change of variables,
\[
  \|W_n\|_{L^3(B_R)}^3
  =\int_{B_R}\rho_n^3|\phi(z_n+\rho_n y)|^3\,dy
  =\int_{z_n+\rho_n B_R}|\phi(x)|^3\,dx.
\]
Set
\[
  m_n:=\int_{z_n+\rho_n B_R}|\phi(x)|^3\,dx.
\]
Since \(\|W_n\|_{L^3(B_R)}^3=m_n\), it is enough to show \(m_n\to0\).
\par
Fix \(\varepsilon>0\).  We split into the two structural alternatives.

\smallskip
\noindent\textit{Case 1: shrinking measure.}
If \(|z_n+\rho_nB_R|\to0\), then by absolute continuity of the integral
\(|\phi|^3\in L^1(\mathbb R^3)\) there exists \(\delta>0\) such that
\[
  |E|<\delta\Longrightarrow \int_E|\phi|^3\,dx<\varepsilon.
\]
Choose \(N\) with \(|z_n+\rho_nB_R|<\delta\) for all \(n\ge N\).  Then
\[
  m_n=\int_{z_n+\rho_nB_R}|\phi|^3\,dx<\varepsilon,\qquad n\ge N.
\]

\smallskip
\noindent\textit{Case 2: escape to infinity.}
Assume now that for every compact set \(K\subset\mathbb R^3\),
\((z_n+\rho_nB_R)\cap K=\varnothing\) for all sufficiently large \(n\).  By
absolute continuity at infinity, choose compact \(K\subset\mathbb R^3\) such that
\[
  \int_{\mathbb R^3\setminus K}|\phi|^3<\varepsilon.
\]
For every such compact \(K\), there is \(N(K)\) with
\((z_n+\rho_nB_R)\cap K=\varnothing\) for all \(n\ge N(K)\).  For these
indices,
\[
  z_n+\rho_nB_R \subset \mathbb R^3\setminus K,
\]
and therefore, for \(n\ge N(K)\),
\[
  m_n=\int_{z_n+\rho_nB_R}|\phi|^3\,dx
  \le
  \int_{\mathbb R^3\setminus K}|\phi|^3
  <\varepsilon.
\]
In either case, for every \(\varepsilon>0\), \(m_n<\varepsilon\) for \(n\) sufficiently
large.  Hence \(m_n\to0\), and therefore
\[
  \|W_n\|_{L^3(B_R)}\to0.
\]
\end{proof}

\begin{lemma}[Separated active profiles vanish in the selected frame]\label{paper3:lem:separated-vanish}
Let \((x^p_n,\lambda^p_n)\) and \((x^q_n,\lambda^q_n)\) be two active frames
with separated physical centers in the sense that
\[
  \frac{|x^q_n-x^p_n|}{\lambda^p_n}\to\infty.
\]
If \(\phi^q\in L^3(\mathbb R^3)\), then the \(q\)-profile observed in the
\(p\)-frame converges to zero locally in \(L^3\): for every \(R<\infty\),
\[
  \left\|
  \frac{\lambda^p_n}{\lambda^q_n}
  \phi^q\left(\frac{x^p_n-x^q_n}{\lambda^q_n}
       +\frac{\lambda^p_n}{\lambda^q_n}y\right)
  \right\|_{L^3(B_R)}\to0.
\]
\end{lemma}

\begin{proof}
The displayed function has the form in \cref{paper3:lem:escaping-l3} with
\[
  z_n=\frac{x^p_n-x^q_n}{\lambda^q_n},
  \qquad
  \rho_n=\frac{\lambda^p_n}{\lambda^q_n}.
\]
The corresponding physical set in the \(q\)-profile variables is
\[
  z_n+\rho_nB_R
  =\frac{x^p_n-x^q_n}{\lambda^q_n}
   +\frac{\lambda^p_n}{\lambda^q_n}B_R.
\]
For every \(y\in B_R\),
\[
  |z_n+\rho_n y|
  \ge |z_n|-\rho_nR.
\]
Set \(E_n:=z_n+\rho_nB_R\).  Then
\[
  \operatorname{dist}(E_n,0)
  =|z_n|-\rho_nR
  =\frac{\lambda^p_n}{\lambda^q_n}\left(\frac{|x^q_n-x^p_n|}{\lambda^p_n}-R\right).
\]
Since \(\frac{|x^q_n-x^p_n|}{\lambda^p_n}\to\infty\), there exists \(N_0(R)\) such that
\[
  \frac{|x^q_n-x^p_n|}{\lambda^p_n}-R\ge1,\quad n\ge N_0(R),
\]
and hence for such \(n\), \(\operatorname{dist}(E_n,0)\ge\rho_n\).
\par
There are two cases.
\begin{enumerate}[label=\textup{(\alph*)}]
\item If \(\rho_n\to0\), then
 \[
   |z_n+\rho_nB_R|=\rho_n^3|B_R|\to0.
 \]
 Then the first alternative in \cref{paper3:lem:escaping-l3} applies.
\item Otherwise, after passing to a subsequence, \(\rho_n\ge\rho_\ast>0\).  Then
 \[
   \operatorname{dist}(z_n+\rho_nB_R,0)\ge \rho_\ast\left(\frac{|x^q_n-x^p_n|}{\lambda^p_n}-R\right)\to\infty,
 \]
 so the sets escape every fixed compact subset of \(\mathbb R^3\).  The
 second alternative in \cref{paper3:lem:escaping-l3} applies.
\end{enumerate}
In both subcases the \(L^3(B_R)\)-norm of the displayed expression converges to
zero.
\end{proof}

\begin{lemma}[Strict outer scales vanish in the innermost frame]\label{paper3:lem:outer-vanish}
Let several active profiles have the same physical center and scales
\(\lambda^1_n,\ldots,\lambda^J_n\).  Suppose \(\lambda^1_n\) is an innermost
scale and
\[
  \frac{\lambda^1_n}{\lambda^j_n}\to0
  \qquad (j\ne1).
\]
Then every outer profile \(\phi^j\in L^3(\mathbb R^3)\), viewed in the
\(\lambda^1_n\)-frame, converges to zero in \(L^3(B_R)\) for each fixed
\(R<\infty\).  Any constant ambient velocity arising from the regular part of
the outer flow is absorbed into the translation modulation.
\end{lemma}

\begin{proof}
In the innermost variables, the \(j\)-th outer contribution has the form
\[
  W^j_n(y)=\frac{\lambda^1_n}{\lambda^j_n}
  \phi^j\left(\frac{x^1_n-x^j_n}{\lambda^j_n}
  +\frac{\lambda^1_n}{\lambda^j_n}y\right).
\]
Set
\[
  \rho_n:=\frac{\lambda^1_n}{\lambda^j_n}\to0,
  \qquad
  z_n^j:=\frac{x^1_n-x^j_n}{\lambda^j_n}.
\]
Then
\[
  W^j_n(y)=\rho_n\phi^j\left(z_n^j+\rho_n y\right),
\]
and by the critical change of variables,
\[
  \|W^j_n\|_{L^3(B_R)}^3
  =\int_{z_n^j+\rho_nB_R}|\phi^j(x)|^3\,dx.
\]
Since \(\rho_n\to0\), the translated domains have
\[
  |z_n^j+\rho_nB_R|=\rho_n^3|B_R|\to0.
\]
By absolute continuity of the integral for \(\phi^j\in L^3(\mathbb R^3)\),
\[
  \|W^j_n\|_{L^3(B_R)}\to0.
\]
To treat a general \(\phi^j\), let \(\phi^j_\varepsilon\in C_c^\infty(\mathbb R^3)\) with
\(\|\phi^j-\phi^j_\varepsilon\|_{L^3}<\varepsilon\).  By critical isometry,
\[
  \|W^j_n-(W^j_n)_\varepsilon\|_{L^3(B_R)}
  \le
  \|\phi^j-\phi^j_\varepsilon\|_{L^3} < \varepsilon,
\]
where \((W^j_n)_\varepsilon\) is the same scaling of \(\phi^j_\varepsilon\).
The compactly supported term has vanishing local norm by the shrinking-domain
argument above, and then letting first \(n\to\infty\), then \(\varepsilon\downarrow0\)
implies the claimed convergence.
A constant vector field from the regular part of an outer flow contributes only
a Galilean drift in the selected frame; it is therefore incorporated into the
translation modulation and does not remain as an exterior perturbation.
\end{proof}

\begin{lemma}[Compactness failure produces another core]\label{paper3:lem:compactness-failure}
If the compact-cylinder suitable estimates are unbounded on a fixed rescaled
cylinder, then after passing to a subsequence the local \(C_n+D_n\) quantity is
unbounded on a comparable cylinder.  Consequently the branch contains an
additional local CKN concentration core.
\end{lemma}

\begin{proof}
Fix \(r>0\).  Apply the local Caccioppoli inequality with outer radius \(2r\).
For each index \(n\), suitable weak solutions satisfy, with \(C_r\) independent
of \(n\),
\[
  A_n(r)+E_n(r)
  \le C_r\bigl(1+C_n(2r)+C_n(2r)^{2/3}+D_n(2r)\bigr),
\]
where the pressure is normalized by subtracting spatial means
\((p_n)_{B_{2r}}(t)\).
\par
Suppose that \(A_n(r)+E_n(r)+C_n(r)+D_n(r)\) is unbounded along a subsequence.
If \(C_n(2r)+D_n(2r)\) were bounded on this subsequence, the preceding estimate
would give uniform boundedness of \(A_n(r)+E_n(r)\).  Monotonicity of the local
integrals also gives \(C_n(r)+D_n(r)\le C_n(2r)+D_n(2r)\), contradicting
unboundedness of the full family on \(Q_r\).  Therefore, after extracting a
subsequence,
\[
  C_n(2r)+D_n(2r)\to\infty .
\]
\par
Define the localized CKN concentration at variable center and radius
\[
  G_n(x,t,\rho):=\rho^{-2}\int_{Q_\rho(x,t)}|u_n|^3\,dx\,ds
  +\rho^{-2}\int_{Q_\rho(x,t)}|p_n-(p_n)_{B_\rho}(s)|^{3/2}\,dx\,ds.
\]
Set
\[
  \mathbf M_n:=\sup\{G_n(x,t,\rho):0<\rho\le 2r,\ Q_\rho(x,t)\subset Q_{4r}\}.
\]
The choice \((x,t,\rho)=(0,0,2r)\) is admissible because \(Q_{2r}(0,0)\subset Q_{4r}\), so
\[
  \mathbf M_n\ge C_n(2r)+D_n(2r)\to\infty.
\]
After extracting a further subsequence, we may assume \(\mathbf M_n\ge1\) for all \(n\).
Choose \((x_n,t_n,\rho_n)\) with \(0<\rho_n\le2r\), \(Q_{\rho_n}(x_n,t_n)\subset Q_{4r}\),
such that
\[
  G_n(x_n,t_n,\rho_n)\ge\tfrac12\,\mathbf M_n\ge\tfrac12.
\]
Rescale around \((x_n,t_n,\rho_n)\):
\[
  v_n(y,s):=\rho_nu_n(x_n+\rho_n y,t_n+\rho_n^2 s),\qquad
  q_n(y,s):=\rho_n^2p_n(x_n+\rho_n y,t_n+\rho_n^2 s).
\]
Using the critical Jacobian identity,
\[
  \int_{Q_1}|v_n|^3\,dy\,ds
  +\int_{Q_1}|q_n-(q_n)_{B_1}(s)|^{3/2}\,dy\,ds
  =G_n(x_n,t_n,\rho_n)\ge\tfrac12.
\]
Hence this new frame carries positive local critical mass, i.e. a new local CKN
concentration core, as required.
\end{proof}

\subsection{Terminal active-frame decoupling}

\begin{definition}[Admissible terminal active family]\label{paper3:def:admissible-terminal}
Fix a compact cylinder \(B_{2R}\times I\) and a selected active frame.  A
terminal active family is admissible on this cylinder if the following hold.
\begin{enumerate}[label=\textup{(\roman*)}]
\item The active profiles satisfy the finite-mass and active-mass conditions of
\cref{paper3:def:active-family}.
\item Every nonselected active profile is either separated from the selected
physical point in the sense of \cref{paper3:lem:separated-vanish}, or has the same
physical point and a strictly outer scale in the sense of
\cref{paper3:lem:outer-vanish}, after comparable same-point scales have been grouped.
\item The rescaled remainder \(r_n\) satisfies
\[
  \|r_n\|_{L^\infty_I L^3(B_{2R})}\to0,
\]
and there exist fields \((\pi_n^r,f_n)\) such that on \(B_{2R}\times I\),
\[
  \partial_t r_n-\Delta r_n+a_n\bigl(r_n+y\cdot\nabla r_n\bigr)+b_n\cdot\nabla r_n
  +\nabla\pi_n^r=f_n,\qquad
  \nabla\cdot r_n=0,
\]
with
\[
  \|f_n\|_{L^1_IH^{-1}(B_R)}\to0,
\]
\[
  \|\nabla\pi_n^r\|_{L^1_IH^{-1}(B_R)}\to0.
\]
\item The selected velocity satisfies
\[
  \sup_n\|U_n\|_{L^\infty_I L^3(B_{2R})}<\infty .
\]
\item The selected modulation coefficients satisfy
\[
  \sup_n\bigl(\|a_n\|_{L^\infty(I)}+\|b_n\|_{L^\infty(I)}\bigr)<\infty .
\]
\item Local pressure decomposition is stable under the splitting: if tensor
coefficients tend to zero in \(L^1_IL^{3/2}(B_R)\), then the corresponding
pressure gradients tend to zero in \(L^1_IH^{-1}(B_R)\), after subtracting
spatial means.
\item Cutoff commutators generated by localizing to \(B_R\) are controlled by
the \(L^\infty_I L^3\), \(L^1_I L^{3/2}\), and pressure-gradient norms appearing
above.  In particular, for any compact \(\chi_R\in C_c^\infty(B_{2R})\) with
\(\chi_R\equiv1\) on \(B_R\), the commutators satisfy
\[
  \|[\Delta,\chi_R]u\|_{H^{-1}(B_R)}
  +\|[\nabla\chi_R,\,\nabla\cdot](u\otimes u)\|_{H^{-1}(B_R)}
  \le C_R\|u\|_{L^3(B_{2R})}^2.
\]
uniformly for \(u\in L^3(B_{2R})\) and the same type of estimate for
\(S_n\)-quadratic commutator terms.
\end{enumerate}
\end{definition}

\begin{theorem}[Terminal profile-evolution decoupling]\label{paper3:thm:terminal-decoupling}
Consider an admissible terminal active family on \(B_{2R}\times I\).  Let
\(U_n\) be the selected-frame velocity and let \(S_n\) be the sum, in that
frame, of all nonselected active profiles together with the rescaled remainder.
Then on \(B_R\times I\):
\begin{enumerate}[label=\textup{(\roman*)}]
\item \(\nabla\cdot S_n=0\) distributionally;
\item \(S_n\to0\) in \(L^\infty_I L^3(B_{2R})\);
\item \(U_n\otimes S_n+S_n\otimes U_n+S_n\otimes S_n\to0\) in
\(L^1_I L^{3/2}(B_R)\);
\item
\[
  \begin{aligned}
  &\exists\,\pi_n,\ F_n:\quad
  \partial_t S_n-\Delta S_n
  +a_n\bigl(y\cdot\nabla S_n + S_n\bigr)
  +b_n\cdot\nabla S_n+\nabla\pi_n =F_n,\\
  &\|F_n\|_{L^1_IH^{-1}(B_R)}\to0,\qquad
  \nabla\cdot S_n=0\quad\text{in }\mathcal D'(B_R\times I).
  \end{aligned}
\]
\[
  \|\nabla\pi_n\|_{L^1_IH^{-1}(B_R)}\to0.
\]
\item every pressure source containing at least one factor of \(S_n\) has
vanishing pressure gradient in \(L^1_IH^{-1}(B_R)\);
\item localization errors introduced by compact cutoffs vanish in
\(L^1_IH^{-1}(B_R)\);
\item the effective modulation parameters of the selected frame remain bounded;
\item after discarding \(S_n\), the selected frame remains a positive finite-mass
local Type II branch with pressure reconstruction and the compact-cylinder
bounds required by the single-core analysis.
\end{enumerate}
\end{theorem}

\begin{proof}
Each profile in the family is divergence-free, and the rescaled remainder is
divergence-free.  Since divergence commutes with the Navier--Stokes critical
rescalings, finite sums remain divergence-free in distributions.  This establishes
\((i)\).

We prove \((ii)\).  By admissibility, every nonselected active profile is either
separated from the selected physical point or is a same-point strictly outer
scale.  The separated profiles vanish locally by \cref{paper3:lem:separated-vanish},
and the same-point outer profiles vanish locally by \cref{paper3:lem:outer-vanish}.
The remainder tends to zero in \(L^\infty_I L^3(B_{2R})\) by
\cref{paper3:def:admissible-terminal}.  The number of active profiles is finite by
\cref{paper3:lem:finite-active-count}.  Since every summand in \(S_n\)
tends to zero in \(L^\infty_I L^3(B_{2R})\), the triangle inequality gives
\[
  \|S_n\|_{L^\infty_I L^3(B_{2R})}
  \le
  \sum_{j\ne j_*}\|S^j_n\|_{L^\infty_I L^3(B_{2R})}
  +\|r_n\|_{L^\infty_I L^3(B_{2R})}\to0.
\]
This establishes \((ii)\).

For \((iii)\), admissibility gives
\[
  \sup_n\|U_n\|_{L^\infty_I L^3(B_{2R})}<\infty.
\]
Hence, by H\"older's inequality,
\[
  \|U_n\otimes S_n\|_{L^1_IL^{3/2}(B_R)}
  \le |I|\|U_n\|_{L^\infty_IL^3(B_R)}
       \|S_n\|_{L^\infty_IL^3(B_R)}\to0,
\]
and the same estimate applies to \(S_n\otimes U_n\).  Finally,
\[
  \|S_n\otimes S_n\|_{L^1_IL^{3/2}(B_R)}
  \le |I|\|S_n\|_{L^\infty_IL^3(B_R)}^2\to0.
\]

For \((iv)\), write the represented Navier--Stokes equation for the full
profile sum \(U_n+S_n\) and subtract the represented equation for the selected
component \(U_n\).  The remaining equation for \(S_n\) has the displayed linear
left-hand side
\[
  \partial_t S_n-\Delta S_n
  +a_n\bigl(y\cdot\nabla S_n+S_n\bigr)+b_n\cdot\nabla S_n.
\]
Collect all right-hand terms into
\[
  F_n:= -\nabla\cdot(U_n\otimes S_n+S_n\otimes U_n+S_n\otimes S_n)
       +f_n+E_n^{\mathrm{cut}}+E_n^{\mathrm{mod}},
\]
where \(f_n\) is the remainder forcing from item \((iii)\) of admissibility, and
\(E_n^{\mathrm{cut}},E_n^{\mathrm{mod}}\) are localization and modulation errors.
The nonlinear stresses tend to zero in \(L^1_IL^{3/2}(B_R)\), hence in
\(L^1_IH^{-1}(B_R)\), because \(L^{3/2}(B_R)\hookrightarrow H^{-1}(B_R)\) on a
bounded three-dimensional ball.  The residual term vanishes by
\cref{paper3:def:admissible-terminal}.  The cutoff terms vanish by the localization
part of admissibility, recorded explicitly in \((vi)\).  This gives the
displayed convergence in \(L^1_IH^{-1}(B_R)\).

For \((v)\), decompose the pressure source of \(U_n+S_n\) into the pure selected
source, the pure exterior source, and the two mixed sources.  The mixed sources
are
\[
  \partial_i\partial_j(U_{n,i}S_{n,j}+S_{n,i}U_{n,j})
  \quad\text{and}\quad
  \partial_i\partial_j(S_{n,i}S_{n,j}).
\]
By \((iii)\), their tensor coefficients tend to zero in
\(L^1_IL^{3/2}(B_R)\).  The pressure-stability part of admissibility,
equivalently local Calderon--Zygmund estimates plus harmonic pressure control
modulo spatial means, gives convergence of the corresponding pressure gradients
to zero in \(L^1_IH^{-1}(B_R)\).  Pure exterior sources are generated by
profiles that are locally invisible in the selected frame; applying the same
pressure stability after \cref{paper3:lem:separated-vanish,paper3:lem:outer-vanish} gives
their vanishing pressure gradients.

For \((vi)\), let \(\chi_R\) be a cutoff equal to one on \(B_R\) and supported
in \(B_{2R}\).  The commutators are finite sums of terms of the following
types:
\[
  (\nabla\chi_R)\cdot\nabla S_n,
  \qquad
  (\Delta\chi_R)S_n,
  \qquad
  (\nabla\chi_R)\cdot (U_n\otimes S_n+S_n\otimes U_n+S_n\otimes S_n),
\]
plus analogous pressure cutoff terms.  Their supports lie in the fixed annulus
\(B_{2R}\setminus B_R\).  By the localization part of admissibility, these
commutators are controlled by the norms already shown to vanish in \((ii)\),
\((iii)\), and \((v)\).  Hence all localization errors vanish in
\(L^1_IH^{-1}(B_R)\).

For \((vii)\), the only exterior contribution that can remain nonzero at first
order in a compact selected frame is a constant ambient velocity.  This has
already been absorbed into the translation modulation.  The remaining exterior
profiles vanish locally, so their contribution to the modulation equations is
small in the same local \(L^3\) topology used above.  The selected modulation
coefficients are bounded by admissibility; hence the effective parameters remain
bounded.

For \((viii)\), the positive CKN mass of the selected profile is unchanged by
subtracting a term that tends to zero in local \(L^3\) and whose
mixed pressure sources vanish by \((v)\).  The pressure reconstruction is stable
under the same decomposition by the pressure-stability condition in
\cref{paper3:def:admissible-terminal}.  The compact-cylinder bounds pass to the
selected component because the full family is bounded and the discarded terms
vanish on compact cylinders.  Thus the selected frame remains a positive
finite-mass local Type II branch with the analytic inputs needed for the
single-core analysis.
\end{proof}

\subsection{Same-point and separated-point reductions}

\begin{theorem}[Same-point cascade reduction]\label{paper3:thm:same-point-reduction}
Every same-point multibubble or strict same-point cascade reduces, in the
selected innermost or compound frame, to one of the following: the single-core
criterion, scale-rigidity, or a perturbative terminal branch covered by
\cref{paper3:thm:terminal-decoupling}, provided the terminal family generated in the
selected innermost frame is admissible in the sense of
\cref{paper3:def:admissible-terminal}.
\end{theorem}

\begin{proof}
\noindent\textit{Step 1: scale-comparability classes.}
Split active same-point profiles into equivalence classes by
\[
  i\sim j \iff \exists c,C\in(0,\infty):
  c\le \frac{\lambda^i_n}{\lambda^j_n}\le C \quad \text{for all }n
\]
after passing to a subsequence.
Thus within each class there are constants \(c_{ij},C_{ij}\) with
\(c_{ij}\le\lambda^i_n/\lambda^j_n\le C_{ij}\) for all \(n\), so one may define
a single common rescaling sequence (choose \(\lambda_n^{(k)}=\lambda_n^{j_k}\)) and write
all class-\(k\) profiles in the same coordinates.
Within each class, finite sums of divergence-free profiles remain divergence-free.

\smallskip
\noindent\textit{Step 2: compound profile reduction inside each class.}
For each class set
\[
  \Phi^k_n:=\sum_{j\in J_k}\lambda^j_n\phi^j(x_n^j+\lambda^j_n y).
\]
If the \(L^3\)-size of this rescaled sum satisfies
\[
  \left\|\sum_{j\in J_k}\phi^j\right\|_{L^3(\mathbb R^3)}<\varepsilon_\star,
\]
with \(\varepsilon_\star\) the perturbative small-data threshold, then standard local
small-data theory for Navier--Stokes in the critical space eliminates this class
from the active alternatives.

If the class is above threshold, it is retained as a single compound active core;
at most one active core from that class is needed for local selection.

\smallskip
\noindent\textit{Step 3: local alternatives.}
If the retained compound core is unique and its scale-centering Jacobian is
nondegenerate, we obtain the single selected core alternative.
If this Jacobian is degenerate, the localized choice of active center or scale is
not unique, which is the gauge-degenerate alternative in
\cref{paper3:def:local-multibubble}.

\smallskip
\noindent\textit{Step 4: strict same-point scale separation.}
Assume now the branch is not reduced to comparable scales and choose the innermost
active scale \(\lambda^1_n\) in that physical point.  By
\cref{paper3:lem:outer-vanish}, every profile with larger scale vanishes in the
\(\lambda^1_n\)-frame in local \(L^3\), up to constant drift terms.
These constant drifts are absorbed into translation modulation by admissibility.
Let \(S_n\) be the sum of all such outer profiles plus the rescaled remainder
\(r_n\).  Then, by Step 2 and \cref{paper3:lem:outer-vanish},
\[
  S_n\to0\quad\text{in }L^\infty_I L^3(B_{2R}).
\]

Subtracting the equation for the innermost selected profile from the full-frame
equation gives a perturbation equation for the selected profile whose forcing
belongs to \(L^1_I H^{-1}(B_R)\) and vanishes.  Thus
\cref{paper3:thm:terminal-decoupling} applies, and one of three outcomes holds:
\begin{enumerate}[label=\textup{(\alph*)}]
\item \textbf{Single-core branch:}
  after discarding \(S_n\), the branch is governed by one active component, so it
  reduces to the compact single-core branch of \cref{paper3:ass:single-core};
\item \textbf{Scale-rigid branch:}
  the relative Jacobian between comparable scales becomes rigid in the sense of
  \cref{paper3:ass:scale-rigidity}, so \cref{paper3:ass:scale-rigidity} excludes it;
\item \textbf{Purely perturbative branch:}
  \(S_n\to0\) and the selected profile satisfies a vanishing forcing perturbation
  in \(L^1_IH^{-1}(B_R)\).  No independent same-point cascade mechanism remains;
  after discarding \(S_n\), the selected component is reduced to the single-core
  or scale-rigidity analysis.
\end{enumerate}

These are the strict same-point scale-separation alternatives, so the
same-point reduction is exhausted.
\end{proof}


\begin{theorem}[Separated-point decoupling]\label{paper3:thm:separated-point-reduction}
Separated active physical points are locally invisible in each other's active
rescaled frames.  Consequently every separated-point multibubble reduces in each
selected active frame to a single-core, scale-rigid, or perturbative terminal
branch, provided the terminal family in that selected active frame is
admissible in the sense of \cref{paper3:def:admissible-terminal}.
\end{theorem}

\begin{proof}
Fix an active point \(p\) with frame \((x^p_n,\lambda^p_n)\).  Let \(q\ne p\)
be another active point.  By separation,
\[
  \frac{|x^q_n-x^p_n|}{\lambda^p_n}\to\infty.
\]
Thus \cref{paper3:lem:separated-vanish} gives local \(L^3\) convergence to zero of the
\(q\)-profile in the \(p\)-frame on every compact ball.  The same conclusion
holds for each profile centered at any physical point different from \(p\).
Since only finitely many active profiles remain after the critical mass
threshold is imposed, the sum of all profiles outside the \(p\)-frame tends to
zero locally in \(L^3\).

The pure pressure generated by a separated profile is locally invisible by the
same change of variables whenever its profile pressure belongs to
\(L^{3/2}_{\mathrm{loc}}\) with the standard Calderon--Zygmund control.  Mixed
pressure terms contain at least one exterior velocity factor, hence vanish in
\(L^1H^{-1}_{\mathrm{loc}}\) by the pressure part of
\cref{paper3:thm:terminal-decoupling}.  Localization and modulation commutators are
localized to \(B_{2R}\setminus B_R\), so their \(H^{-1}\)-norm is controlled by
the \(L^3\)-smallness already proved in \cref{paper3:lem:separated-vanish,paper3:lem:outer-vanish}
and by admissibility item \((\text{vii})\).

Therefore, in the selected \(p\)-frame, all other physical points contribute
only a perturbative forcing tending to zero in \(L^1H^{-1}\) on compact
cylinders.  The selected frame remains a positive finite-mass branch with the
same compact-cylinder and pressure reconstruction properties.  Hence the branch
reduces to the single-core case, the scale-rigid case, or the perturbative
terminal branch.  Since \(p\) was arbitrary, the same reduction holds in every
active frame.
\end{proof}

\subsection{Multibubble and cascade exclusion}

The following assumptions isolate the analytic inputs needed for the
multibubble reduction: the single-core exclusion, the scale-rigidity reduction,
and terminal active-frame admissibility.  They are stated explicitly so that the
following theorem uses no unstated compactness, pressure, or profile completeness
conclusion.

\begin{assumption}[Local single-core exclusion]\label{paper3:ass:single-core}
Every compact single-core branch with positive local CKN concentration,
nondegenerate scale-center selection, pressure reconstruction, bounded
compact-cylinder suitable estimates, bounded modulation, and vanishing localized
dissipation is excluded from the local Type II regime.
\end{assumption}

\begin{assumption}[Local scale-rigidity exclusion]\label{paper3:ass:scale-rigidity}
Every scale-rigid local branch produced by the same-point or separated-point
reductions, including a branch with degenerate relative scale-center Jacobian,
either has a nonzero bounded selected-window limit or reduces to the single-core
branch in \cref{paper3:ass:single-core}.  Consequently no such branch is compatible
with the local Type II condition.
\end{assumption}

\begin{assumption}[Terminal admissibility in active frames]\label{paper3:ass:terminal-admissibility}
Every terminal family produced by the same-point and separated-point reductions
is admissible on each compact selected cylinder in the sense of
\cref{paper3:def:admissible-terminal}.
\end{assumption}

\begin{theorem}[Local multibubble and cascade exclusion]\label{paper3:thm:multi}
Assume the local single-core and scale-rigidity exclusions
\cref{paper3:ass:single-core,paper3:ass:scale-rigidity}, and assume terminal admissibility in
active frames, \cref{paper3:ass:terminal-admissibility}.  Let a positive local Type II
branch admit an active family satisfying the hypotheses of
\cref{paper3:thm:minimal-bad-extraction}.  Then no local multibubble, cascade, or
gauge-degenerate Type II case can occur in this class.
\end{theorem}

\begin{proof}
Let a local Type II branch be in the multibubble/cascade class of
\cref{paper3:def:local-multibubble}.  Apply \cref{paper3:thm:minimal-bad-extraction}.  The
positive local concentration is carried by an active frame or compound profile,
and one of the alternatives listed there occurs.
\par
\textit{Case 1: compactness-failure reduction.}
Compactness failure yields, after changing radius by a fixed factor and
rescaling as in \cref{paper3:lem:compactness-failure}, a bounded-cylinder concentration
branch.  Relabel this branch as an active core.  After passing to a subsequence,
the new core and every previously retained active core have one of three
relations: comparable scale at the same physical point, strictly separated scale
at the same physical point, or separated physical center.  Comparable same-point
cores are grouped into compound profiles.  A perturbative compound profile is
removed by small-data stability; an active compound profile is treated as one
selected core in its active frame.  If the scale-center selection inside such a
compound core is degenerate, the branch is gauge-degenerate and belongs to the
non-uniqueness alternative in \cref{paper3:def:local-multibubble}, rather than to
a separate core dynamics.
\par
\textit{Case 1a: gauge degeneracy.}
If the scale-center selection of the selected compound or terminal frame is
degenerate, then the branch is in the gauge-degenerate alternative of
\cref{paper3:def:local-multibubble}.  By the repaired-gauge representation
from \cref{sec:repaired-gauge-construction}, nondegeneracy is the hypothesis needed to enter the
single-core branch.  Persistent degeneracy therefore belongs to the
scale-rigidity alternative in \cref{paper3:ass:scale-rigidity}; if the degeneracy
disappears after passing to a subsequence, the branch enters the single-core or
terminal-frame alternatives below.
\par
\textit{Case 1b: same-point comparable scales.}
Same-point comparable-scale active frames are grouped into the compound profile
of \cref{paper3:def:compound-profile}.  If the compound profile is below the critical
small-data threshold, it is part of the perturbative remainder.  If it is above
threshold and its scale-center selection is nondegenerate, it is one selected
active core and enters the single-core branch.  If the corresponding
scale-center selection is degenerate, it is covered by Case 1a.  Thus
comparable same-point profiles do not produce an additional multibubble
mechanism.
\par
\textit{Case 2: same-point cascading.}
Strict same-point scale cascades are treated by
\cref{paper3:thm:same-point-reduction}.  In the innermost frame, all outer scales are
locally invisible after constant drifts are absorbed.  The branch therefore
reduces to one of the three terminal alternatives
\[
  \text{(A) single-core},\quad \text{(B) scale-rigid},\quad \text{(C) terminal perturbative}.
\]
\par
\textit{Case 3: separated-point interactions.}
For separated physical points, \cref{paper3:thm:separated-point-reduction} gives the same
three terminal alternatives as in Case 2:
\[
  \text{(A) single-core},\quad \text{(B) scale-rigid},\quad \text{(C) terminal perturbative}.
\]

The single-core alternative is excluded by \cref{paper3:ass:single-core}.  The
scale-rigid alternative is excluded by \cref{paper3:ass:scale-rigidity}, which includes
the scale-collapse and autonomous-limit exclusions in \cref{sec:scale-collapse-cost} in this setting.  The
terminal perturbative alternative gives no independent multibubble obstruction:
\cref{paper3:thm:terminal-decoupling} removes all exterior profiles and remainders from
the selected active equation in \(L^1H^{-1}_{\mathrm{loc}}\), preserves the positive
finite-mass selected branch, and leaves only the single-core or scale-rigid
possibilities already excluded.  Consequently every alternative in
\cref{paper3:def:local-multibubble} is incompatible with the local Type II regime.
\end{proof}

\begin{corollary}[Removal of the multibubble alternative]\label{paper3:cor:multi-removal}
Under the hypotheses of \cref{paper3:thm:multi}, every failure of the single
nondegenerate compact core alternative is excluded.
\end{corollary}

\begin{proof}
By \cref{paper3:def:local-multibubble}, failure of the single nondegenerate compact
core alternative means that at least one of the following occurs: the
compact-cylinder suitable estimates are unbounded; positive local CKN mass splits
between active cores; a strict same-point scale cascade occurs; separated
physical points remain active in distinct local frames; or the localized scale
or center selection is degenerate.  The compactness-failure case produces an
additional core by \cref{paper3:lem:compactness-failure}.  Same-point and separated
cases are reduced by \cref{paper3:thm:same-point-reduction,paper3:thm:separated-point-reduction}.
\Cref{paper3:thm:multi} excludes the remaining cases.  Hence no failure of the
single nondegenerate compact core alternative remains.
\end{proof}

\clearpage
\section{Local Caccioppoli and Rough-Core Exclusion}\label{sec:rough-core-exclusion}
\subsection{Analytic inputs and notation}

The renormalized local energy inequality from the preceding renormalized-chart
analysis is used as an analytic input.  The renormalized suitable structure and
pressure reconstruction are not rederived.  All estimates are conditional on
the following local hypotheses on the compact renormalized cylinders under
consideration.

\begin{assumption}[Renormalized local energy inequality]\label{paper4:ass:lei}
Let $(V,P,a,b)$ be defined on a compact renormalized cylinder
$Q_{2R,J}:=B_{2R}\times J$, where $J\subset\mathbb R$ is a bounded interval.
Assume
\begin{equation}\label{paper4:eq:renorm-ns}
\partial_\tau V+(V\cdot\nabla)V+\nabla P-\nu\Delta V
= a(\tau)V+a(\tau)(y\cdot\nabla V)+b(\tau)\cdot\nabla V,
\qquad
\nabla\cdot V=0,
\end{equation}
with $\nu>0$, $\nabla\cdot V=0$, and $a,b\in L^\infty(J)$. For every
nonnegative $\phi\in C_c^\infty(Q_{2R,J})$ and every $\tau_1<\tau_2$ in $J$,
\begin{equation}\label{paper4:eq:local-energy}
\begin{aligned}
&\int_{B_{2R}}\frac{|V(\cdot,\tau_2)|^2}{2}\phi(\cdot,\tau_2)
-\int_{B_{2R}}\frac{|V(\cdot,\tau_1)|^2}{2}\phi(\cdot,\tau_1)
\\
&\quad+\nu\int_{\tau_1}^{\tau_2}\!\int_{B_{2R}}\phi\,|\nabla V|^2
\le
\int_{\tau_1}^{\tau_2}\!\int_{B_{2R}}\frac{|V|^2}{2}(\partial_\tau\phi+\nu\Delta\phi)
\\
&\quad+\int_{\tau_1}^{\tau_2}\!\int_{B_{2R}}\frac{|V|^2}{2}V\cdot\nabla\phi
+\int_{\tau_1}^{\tau_2}\!\int_{B_{2R}}PV\cdot\nabla\phi
\\
&\quad+\int_{\tau_1}^{\tau_2}\!\int_{B_{2R}}
a\frac{|V|^2}{2}(-\phi-y\cdot\nabla\phi)
\\
&\quad-\int_{\tau_1}^{\tau_2}\!\int_{B_{2R}}\frac{|V|^2}{2}b\cdot\nabla\phi.
\end{aligned}
\end{equation}
The pressure is understood modulo functions of time, and the mean-normalized
pressure on compact balls belongs to the local class used below.
\end{assumption}

Throughout, for a space--time set $E\subset\mathbb R^3\times\mathbb R$, write
\[
\int_E f:=\int_E f(y,\tau)\,dy\,d\tau.
\]
For $r>0$ and a bounded interval $J\subset\mathbb R$, define
\[
Q_{r,J}:=B_r\times J,\qquad |Q_{r,J}|:=|B_r|\,|J|.
\]
For a ball $B\subset\mathbb R^3$, denote the spatial mean of pressure by
\[
(P)_B(\tau):=\frac1{|B|}\int_B P(y,\tau)\,dy.
\]
The local quantities are
\[
\mathcal A_J(r):=\operatorname*{ess\,sup}_{\tau\in J}\int_{B_r}|V(y,\tau)|^2\,dy,
\]
\[
\mathcal H_J(r):=\int_J\!\int_{B_r}|\nabla V|^2\,dy\,d\tau,
\]
\[
\mathcal C_J(r):=\int_J\!\int_{B_r}|V|^3\,dy\,d\tau,
\qquad
\mathcal D_J(r):=\int_J\!\int_{B_r}|P-(P)_{B_r}(\tau)|^{3/2}\,dy\,d\tau.
\]

\begin{lemma}[Finite-measure conversion]\label{paper4:lem:l2-from-l3}
For every bounded interval $J$ and every $r>0$,
\[
\int_J\!\int_{B_r}|V|^2
\le |Q_{r,J}|^{1/3}\mathcal C_J(r)^{2/3}
\le |Q_{r,J}|^{1/3}\bigl(1+\mathcal C_J(r)\bigr).
\]
\end{lemma}

\begin{proof}
H\"older's inequality on $Q_{r,J}$ gives
\[
\int_J\!\int_{B_r}|V|^2
\le |Q_{r,J}|^{1/3}\left(\int_J\!\int_{B_r}|V|^3\right)^{2/3}.
\]
The second inequality follows from $s^{2/3}\le 1+s$ for $s\ge0$.
\end{proof}

\begin{lemma}[Comparison of pressure normalizations]\label{paper4:lem:pressure-means}
Let $1\le p<\infty$ and $0<r\le R$. For almost every $\tau$ such that
$P(\tau)\in L^p(B_R)$,
\[
\|P(\tau)-(P)_{B_r}(\tau)\|_{L^p(B_r)}
\le 2\|P(\tau)-(P)_{B_R}(\tau)\|_{L^p(B_R)}.
\]
Consequently,
\[
\mathcal D_J(r)\le 2^{3/2}\mathcal D_J(R)
\]
whenever the right-hand side is finite.
\end{lemma}

\begin{proof}
Set $c=(P)_{B_R}(\tau)$. Then
\[
\|P-(P)_{B_r}\|_{L^p(B_r)}
\le \|P-c\|_{L^p(B_r)}+|B_r|^{1/p}|(P)_{B_r}-c|.
\]
Since
\[
|(P)_{B_r}-c|\le |B_r|^{-1}\int_{B_r}|P-c|
\le |B_r|^{-1/p}\|P-c\|_{L^p(B_r)},
\]
the displayed pointwise estimate follows. Raising to the power $3/2$ and
integrating in time gives the last assertion.
\end{proof}

\subsection{Compact-cylinder Caccioppoli estimate}

\begin{theorem}[Compact-cylinder Caccioppoli estimate]\label{paper4:thm:caccioppoli}
Let $R>0$, let $J\subset\mathbb R$ be a bounded open interval, and let
$J'\Subset J$ be a compact subinterval. Assume
\cref{paper4:ass:lei} on $Q_{2R,J}$ and assume
\[
\mathcal C_J(2R)+\mathcal D_J(2R)<\infty.
\]
Then there is a constant
\[
C=C\bigl(R,|J|,\nu,\|a\|_{L^\infty(J)},\|b\|_{L^\infty(J)},
\operatorname{dist}(J',\partial J)\bigr)
\]
such that
\[
\mathcal A_{J'}(R)+\mathcal H_{J'}(R)
\le C\Bigl(1+
\mathcal C_J(2R)+\mathcal D_J(2R)
\Bigr).
\]
The estimate is outer-to-inner: all right-hand quantities are measured on
$B_{2R}\times J$, while the conclusion is on $B_R\times J'$.
\end{theorem}

\begin{proof}
\textbf{Step 1: localized cut-off and the two energy inequalities.}
Set
\[
\delta:=\operatorname{dist}(J',\partial J)>0.
\]
Choose $\eta\in C_c^\infty(J)$ and $\zeta\in C_c^\infty(B_{2R})$ with
\[
0\le\eta\le1,\qquad
0\le\zeta\le1,\qquad
\eta\equiv1\ \text{on}\ J',\qquad
\zeta\equiv1\ \text{on}\ B_R,
\]
and
\[
\|\eta'\|_{L^\infty(J)}\le C_\eta\,\delta^{-1},
\qquad
\|\nabla\zeta\|_{L^\infty(B_{2R})}\le C_\zeta R^{-1},
\qquad
\|\Delta\zeta\|_{L^\infty(B_{2R})}\le C_\zeta R^{-2}.
\]
Set
\[
\phi(y,\tau):=\eta(\tau)\,\zeta(y)^2.
\]
Then $\phi\in C_c^\infty(Q_{2R,J})$, $0\le\phi\le1$, and
\[
\phi\equiv1\text{ on }B_R\times J'.
\]
Choose $\tau_-<\inf J'$ and $\tau_+>\sup J'$ in $J$ with
$\eta(\tau_-)=\eta(\tau_+)=0$. Applying \eqref{paper4:eq:local-energy} on
$(\tau_-,\tau_+)$ gives
\[
\nu\int_{\tau_-}^{\tau_+}\!\int_{B_{2R}}\phi\,|\nabla V|^2
\le |R_\tau|+|R_\Delta|+|R_{\mathrm{conv}}|+|R_{\mathrm{pres}}|+|R_{\mathrm{mod}}|,
\]
where
\[
R_\tau:=\int_J\!\int_{B_{2R}}\frac{|V|^2}{2}\,\partial_\tau\phi,\qquad
R_\Delta:=\nu\int_J\!\int_{B_{2R}}\frac{|V|^2}{2}\,\Delta\phi,
\]
\[
R_{\mathrm{conv}}:=\int_J\!\int_{B_{2R}}\bigl(\tfrac{|V|^2}{2}V\bigr)\cdot\nabla\phi,\qquad
R_{\mathrm{pres}}:=\int_J\!\int_{B_{2R}}\bigl(P-(P)_{B_{2R}}(\tau)\bigr)V\cdot\nabla\phi,
\]
\[
R_{\mathrm{mod}}:=\int_J\!\int_{B_{2R}}
a\frac{|V|^2}{2}(-\phi-y\cdot\nabla\phi)
-\int_J\!\int_{B_{2R}}\frac{|V|^2}{2}b\cdot\nabla\phi.
\]
For almost every $\sigma\in J'$, applying \eqref{paper4:eq:local-energy} on
$(\tau_-,\sigma)$ gives the same right-hand estimates and leaves the terminal
term $\frac12\int_{B_R}|V(\sigma)|^2$ on the left. Therefore the bounds below
control both $\mathcal H_{J'}(R)$ and $\mathcal A_{J'}(R)$.
Define
\[
L:=\int_J\!\int_{B_{2R}}|V|^2,\qquad
C_3:=\mathcal C_J(2R),\qquad D_{3/2}:=\mathcal D_J(2R).
\]
By \cref{paper4:lem:l2-from-l3},
\[
L\le |Q_{2R,J}|^{1/3}C_3^{2/3}\le C(R,|J|)(1+C_3).
\]

Since $\phi(y,\tau)=\eta(\tau)\zeta(y)^2$,
\[
\partial_\tau\phi=\eta'(\tau)\zeta(y)^2,
\qquad
\nabla\phi=2\eta(\tau)\zeta(y)\nabla\zeta(y),
\qquad
\Delta\phi=\eta(\tau)\Delta(\zeta(y)^2),
\]
hence
\[
\|\partial_\tau\phi\|_{L^\infty(Q_{2R,J})}\le C_\eta\delta^{-1},
\qquad
\|\nabla\phi\|_{L^\infty(Q_{2R,J})}\le 2C_\zeta R^{-1},
\qquad
\|\Delta\phi\|_{L^\infty(Q_{2R,J})}\le C_\zeta R^{-2}.
\]

\textbf{Step 2: derivative terms from the cut-off.}
\[
|R_\tau|
\le\frac12\|\partial_\tau\phi\|_{L^\infty(Q_{2R,J})}L
\le C\delta^{-1}(1+C_3),
\]
\[
|R_\Delta|
\le\frac{\nu}{2}\|\Delta\phi\|_{L^\infty(Q_{2R,J})}L
\le C\nu R^{-2}(1+C_3).
\]

\textbf{Step 3: convection term.}
Using $|\nabla\phi|\le CR^{-1}$,
\[
|R_{\mathrm{conv}}|
\le \frac12\|\nabla\phi\|_{L^\infty(Q_{2R,J})}\int_J\!\int_{B_{2R}}|V|^3
\le CR^{-1}C_3.
\]

\textbf{Step 4: pressure term.}
For each time $\tau$,
\[
\int_{B_{2R}}(P-(P)_{B_{2R}}(\tau))V\cdot \nabla\phi
=\int_{B_{2R}}PV\cdot \nabla\phi,
\]
with
\[
\int_{B_{2R}}(P)_{B_{2R}}(\tau)V\cdot \nabla\phi
=(P)_{B_{2R}}(\tau)\int_{B_{2R}}V\cdot \nabla\phi
=-(P)_{B_{2R}}(\tau)\int_{B_{2R}}\phi\,\nabla\cdot V=0
\]
(using $\phi\in C_c^\infty(B_{2R})$ and $\nabla\cdot V=0$).
Thus subtracting any time-dependent pressure constant is legitimate in the
pressure flux term.
Hence
\[
|R_{\mathrm{pres}}|
\le\|\nabla\phi\|_{L^\infty(Q_{2R,J})}
\int_J
\|P-(P)_{B_{2R}}(\tau)\|_{L_x^{3/2}(B_{2R})}
\|V(\tau,\cdot)\|_{L_x^3(B_{2R})}
\,d\tau
\le CR^{-1}D_{3/2}^{2/3}C_3^{1/3}.
\]
Young's inequality gives
\[
|R_{\mathrm{pres}}|\le C(R)(C_3+D_{3/2}).
\]

\textbf{Step 5: modulation terms.}
The local energy inequality already contains the modulation terms
after the spatial integrations by parts from the repaired-gauge equation:
\[
\int_{B_{2R}}(y\cdot\nabla V)\cdot V\,\phi
=\frac12\int_{B_{2R}}y\cdot\nabla(|V|^2)\phi
=-\frac12\int_{B_{2R}}|V|^2\nabla\cdot(y\phi),
\]
and
\[
\int_{B_{2R}}(b(\tau)\cdot\nabla V)\cdot V\,\phi
=\frac12\int_{B_{2R}}b(\tau)\cdot\nabla(|V|^2)\phi
=-\frac12\int_{B_{2R}}|V|^2b(\tau)\cdot\nabla\phi.
\]
Thus we estimate the additional terms directly. Split
\[
R_{\mathrm{mod}}=R_a+R_{ay}+R_b,
\]
where
\[
R_a:=-\int_J\!\int_{B_{2R}}a\,\frac{|V|^2}{2}\phi,
\quad
R_{ay}:=-\int_J\!\int_{B_{2R}}a\,\frac{|V|^2}{2}\,y\cdot\nabla\phi,
\quad
R_b:=-\int_J\!\int_{B_{2R}}\frac{|V|^2}{2}b\cdot\nabla\phi.
\]
\[
M_a:=\|a\|_{L^\infty(J)},\qquad M_b:=\|b\|_{L^\infty(J)}.
\]
For $R_a$,
\[
|R_a|
\le \frac12 M_aL
\le C M_a(1+C_3).
\]
For the scale-gradient part, $|y|\le2R$ on $B_{2R}$ and
$\|\nabla\phi\|_{L^\infty(Q_{2R,J})}\le CR^{-1}$, hence
\[
|y\cdot\nabla\phi|\le C.
\]
Therefore
\[
|R_{ay}|\le C M_aL\le C M_a(1+C_3).
\]
For $R_b$,
\[
|R_b|\le C M_bR^{-1}L\le C M_bR^{-1}(1+C_3).
\]

\textbf{Step 6: absorption and conclusion.}
Collecting estimates from Steps~2--5 yields
\[
\nu\int_{\tau_-}^{\tau_+}\!\int_{B_{2R}}\phi|\nabla V|^2
\le C\bigl(1+C_3+D_{3/2}\bigr),
\]
where $C$ depends only on
\[
R,|J|,\nu,\|a\|_{L^\infty(J)},\|b\|_{L^\infty(J)},\delta.
\]
Since $\phi\equiv1$ on $B_R\times J'$,
\[
\mathcal H_{J'}(R)
\le C\bigl(1+\mathcal C_J(2R)+\mathcal D_J(2R)\bigr).
\]
Repeating the same estimates with terminal time $\sigma\in J'$ gives
\[
\frac12\int_{B_R}|V(y,\sigma)|^2\,dy
\le C\bigl(1+\mathcal C_J(2R)+\mathcal D_J(2R)\bigr)
\]
for almost every $\sigma\in J'$. Taking the essential supremum over
$\sigma\in J'$ and renaming $C$ proves the theorem.
\end{proof}

\begin{definition}[Local compact-cylinder control]
Local Caffarelli--Kohn--Nirenberg control on $Q_{r,J}$ means
\[
\mathcal C_J(r)+\mathcal D_J(r)<\infty.
\]
Local windowed $H^1$ control on $Q_{r,J}$ means
\[
\mathcal H_J(r)<\infty.
\]
\end{definition}

\begin{corollary}[Outer control implies inner suitable control]\label{paper4:cor:outer-inner}
Let $R>0$, let $J\subset\mathbb R$ be a bounded open interval, and let
$J'\Subset J$ be a compact subinterval. Assume \cref{paper4:ass:lei} on $Q_{2R,J}$ and
$\mathcal C_J(2R)+\mathcal D_J(2R)<\infty$. Then
\[
\mathcal A_{J'}(R)+\mathcal H_{J'}(R)<\infty.
\]
More quantitatively, the bound is the one in \cref{paper4:thm:caccioppoli}. If an
argument uses pressure means on a smaller ball, \cref{paper4:lem:pressure-means}
converts those normalizations to the outer-ball normalization above.
\end{corollary}

\begin{proof}
The estimate follows from \cref{paper4:thm:caccioppoli}. The comparison of pressure normalizations follows from
\cref{paper4:lem:pressure-means}.
\end{proof}

\begin{corollary}[Rough-core exclusion on compact windows]\label{paper4:cor:rough-core}
Let $R>0$, let $J\subset\mathbb R$ be a bounded open interval, and let
$J'\Subset J$ be a compact subinterval. Assume \cref{paper4:ass:lei} on $Q_{2R,J}$. If
\[
\mathcal C_J(2R)+\mathcal D_J(2R)<\infty,
\]
then
\[
\mathcal H_{J'}(R)<\infty.
\]
Equivalently, if $\mathcal H_{J'}(R)=+\infty$ for an inner cylinder
$Q_{R,J'}$, then for every enclosing cylinder $Q_{2R,J}$ with
$J'\Subset J$ on which the local energy inequality and bounded
modulation hypotheses hold,
\[
\mathcal C_J(2R)=+\infty
\qquad\text{or}\qquad
\mathcal D_J(2R)=+\infty.
\]
\end{corollary}

\begin{proof}
The direct implication is the gradient part of \cref{paper4:thm:caccioppoli}. The
last assertion is its contrapositive with the same outer-to-inner geometry:
finite outer $\mathcal C+\mathcal D$ on $Q_{2R,J}$ forces finite inner
$\mathcal H$ on $Q_{R,J'}$.
\end{proof}

\begin{remark}[Meaning of rough core]
Here a rough core means failure of local windowed $H^1$
control on compact inner renormalized cylinders while the modulation
coefficients are bounded. \Cref{paper4:cor:rough-core} shows that such a failure is
not independent of compact-cylinder Caffarelli--Kohn--Nirenberg control: it can
occur only if the velocity-pressure bounds $\mathcal C+\mathcal D$ fail on
every relevant larger enclosing compact cylinder.
\end{remark}

\subsection{Consequences for the assembly}

We shall use the following consequences in the final assembly.
\begin{enumerate}[label=\textup{(\arabic*)}]
\item On every compact renormalized window satisfying the local energy
inequality and bounded modulation, finite $\mathcal C+\mathcal D$ on
$Q_{2R,J}$ implies finite $\mathcal A+\mathcal H$ on $Q_{R,J'}$ whenever
$J'\Subset J$.
\item The pressure normalization may be taken on any comparable ball, after
using \cref{paper4:lem:pressure-means}.
\item Failure of local windowed $H^1$ control on an inner compact cylinder
forces failure of the compact-cylinder $\mathcal C+\mathcal D$ bounds on each
larger enclosing cylinder to which \cref{paper4:thm:caccioppoli} applies.
\end{enumerate}

\clearpage
\section{Scale-Collapse Cost and Autonomous Limits}\label{sec:scale-collapse-cost}
\medskip
\noindent\textit{Notation and functional conventions.}
For a bounded spatial region \(K\subset\mathbb R^3\), we use
\[
  (f,g)_{L^2(K)}:=\int_K f\cdot g\,dx,\qquad
  \langle F,\phi\rangle_{H^{-1}(K),H^1_0(K)}
  :=\int_0^T\langle F(\cdot,s),\phi(\cdot,s)\rangle_{H^{-1}(K),H^1_0(K)}\,ds
\]
whenever \(F\in L^1(0,T;H^{-1}(K))\) and \(\phi\in L^\infty(0,T;H^1_0(K))\). Weak limits and weak convergences are taken in the distributional sense on \(C_c^\infty\)-test functions.
\[
  \|f\|_{L^p_tL^q_x}:=\left(\int_0^T\|f(\cdot,s)\|_{L^q(\mathbb R^3)}^p\,ds\right)^{1/p},
\]
for any fixed \(T>0\), with the usual modification \(p=\infty\).

\subsection{Renormalized scale-collapse setting}

\medskip
\noindent\textit{Position of the result.}
The scale-collapse analysis has two components:
\begin{enumerate}[label=(\arabic*)]
  \item Sections~\ref{paper5:sec:cost}--\ref{paper5:sec:finite-cost} contain the cost
    exclusion arguments, which do not use compactness.
  \item Section~\ref{paper5:sec:autonomous} contains the autonomous compactness
    statement under explicitly stated compactness hypotheses.
\end{enumerate}

Throughout the scale-collapse analysis, the renormalized representation established earlier is assumed: on the
time interval under consideration, the fields \((V,P)\) satisfy \eqref{paper5:eq:renorm-ns} and its
distributional weak formulation.  In particular, the solenoidal formulation is obtained by
restricting the test fields below to divergence-free fields.

Let \((u,p)\) be a divergence-free suitable weak solution of
\[
  \partial_tu+(u\cdot\nabla)u+\nabla p=\nu\Delta u,
  \qquad \nabla\cdot u=0,
\]
in \(\mathbb R^3\), with viscosity \(\nu>0\).  Let
\(\lambda(\tau)>0\), \(X(\tau)\in\mathbb R^3\), and a measurable gauge function
\(c(t)\) be given, and define the renormalized fields by
\[
  u(x,t)=\lambda(\tau)^{-1}V\left(\frac{x-X(\tau)}{\lambda(\tau)},\tau\right),
  \qquad
  p(x,t)=\lambda(\tau)^{-2}P\left(\frac{x-X(\tau)}{\lambda(\tau)},\tau\right)+c(t).
\]
Pressure is defined only up to an arbitrary function of time; \(c(t)\) has no effect on
\(\nabla p\), so it is irrelevant for all subsequent estimates.
Assume that \((V,P)\) satisfies on \((0,\infty)\times\mathbb R^3\)
\begin{equation}\label{paper5:eq:renorm-ns}
  \partial_\tau V+(V\cdot\nabla)V+\nabla P
  =\nu\Delta V+a(\tau)(V+y\cdot\nabla V)+b(\tau)\cdot\nabla V,
  \qquad \nabla\cdot V=0.
\end{equation}
For every \(\psi\in C_c^\infty((0,\infty)\times\mathbb R^3;\mathbb R^3)\) with \(\nabla\cdot\psi=0\),
the pair \((V,P)\) satisfies the weak formulation
\[
\begin{aligned}
0={}&\int_0^\infty\!\!\int_{\mathbb R^3}
\bigl(
-V\cdot\partial_\tau\psi
-(V\otimes V):\nabla\psi
+\nu\nabla V:\nabla\psi
\bigr)\,dy\,d\tau \\
&+\int_0^\infty\!\!\int_{\mathbb R^3}
\bigl(
-a(\tau)(V+y\cdot\nabla V)\cdot\psi
-(b(\tau)\cdot\nabla V)\cdot\psi
+P\,\nabla\cdot\psi
\bigr)\,dy\,d\tau .
\end{aligned}
\]
For every \(0<T<\infty\), restricting the \(\tau\)-integration to \((0,T)\) gives the same identity on \((0,T)\times\mathbb R^3\).
Equivalently, for every
\(\psi\in C_c^\infty((0,\infty)\times\mathbb R^3;\mathbb R^3)\), not necessarily
divergence free,
\[
\begin{aligned}
0={}&\int_0^\infty\!\!\int_{\mathbb R^3}
\bigl(
-V\cdot\partial_\tau\psi
-(V\otimes V):\nabla\psi
\nu\nabla V:\nabla\psi
-P\,\nabla\cdot\psi
\bigr)\,dy\,d\tau \\
&+\int_0^\infty\!\!\int_{\mathbb R^3}
\bigl(
-a(\tau)(V+y\cdot\nabla V)\cdot\psi
-(b(\tau)\cdot\nabla V)\cdot\psi
\bigr)\,dy\,d\tau .
\end{aligned}
\]

Define positive and negative parts of the scalar field \(a\) by
\[
  a_+(\tau):=\max\{a(\tau),0\},
  \qquad
  a_-(\tau):=\max\{-a(\tau),0\}.
\]
Then for a.e. \(\tau\), \(a=a_+-a_-\), \(a_+,a_-\ge0\), and \(a_+a_-=0\).  Assume
\[
  a\in L^1_{\mathrm{loc}}([\tau_0,\infty)),
  \qquad
  b\in L^1_{\mathrm{loc}}([\tau_0,\infty);\mathbb R^3),
\]
and the scale convention
\begin{equation}\label{paper5:eq:log-scale-convention}
  \frac{d}{d\tau}\log\lambda(\tau)=a(\tau)
\end{equation}
for a.e. \(\tau\ge\tau_0\), with \(\log\lambda\) absolutely continuous on compact intervals.

All integrals below are interpreted in \([0,\infty]\). When we write
\(\int_{\tau_0}^{\infty}\), partial integrals are taken over \([\tau_0,T]\) and
monotone convergence is used as \(T\to\infty\).

Fix \(R>0\).  Choose
\(
  \Phi\in C_c^\infty(\mathbb R^3;[0,1])
\)
satisfying
\begin{equation}\label{paper5:eq:Phi-properties}
  \Phi(y)=1\ \text{if }|y|\le R,
  \quad
  \operatorname{supp}_y\Phi\subset B_{2R}.
\end{equation}
Define the localized kinetic mass
\[
  M_R(\tau):=\int_{\mathbb R^3}|V(y,\tau)|^2\Phi(y)
  \,dy.
\]
Since \(\Phi\) is fixed in \(\tau\), \(M_R\) is measurable in \(\tau\), and only this
measurability is used below.

\begin{definition}[Nonvanishing localized core]
A renormalized branch has nonvanishing localized core if there exist
constants \(m_0>0\) and \(\tau_1\ge\tau_0\) such that
\[
  M_R(\tau)\ge m_0\quad\text{for a.e. }\tau\in[\tau_1,\infty).
\]
In applications, this lower bound is supplied by the selected-core construction; here it is stated as an explicit hypothesis.
\end{definition}

\begin{definition}[Scale-collapse and absolute scale-drift densities]
Define
\[
  \mathcal C_{\mathrm{sc}}(\tau):=a_-(\tau)M_R(\tau),
  \qquad
  \mathcal C_{\mathrm{abs}}(\tau):=\nu\int_{\mathbb R^3}|\nabla V(y,\tau)|^2\Phi(y)
  \,dy+|a(\tau)|M_R(\tau).
\]
\(\mathcal C_{\mathrm{sc}}\) is the negative drift weighted by localized core mass, while
\(\mathcal C_{\mathrm{abs}}\) controls both collapse drift and local dissipation.
Their total costs are extended real-valued by
\[
  \int_{\tau_0}^{\infty}\mathcal C_{\mathrm{sc}}(\tau)
  \,d\tau,
  \qquad
  \int_{\tau_0}^{\infty}\mathcal C_{\mathrm{abs}}(\tau)
  \,d\tau.
\]
\end{definition}

\begin{definition}[Finite-cost branch]\label{paper5:def:finite-cost-branch}
A branch has finite scale-collapsing cost if
\[
  \int_{\tau_0}^{\infty}\mathcal C_{\mathrm{sc}}(\tau)
  \,d\tau<\infty.
\]
It has finite absolute cost if
\[
  \int_{\tau_0}^{\infty}\mathcal C_{\mathrm{abs}}(\tau)
  \,d\tau<\infty.
\]
\end{definition}

\subsection{Elementary alternative}
\label{paper5:sec:cost}

\begin{lemma}[Finite-or-infinite alternative]\label{paper5:lem:finite-infinite}
Let \(f\ge0\) be measurable on \([\tau_0,\infty)\). Then exactly one of the following holds:
\[
  \int_{\tau_0}^{\infty}f(\tau)
  \,d\tau<\infty
  \quad\text{or}\quad
  \int_{\tau_0}^{\infty}f(\tau)
  \,d\tau=\infty.
\]
\end{lemma}

\begin{proof}
\begin{enumerate}[label=(\arabic*)]
  \item For \(T>\tau_0\), define partial costs
  \[
    F(T):=\int_{\tau_0}^{T}f(\tau)
    \,d\tau\in[0,\infty].
  \]

  \item If \(T_2>T_1\), then by nonnegativity of \(f\),
  \[
    F(T_2)-F(T_1)=\int_{T_1}^{T_2}f(\tau)
    \,d\tau\ge0,
  \]
  so \(F\) is monotone nondecreasing.

  \item A monotone function on \([\tau_0,\infty)\) has an extended limit
  \(L\in[0,\infty]\),
  \[
    L:=\lim_{T\to\infty}F(T).
  \]
  By definition of improper integral, this limit equals
  \(\int_{\tau_0}^{\infty}f(\tau)\,d\tau\).

  \item Therefore \(L\) is either finite or equal to \(+\infty\), and these
  cases are mutually exclusive.
\end{enumerate}
Hence exactly one item in the lemma holds.
\end{proof}

\begin{lemma}[Logarithmic scale identity]\label{paper5:lem:log-scale}
Under \eqref{paper5:eq:log-scale-convention}, for all \(\tau_1,\tau_2\) with
\(\tau_0\le\tau_1<\tau_2<\infty\),
\[
  \int_{\tau_1}^{\tau_2}a(\tau)
  \,d\tau
  =\log\lambda(\tau_2)-\log\lambda(\tau_1).
\]
\end{lemma}

\begin{proof}
\begin{enumerate}[label=(\arabic*)]
  \item Since \(a\in L^1_{\mathrm{loc}}\), the map
  \(\log\lambda\) is absolutely continuous on \([\tau_1,\tau_2]\), hence has a
  representative with a.e. derivative equal to \(a\).

  \item By the fundamental theorem of calculus for absolutely continuous functions,
  for every \(\tau\in[\tau_1,\tau_2]\),
  \[
    \log\lambda(\tau)=\log\lambda(\tau_1)+\int_{\tau_1}^{\tau}a(s)
    \,ds.
  \]

  \item Setting \(\tau=\tau_2\) yields the identity.
\end{enumerate}
\end{proof}

\subsection{Scale collapse versus localized core lower bounds}

\begin{theorem}[Collapse with nonvanishing localized core forces infinite
scale-collapsing cost]\label{paper5:thm:collapse-core-floor}
Assume:
\begin{enumerate}[label=(\arabic*)]
  \item \(\lambda(\tau)\to0\) as \(\tau\to\infty\);
  \item the branch has nonvanishing localized core.
\end{enumerate}
Then
\[
  \int_{\tau_0}^{\infty}\mathcal C_{\mathrm{sc}}(\tau)
  \,d\tau
  =\int_{\tau_0}^{\infty}a_-(\tau)M_R(\tau)
  \,d\tau=\infty.
\]
\end{theorem}

\begin{proof}
\begin{enumerate}[label=(\arabic*)]
  \item Choose \(m_0>0\) and \(\tau_1\ge\tau_0\) from the nonvanishing
  core assumption so that
  \[
    M_R(\tau)\ge m_0\quad\text{for a.e. }\tau\in[\tau_1,\infty).
  \]

  \item By Lemma~\ref{paper5:lem:log-scale}, for any \(\tau_2>\tau_1\),
  \[
    \int_{\tau_1}^{\tau_2}a(\tau)
    \,d\tau
    =\log\lambda(\tau_2)-\log\lambda(\tau_1).
  \]
  Taking \(\tau_2\to\infty\) and using \(\lambda(\tau_2)\to0\),
  \[
    \lim_{\tau_2\to\infty}\int_{\tau_1}^{\tau_2}a(\tau)
    \,d\tau=-\infty.
  \]

  \item Decompose \(a=a_+-a_-\).  Then for each \(\tau_2>\tau_1\),
  \[
    \int_{\tau_1}^{\tau_2}a(\tau)
    \,d\tau
    =\int_{\tau_1}^{\tau_2}a_+(\tau)
    \,d\tau-\int_{\tau_1}^{\tau_2}a_-(\tau)
    \,d\tau
    \ge-\int_{\tau_1}^{\tau_2}a_-(\tau)
    \,d\tau.
    \]
    Hence
    \[
    \int_{\tau_1}^{\tau_2}a_-(\tau)
    \,d\tau\ge
    -\int_{\tau_1}^{\tau_2}a(\tau)
    \,d\tau.
    \]

  \item Taking limits and using Step~2 yields
  \[
    \int_{\tau_1}^{\infty}a_-(\tau)
    \,d\tau=\infty.
  \]
  Indeed, by Step~2,
  \[
    -\int_{\tau_1}^{\tau_2}a(\tau)\,d\tau\to+\infty.
  \]
  Since the partial integrals of \(a_-\) dominate this quantity, they are unbounded.
  Because \(a_-\ge0\), monotone convergence gives
  \(\int_{\tau_1}^{\infty}a_-(\tau)\,d\tau=\infty\).

  \item Split the cost into early and late times:
  \[
    \int_{\tau_0}^{\infty}a_-M_R
    \,d\tau
    =\int_{\tau_0}^{\tau_1}a_-M_R
    \,d\tau+
    \int_{\tau_1}^{\infty}a_-M_R
    \,d\tau.
  \]
  The first term is finite or infinite independently of the argument; the second term
  satisfies
  \[
    \int_{\tau_1}^{\infty}a_-(\tau)M_R(\tau)
    \,d\tau
    \ge m_0\int_{\tau_1}^{\infty}a_-(\tau)
    \,d\tau
    =\infty
  \]
  by the bound in (1) and Step~4.  Hence the total integral is infinite.
\end{enumerate}
\end{proof}

\subsection{Finite-cost exclusions}
\label{paper5:sec:finite-cost}

\begin{theorem}[Scale-collapse cost exclusion]\label{paper5:thm:cost-exclusion}
No branch with finite scale-collapsing cost (Definition~\ref{paper5:def:finite-cost-branch}) can satisfy simultaneously
\(\lambda(\tau)\to0\) and nonvanishing localized core.
\end{theorem}

\begin{proof}
Assume, for contradiction, that such a branch exists.  The two hypotheses are exactly
those of Theorem~\ref{paper5:thm:collapse-core-floor}; therefore
\[
  \int_{\tau_0}^{\infty}\mathcal C_{\mathrm{sc}}(\tau)
  \,d\tau=\infty.
\]
This contradicts finite scale-collapsing cost.  Therefore no such branch exists.
\end{proof}

\begin{lemma}[Absolute-cost domination]\label{paper5:lem:absolute-cost}
For a.e. \(\tau\in[\tau_0,\infty)\),
\[
  \mathcal C_{\mathrm{abs}}(\tau)\ge\mathcal C_{\mathrm{sc}}(\tau).
\]
Hence, pointwise and after integrating in \(\tau\),
\[
  \int_{\tau_0}^{\infty}\mathcal C_{\mathrm{sc}}(\tau)\,d\tau
  \le
  \int_{\tau_0}^{\infty}\mathcal C_{\mathrm{abs}}(\tau)\,d\tau,
\]
as extended real numbers.
In particular, finite absolute cost implies finite scale-collapsing cost, and
equivalently
\[
  \int_{\tau_0}^{\infty}\mathcal C_{\mathrm{sc}}(\tau)\,d\tau=\infty
  \quad\Longrightarrow\quad
  \int_{\tau_0}^{\infty}\mathcal C_{\mathrm{abs}}(\tau)\,d\tau=\infty.
\]
\end{lemma}

\begin{proof}
At each time \(\tau\),
\[
  \mathcal C_{\mathrm{abs}}(\tau)
  =\nu\int|\nabla V|^2\Phi
  +(a_++a_-)M_R
  \ge(a_++a_-)M_R\ge a_-M_R
  =\mathcal C_{\mathrm{sc}}(\tau),
\]
since \(a_+,a_-\ge0\) and \(\nu\int|\nabla V|^2\Phi\ge0\).  Integrating over
\([\tau_0,\infty)\) yields the claimed implication.
\end{proof}

\begin{corollary}[Absolute scale-cost exclusion]\label{paper5:cor:absolute-cost-exclusion}
No branch with finite absolute cost in the sense of Definition~\ref{paper5:def:finite-cost-branch}
can satisfy simultaneously \(\lambda(\tau)\to0\) and nonvanishing localized core.
\end{corollary}

\begin{proof}
Assume instead such a branch exists.  By Theorem~\ref{paper5:thm:collapse-core-floor},
\(\int\mathcal C_{\mathrm{sc}}=\infty\).  By Lemma~\ref{paper5:lem:absolute-cost}, the
absolute cost is pointwise above \(\mathcal C_{\mathrm{sc}}\), so its integral must also be
infinite, contradicting finite absolute cost.
\end{proof}

\begin{remark}
The exclusions above are purely integral and use neither compactness nor any
autonomous-limit argument.
\end{remark}

\subsection{Autonomous reduced-limit theorem from the selected-window estimates}
\label{paper5:sec:autonomous}

The preceding exclusions are integral statements and do not rely on compactness.  On a fixed finite autonomous window, compactness requires not only the cost bound but also the selected-window estimates supplied by the preceding renormalized representation, pressure reconstruction, and local Caccioppoli estimates.

\medskip
\noindent\textit{Analytic inputs.}
The preceding sections supply the following ingredients on selected compact
renormalized windows:
\begin{enumerate}[label=(\arabic*)]
  \item the repaired-gauge renormalized equation and weak formulation for \((V,P)\);
  \item pressure reconstruction modulo spatial means, with local \(L^{3/2}\) bounds on compact
    cylinders;
  \item local windowed bounds
  \(V\in L^\infty_tL^2_x\cap L^2_tH^1_x\) on compact cylinders;
  \item a retained compact core lower bound preventing the selected profile from vanishing.
\end{enumerate}
Together with the thick-window convergence of the modulation coefficients, these estimates imply
the compactness and limit passage below.

\begin{definition}[Thick autonomous windows]\label{paper5:def:thick-autonomous-windows}
A branch has \emph{thick autonomous windows} if there exist constants
\(T>0\), \(c>0\), \(M_1\in(0,\infty)\), a sequence \((\tau_n)\) with
\(\tau_n\to\infty\), and limits \(a_\infty\in\mathbb R\),
\(b_\infty\in\mathbb R^3\), such that for every \(n\):
\begin{enumerate}[label=(\roman*)]
  \item
  \[
    \frac1T\int_{\tau_n}^{\tau_n+T}a_-(\tau)M_R(\tau)
    \,d\tau\ge c;
  \]
  \item
  \[
    0\le M_R(\tau)\le M_1\quad\text{for a.e. }\tau\in[\tau_n,\tau_n+T];
  \]
  \item
  \[
    a(\tau_n+\cdot)\to a_\infty\ \text{in }L^1(0,T),
    \qquad
    b(\tau_n+\cdot)\to b_\infty\ \text{in }L^1(0,T;\mathbb R^3).
  \]
\end{enumerate}
\end{definition}
\begin{remark}
This notion isolates a limiting drift coefficient \(a_\infty<0\).  The
compactness to extract a profile limit is obtained from the selected-window compactness estimates,
not from the scalar cost inequalities of Sections~\ref{paper5:sec:cost}--\ref{paper5:sec:finite-cost}.
\end{remark}

\begin{lemma}[Thick drift implies negative autonomous limit parameter]\label{paper5:lem:negative-parameter}
On thick autonomous windows, the limiting autonomous coefficient is negative:
\[
  a_\infty\le-\frac{c}{M_1}<0.
\]
\end{lemma}

\begin{proof}
\begin{enumerate}[label=(\arabic*)]
  \item From (ii), for a.e. \(\tau\in[\tau_n,\tau_n+T]\),
  \(0\le M_R(\tau)\le M_1\).  Therefore
  \[
    \frac{a_-(\tau)M_R(\tau)}{M_1}
    \le a_-(\tau).
  \]
  Therefore
  \[
    \frac1T\int_{\tau_n}^{\tau_n+T}a_-(\tau)
    \,d\tau
    \ge \frac1{M_1T}\int_{\tau_n}^{\tau_n+T}a_-(\tau)M_R(\tau)
    \,d\tau
    \ge \frac{c}{M_1}.
  \]

  \item Set \(a_n(s):=a(\tau_n+s)\).  By (iii), \(a_n\to a_\infty\) in
  \(L^1(0,T)\). The map \(r\mapsto r^-\) is 1-Lipschitz, so
  \(a_n^-\to a_\infty^-\) in \(L^1(0,T)\):
  \[
    \|a_n^--a_\infty^-\|_{L^1(0,T)}\le
    \|a_n-a_\infty\|_{L^1(0,T)}\to0.
  \]
  Thus
  \[
  \lim_{n\to\infty}\frac1T\int_{\tau_n}^{\tau_n+T}a_-(\tau)
    \,d\tau
    =\frac1T\int_0^T(a_\infty)^-
    \,ds=(a_\infty)^-.
  \]

  \item Combine Step~1 and Step~2:
  \((a_\infty)^-\ge c/M_1\), i.e.\ \(a_\infty\le-c/M_1<0\).
\end{enumerate}
\end{proof}

\begin{assumption}[Selected-window suitable estimates]\label{paper5:ass:selected-window-estimates}
Let \((\tau_n)\) be a thick autonomous-window sequence and set
\[
  V_n(y,s):=V(y,\tau_n+s),\qquad
  P_n(y,s):=P(y,\tau_n+s),
\]
\[
  \alpha_n(s):=a(\tau_n+s),\qquad
  \beta_n(s):=b(\tau_n+s),
  s\in(0,T).
\]
For every bounded Lipschitz domain \(K\Subset\mathbb R^3\), assume there is a constant
\(C_K<\infty\), independent of \(n\), such that
\[
  \|V_n\|_{L^\infty(0,T;L^2(K))}
  +\|V_n\|_{L^2(0,T;H^1(K))}
  +\|P_n-(P_n)_K\|_{L^{3/2}((0,T)\times K)}
  \le C_K,
\]
where
\[
  (P_n)_K(s):=\frac1{|K|}\int_KP_n(y,s)\,dy.
\]
Assume also that each \(V_n\) is divergence free and that there exist
\(\chi\in C_c^\infty(B_R)\), \(\chi\ge0\), \(\eta>0\), and \(N\in\mathbb N\) such that
for all \(n\ge N\),
\[
  \int_{\mathbb R^3}|V_n(y,s)|^2\chi(y)\,dy\ge\eta
  \quad\text{for a.e. }s\in(0,T).
\]
\end{assumption}

\begin{proposition}[Derivation of the selected-window estimates from the preceding sections]
\label{paper5:prop:selected-window-derived}
Let a repaired-gauge branch satisfy the renormalized representation and pressure reconstruction
of the compact-window representation theorem, \cref{paper2:thm:C}, on the windows \([\tau_n,\tau_n+T]\).  Assume that on
each compact cylinder \((0,T)\times K\) the local compact-cylinder Caffarelli--Kohn--Nirenberg
quantities are finite uniformly in \(n\), and that the selected core has the localized lower
bound stated in Assumption~\ref{paper5:ass:selected-window-estimates}.  Then
Assumption~\ref{paper5:ass:selected-window-estimates} holds.
\end{proposition}

\begin{proof}
The renormalized representation theorem supplies the shifted equation, divergence-free condition,
and pressure reconstruction modulo functions of time.  The pressure reconstruction gives, after
subtracting spatial means,
\[
  \sup_n\|P_n-(P_n)_K\|_{L^{3/2}((0,T)\times K)}<\infty
\]
on every compact \(K\).  The local Caccioppoli estimate applied on a slightly larger compact
cylinder gives the inner bounds
\[
  \sup_n\|V_n\|_{L^\infty(0,T;L^2(K))}
  +\sup_n\|V_n\|_{L^2(0,T;H^1(K))}<\infty.
\]
These are precisely the local velocity and pressure estimates in
Assumption~\ref{paper5:ass:selected-window-estimates}.  The retained selected-core lower bound is the
last condition in that assumption.  Hence all items in
Assumption~\ref{paper5:ass:selected-window-estimates} follow from the representation,
pressure, and local Caccioppoli results.
\end{proof}

\begin{lemma}[Local time regularity on autonomous windows]\label{paper5:lem:window-time-regularity}
Assume thick autonomous windows and Assumption~\ref{paper5:ass:selected-window-estimates}.  Let
\(K_0\Subset K_1\Subset\mathbb R^3\) be bounded Lipschitz domains, and choose an integer
\(m\) so large that \(H^m_0(K_1)\hookrightarrow W^{1,\infty}_0(K_1)\).  Then
\[
  \{\partial_sV_n\}_{n\ge1}
  \quad\text{is bounded in }L^1(0,T;H^{-m}(K_1)).
\]
\end{lemma}

\begin{proof}
Fix \(\zeta\in C_c^\infty(K_1)\) with \(\zeta\equiv1\) on a neighborhood of \(K_0\).
Let \(\phi\in H^m_0(K_1;\mathbb R^3)\).  Since \(H^m_0(K_1)\hookrightarrow W^{1,\infty}_0(K_1)\),
there is \(C_m\) such that
\[
  \|\phi\|_{L^\infty(K_1)}+\|\nabla\phi\|_{L^\infty(K_1)}
  \le C_m\|\phi\|_{H^m(K_1)}.
\]
The shifted equation is
\[
  \partial_sV_n
  =-(V_n\cdot\nabla)V_n-\nabla P_n+\nu\Delta V_n
  +\alpha_n(V_n+y\cdot\nabla V_n)+\beta_n\cdot\nabla V_n
\]
in distributions.  We estimate each term against \(\phi\).

For the convection term, integration by parts gives
\[
  \left|\langle (V_n\cdot\nabla)V_n,\phi\rangle\right|
  =
  \left|\int_{K_1}V_{n,i}V_{n,j}\partial_j\phi_i\,dy\right|
  \le C_m\|V_n(s)\|_{L^2(K_1)}^2\|\phi\|_{H^m(K_1)}.
\]
For the viscous term,
\[
  \left|\langle \nu\Delta V_n,\phi\rangle\right|
  =
  \nu\left|\int_{K_1}\nabla V_n:\nabla\phi\,dy\right|
  \le \nu C_m\|\nabla V_n(s)\|_{L^2(K_1)}|K_1|^{1/2}\|\phi\|_{H^m(K_1)}.
\]
For the pressure term, subtracting the spatial mean does not change the gradient:
\[
  \left|\langle\nabla P_n,\phi\rangle\right|
  =
  \left|\int_{K_1}(P_n-(P_n)_{K_1})\,\nabla\cdot\phi\,dy\right|
  \le C_m |K_1|^{1/3}
  \|P_n-(P_n)_{K_1}\|_{L^{3/2}(K_1)}
  \|\phi\|_{H^m(K_1)}.
\]
For the \(V_n\)-part of the scaling drift,
\[
  |\alpha_n(s)|\,\left|\int_{K_1}V_n\cdot\phi\,dy\right|
  \le C_m|\alpha_n(s)|\,|K_1|^{1/2}
  \|V_n(s)\|_{L^2(K_1)}\|\phi\|_{H^m(K_1)}.
\]
For the \(y\cdot\nabla V_n\)-part, integrate by parts:
\[
  \int_{K_1}(y\cdot\nabla V_n)\cdot\phi\,dy
  =
  -\int_{K_1}V_n\cdot(y\cdot\nabla\phi+3\phi)\,dy,
\]
because \(\phi\) has zero trace.  Hence
\[
  |\alpha_n(s)|\,\left|\int_{K_1}(y\cdot\nabla V_n)\cdot\phi\,dy\right|
  \le C_{K_1,m}|\alpha_n(s)|\|V_n(s)\|_{L^2(K_1)}\|\phi\|_{H^m(K_1)}.
\]
Similarly,
\[
  \left|\int_{K_1}(\beta_n\cdot\nabla V_n)\cdot\phi\,dy\right|
  =
  \left|\int_{K_1}V_n\cdot(\beta_n\cdot\nabla\phi)\,dy\right|
  \le C_m|\beta_n(s)|\,|K_1|^{1/2}
  \|V_n(s)\|_{L^2(K_1)}\|\phi\|_{H^m(K_1)}.
\]
Combining the estimates and taking the supremum over
\(\|\phi\|_{H^m(K_1)}\le1\), we obtain
\[
\begin{aligned}
  \|\partial_sV_n(s)\|_{H^{-m}(K_1)}
  \le C_{K_1,m}\big(&
  \|V_n(s)\|_{L^2(K_1)}^2
  +\|\nabla V_n(s)\|_{L^2(K_1)}\\
  &+\|P_n(s)-(P_n)_{K_1}(s)\|_{L^{3/2}(K_1)}
  +( |\alpha_n(s)|+|\beta_n(s)| )\|V_n(s)\|_{L^2(K_1)}
  \big).
\end{aligned}
\]
Integrating in \(s\in(0,T)\), using Assumption~\ref{paper5:ass:selected-window-estimates}, and using
\(\alpha_n\to a_\infty\), \(\beta_n\to b_\infty\) in \(L^1\), hence uniform \(L^1\)-bounds
for \(\alpha_n,\beta_n\), proves the claimed \(L^1(0,T;H^{-m})\) bound.
\end{proof}

\begin{lemma}[Local compactness on autonomous windows]\label{paper5:lem:autonomous-window-compactness}
Assume thick autonomous windows and Assumption~\ref{paper5:ass:selected-window-estimates}.  Then after
passing to a subsequence there exist
\[
  U\in L^\infty(0,T;L^2_{\mathrm{loc}}(\mathbb R^3))
  \cap L^2(0,T;H^1_{\mathrm{loc}}(\mathbb R^3)),
  \qquad
  \Pi\in L^{3/2}_{\mathrm{loc}}((0,T)\times\mathbb R^3),
\]
such that, for every compact \(K\Subset\mathbb R^3\),
\begin{enumerate}[label=(\alph*)]
  \item
  \[
    V_n\to U
    \quad\text{strongly in }L^2(0,T;L^2(K));
  \]
  \item
  \[
    V_n\rightharpoonup U
    \quad\text{weakly in }L^2(0,T;H^1(K));
  \]
  \item
  \[
    V_n\rightharpoonup^\ast U
    \quad\text{weakly-* in }L^\infty(0,T;L^2(K));
  \]
  \item
  \[
    P_n-(P_n)_K\rightharpoonup\Pi_K
    \quad\text{weakly in }L^{3/2}((0,T)\times K),
  \]
  where the local distributions \(\nabla\Pi_K\) agree on overlaps and define
  a pressure gradient \(\nabla\Pi\) modulo functions of time.
\end{enumerate}
\end{lemma}

\begin{proof}
Fix \(K_0\Subset K_1\Subset\mathbb R^3\).  Assumption~\ref{paper5:ass:selected-window-estimates}
gives boundedness of \(V_n\) in \(L^2(0,T;H^1(K_1))\), and
Lemma~\ref{paper5:lem:window-time-regularity} gives boundedness of \(\partial_sV_n\) in
\(L^1(0,T;H^{-m}(K_1))\) for \(m\) large.  Since
\[
  H^1(K_1)\Subset L^2(K_0)\hookrightarrow H^{-m}(K_1),
\]
the Aubin--Lions--Simon compactness theorem implies, after extracting a subsequence,
\[
  V_n\to U
  \quad\text{strongly in }L^2(0,T;L^2(K_0)).
\]
The local \(L^\infty L^2\) and \(L^2H^1\) bounds also give, after extraction,
\[
  V_n\rightharpoonup^\ast U
  \quad\text{in }L^\infty(0,T;L^2(K_0)),
  \qquad
  V_n\rightharpoonup U
  \quad\text{in }L^2(0,T;H^1(K_0)).
\]
The pressure bound in Assumption~\ref{paper5:ass:selected-window-estimates} gives weak compactness of
\(P_n-(P_n)_{K_0}\) in \(L^{3/2}((0,T)\times K_0)\).  A diagonal extraction over balls
\(B_j\) gives a single subsequence for every compact \(K\).  On overlaps, two pressure limits
can differ only by a function of \(s\), because both represent the same distributional pressure
gradient in the limit equation.  Thus the gradients assemble into \(\nabla\Pi\), with \(\Pi\)
defined modulo functions of time.
\end{proof}

\begin{lemma}[Compatibility of local pressure limits]\label{paper5:lem:pressure-compatibility}
Let \(K_1,K_2\Subset\mathbb R^3\) be bounded Lipschitz domains with
nonempty connected overlap \(K_{12}:=K_1\cap K_2\).  Suppose that, after passing to a common
subsequence,
\[
  P_n-(P_n)_{K_i}\rightharpoonup \Pi_i
  \quad\text{in }L^{3/2}((0,T)\times K_i),
  \qquad i=1,2.
\]
Then there exists \(c_{12}\in L^{3/2}(0,T)\) such that
\[
  \Pi_1(y,s)-\Pi_2(y,s)=c_{12}(s)
  \quad\text{for a.e. }(y,s)\in K_{12}\times(0,T).
\]
Consequently the pressure gradient \(\nabla\Pi\) obtained from the local limits is globally
well-defined on \((0,T)\times\mathbb R^3\), with the pressure itself defined modulo functions
of time.
\end{lemma}

\begin{proof}
For each \(n\),
\[
  \bigl(P_n-(P_n)_{K_1}\bigr)-\bigl(P_n-(P_n)_{K_2}\bigr)
  =(P_n)_{K_2}(s)-(P_n)_{K_1}(s)
\]
on \(K_{12}\).  Hence its spatial gradient is zero in
\(\mathcal D'((0,T)\times K_{12})\).  Passing to the weak limit gives
\[
  \nabla_y(\Pi_1-\Pi_2)=0
  \quad\text{in }\mathcal D'((0,T)\times K_{12}).
\]
Since \(K_{12}\) is connected, the standard distributional characterization of functions with
zero spatial gradient implies that \(\Pi_1-\Pi_2=c_{12}(s)\) for some
\(c_{12}\in L^{3/2}(0,T)\).  The last assertion follows because adding a time-dependent scalar
to pressure does not change \(\nabla\Pi\).
\end{proof}

\begin{theorem}[Autonomous reduced-limit theorem]\label{paper5:thm:autonomous-limit}
Assume that a branch has thick autonomous windows and satisfies
Assumption~\ref{paper5:ass:selected-window-estimates} on the same window sequence.  Let
\((U,\Pi)\) be the limit extracted in Lemma~\ref{paper5:lem:autonomous-window-compactness}.  Then
\[
  \nabla\cdot U=0,
\]
\[
  U\not\equiv0
\]
and \((U,\Pi)\) satisfies
\begin{equation}\label{paper5:eq:autonomous-limit}
  \partial_sU+(U\cdot\nabla)U+\nabla\Pi
  =\nu\Delta U+a_\infty(U+y\cdot\nabla U)+b_\infty\cdot\nabla U,
  \qquad
  \nabla\cdot U=0,
\end{equation}
in \(\mathcal D'((0,T)\times\mathbb R^3)\).  Equivalently, for every
\(\varphi\in C_c^\infty((0,T)\times\mathbb R^3;\mathbb R^3)\),
\[
  \int_0^T\!\int
  \bigl(
  -U\cdot\partial_s\varphi
  -(U\otimes U):\nabla\varphi
  +\nu\nabla U:\nabla\varphi
  -\Pi\,\nabla\cdot\varphi
  -a_\infty(U+y\cdot\nabla U)\cdot\varphi
  -(b_\infty\cdot\nabla U)\cdot\varphi
  \bigr)\,dy\,ds=0.
\]
\end{theorem}

\begin{remark}
The theorem asserts only the compact finite-window autonomous limit from the
selected-window estimates.  A rigidity contradiction requires a separate
Liouville theorem for the class containing \((U,\Pi)\).
\end{remark}

\begin{proof}
\emph{Step 1: divergence-free property of the limit.}

Each \(V_n\) is divergence free by Assumption~\ref{paper5:ass:selected-window-estimates}, so for any compactly supported scalar
\(\psi\in C_c^\infty((0,T)\times\mathbb R^3)\),
\[
  \int_{(0,T)\times\mathbb R^3}V_n\cdot\nabla\psi\,dy\,ds=0.
\]
Hence, for every such \(\psi\), since \(\nabla\psi\in C_c^\infty((0,T)\times\mathbb R^3)\) and \(V_n\to U\) in
\(L^2(0,T;L^2_{\mathrm{loc}})\),
\[
  \int_{(0,T)\times\mathbb R^3}(V_n-U)\cdot\nabla\psi\,dy\,ds\to0,
\]
we deduce
\[
  \int_{(0,T)\times\mathbb R^3}U\cdot\nabla\psi\,dy\,ds=0.
\]
Hence \(\nabla\cdot U=0\) in \(\mathcal D'((0,T)\times\mathbb R^3)\).


\emph{Step 2: nontriviality of the limit profile.}
Define for \(s\in(0,T)\)
\[
  f_n(s):=\int_{\mathbb R^3}|V_n(y,s)|^2\chi(y)
  \,dy,
  \qquad
  f(s):=\int_{\mathbb R^3}|U(y,s)|^2\chi(y)
  \,dy.
\]
Here \(K_\chi:=\operatorname{supp}\chi\), so
\(\chi^{1/2}V_n,\chi^{1/2}U\in L^2(0,T;L^2(K_\chi))\).
Then
\[
  f_n-f=\int_{\mathbb R^3}(V_n-U)\cdot(V_n+U)\chi
  \,dy.
\]
Applying Cauchy--Schwarz in space,
\[
  |f_n(s)-f(s)|\le
  \|\chi^{1/2}(V_n(s)-U(s))\|_{L^2(\mathbb R^3)}
  \|\chi^{1/2}(V_n(s)+U(s))\|_{L^2(\mathbb R^3)}.
\]
Integrating in \(s\) and applying Cauchy--Schwarz in time gives
\[
  \|f_n-f\|_{L^1(0,T)}
  \le
  \|\chi^{1/2}(V_n-U)\|_{L^2(0,T;L^2)}
  \|\chi^{1/2}(V_n+U)\|_{L^2(0,T;L^2)}.
\]
Since \(K_\chi\) is compact, Assumption~\ref{paper5:ass:selected-window-estimates} and
the membership \(U\in L^2(0,T;L^2(K_\chi))\) give
\[
  \sup_n\|\chi^{1/2}(V_n+U)\|_{L^2(0,T;L^2)}
  \le C_\chi\left(\sup_n\|V_n\|_{L^2(0,T;L^2(K_\chi))}
  +\|U\|_{L^2(0,T;L^2(K_\chi))}\right)<\infty.
\]
Moreover
\[
  \|\chi^{1/2}(V_n-U)\|_{L^2(0,T;L^2)}
  \le \|\chi\|_{L^\infty}^{1/2}
  \|V_n-U\|_{L^2(0,T;L^2(K_\chi))}\to0.
\]
Therefore
\[
  \|f_n-f\|_{L^1(0,T)}\to0.
\]
Hence \(f_n\to f\) in \(L^1(0,T)\).

Assumption~\ref{paper5:ass:selected-window-estimates} gives for large \(n\),
\(f_n(s)\ge\eta\) for a.e.\(s\). Therefore
\[
  \int_0^T f(s)
  \,ds=\lim_{n\to\infty}\int_0^Tf_n(s)
  \,ds\ge\eta T>0.
\]
Hence \(f\not\equiv0\), so \(U\not\equiv0\).

\emph{Step 3: shifted weak formulation for \(V_n\).}
Take any test field
\(\varphi\in C_c^\infty((0,T)\times\mathbb R^3;\mathbb R^3)\), and set
\[
K_\varphi:=\overline{\{x\in\mathbb R^3:\exists s\in(0,T),\ \varphi(x,s)\neq0\}}
\subset\mathbb R^3,\qquad K_\varphi\ \text{bounded}.
\]
Since \((V,P)\) satisfies~\eqref{paper5:eq:renorm-ns} in distributions, the change of variables
\(\tau=\tau_n+s\) gives the shifted weak formulation
\begin{equation}\label{paper5:eq:weak-n}
\begin{aligned}
0={}&\int_0^T\!\int_{\mathbb R^3}
\Bigl(
-V_n\cdot\partial_s\varphi
-(V_n\otimes V_n):\nabla\varphi
+\nu\nabla V_n:\nabla\varphi
-P_n\nabla\cdot\varphi
\Bigr)\,dy\,ds\\
&+\int_0^T\!\int_{\mathbb R^3}
\Bigl(
-a(\tau_n+s)(V_n+y\cdot\nabla V_n)\cdot\varphi
-(b(\tau_n+s)\cdot\nabla V_n)\cdot\varphi
\Bigr)\,dy\,ds.
\end{aligned}
\end{equation}

\emph{Step 4: term-by-term passage to the limit.}
Write \(\alpha_n(s):=a(\tau_n+s)\), \(\beta_n(s):=b(\tau_n+s)\).

\medskip
\noindent\textit{(i) Time derivative.}

By strong convergence in \(L^2(0,T;L^2(K_\varphi))\),
\[
  \left|\int_0^T\!\int (V_n-U)\cdot\partial_s\varphi
  \,dy\,ds\right|
  \le
  \|V_n-U\|_{L^2(0,T;L^2(K_\varphi))}
  \|\partial_s\varphi\|_{L^2(0,T;L^2(K_\varphi))}
  \to0.
\]
Hence
\[\int V_n\cdot\partial_s\varphi\to\int U\cdot\partial_s\varphi.
\]

\medskip
\noindent\textit{(ii) Convection term.}
Decompose
\[
  V_n\otimes V_n-U\otimes U=(V_n-U)\otimes V_n+U\otimes(V_n-U).
\]
Then
\[
\begin{aligned}
&\left|\int_0^T\!\int((V_n\otimes V_n-U\otimes U):\nabla\varphi)
  \,dy\,ds\right|\\
&\quad\le
  \|\nabla\varphi\|_{L^\infty_{x,t}}
  \|V_n-U\|_{L^2(0,T;L^2(K_\varphi))}
  \bigl(\|V_n\|_{L^2(0,T;L^2(K_\varphi))}
  +\|U\|_{L^2(0,T;L^2(K_\varphi))}\bigr)
  \to0,
\end{aligned}
\]
using boundedness of \(V_n\) in \(L^2(0,T;L^2(K_\varphi))\). Therefore
\(
  \int(V_n\otimes V_n):\nabla\varphi
  \to\int(U\otimes U):\nabla\varphi.
\)

\medskip
\noindent\textit{(iii) Viscous term.}
From weak convergence of \(V_n\) in \(L^2(0,T;H^1_{\mathrm{loc}})\),
\[
  \int_0^T\!\int\nabla V_n:\nabla\varphi
  \,dy\,ds
  \to
  \int_0^T\!\int\nabla U:\nabla\varphi
  \,dy\,ds.
\]

\medskip
\noindent\textit{(iv) Pressure term.}
Let \(K\) be a bounded Lipschitz domain with \(K_\varphi\Subset K\).  Since
\(\varphi\) is compactly supported in \(K\),
\[
  \int_K (P_n)_K(s)\,\nabla\cdot\varphi(y,s)\,dy
  =
  (P_n)_K(s)\int_K\nabla\cdot\varphi(y,s)\,dy=0.
\]
Therefore
\[
  \int_0^T\!\int_{\mathbb R^3}P_n\nabla\cdot\varphi\,dy\,ds
  =
  \int_0^T\!\int_K(P_n-(P_n)_K)\nabla\cdot\varphi\,dy\,ds.
\]
By Lemma~\ref{paper5:lem:autonomous-window-compactness},
\[
  P_n-(P_n)_K\rightharpoonup\Pi_K
  \quad\text{in }L^{3/2}((0,T)\times K),
\]
and \(\nabla\cdot\varphi\in L^3((0,T)\times K)\).  The compatibility
Lemma~\ref{paper5:lem:pressure-compatibility} permits us to write the resulting pressure distribution
as \(\Pi\), modulo functions of time.  Hence
\[
  \int_0^T\!\int_{\mathbb R^3}P_n\nabla\cdot\varphi\,dy\,ds
  \to
  \int_0^T\!\int_K\Pi_K\nabla\cdot\varphi\,dy\,ds.
\]
The value of the last integral is independent of the chosen \(K\), because replacing
\(\Pi_K\) by \(\Pi_K+c(s)\) changes the integral by
\(\int_0^Tc(s)\int_K\nabla\cdot\varphi\,dy\,ds=0\).  The limit is therefore the pressure term
\(\int\Pi\nabla\cdot\varphi\).

\medskip
\noindent\textit{(v) Drift term with \(a\): part involving \(V_n\).}
Define
\[
  Y_n^0(s):=\int_{\mathbb R^3}V_n(y,s)\cdot\varphi(y,s)
  \,dy,
  \qquad
  Y^0(s):=\int_{\mathbb R^3}U(y,s)\cdot\varphi(y,s)
  \,dy.
\]
For each \(s\), Cauchy--Schwarz gives
\[
  |Y_n^0(s)|
  \le
  \|V_n(s)\|_{L^2(K_\varphi)}
  \|\varphi(s)\|_{L^2(K_\varphi)}.
\]
By Assumption~\ref{paper5:ass:selected-window-estimates}, there is \(C_0<\infty\)
such that \(\sup_n\|Y_n^0\|_{L^\infty(0,T)}\le C_0\). Moreover,
\[
  Y_n^0\to Y^0\quad\text{in }L^2(0,T),
\]
because \(V_n\to U\) in \(L^2(0,T;L^2(K_\varphi))\).
Now
\[
  \int_0^T\alpha_n(s)Y_n^0(s)
  \,ds
  =\int_0^T(\alpha_n(s)-a_\infty)Y_n^0(s)
  \,ds+a_\infty\int_0^TY_n^0(s)
  \,ds.
\]
For the first term,
\[
  \left|\int_0^T(\alpha_n-a_\infty)Y_n^0\,ds\right|\le
\sup_{s\in(0,T)}|Y_n^0(s)|\,\|\alpha_n-a_\infty\|_{L^1(0,T)}\to0.
\]
For the second term, strong convergence in \(L^2(0,T)\) implies
\(\int_0^TY_n^0\to\int_0^TY^0\), so
\[
  \int_0^T\alpha_n(s)\int_{\mathbb R^3}V_n\cdot\varphi
  \,dy\,ds
  \to a_\infty\int_0^T\int_{\mathbb R^3}U\cdot\varphi
  \,dy\,ds.
\]
Thus the signed contribution in \eqref{paper5:eq:weak-n} satisfies
\[
  -\int_0^T\alpha_n(s)\int_{\mathbb R^3}V_n\cdot\varphi
  \,dy\,ds
  \to -a_\infty\int_0^T\int_{\mathbb R^3}U\cdot\varphi
  \,dy\,ds.
\]

\medskip
\noindent\textit{(vi) Drift term with \(a\): part involving \(y\cdot\nabla V_n\).}
Since \(\varphi\in C_c^\infty\), integrating by parts in \(y\) gives
\[
  \int_{\mathbb R^3}(y\cdot\nabla V_n)\cdot\varphi
  =-\int_{\mathbb R^3}V_n\cdot(y\cdot\nabla\varphi+3\varphi)
  \,dy.
\]
Define
\[
  Y_n^1(s):=\int_{\mathbb R^3}V_n\cdot(y\cdot\nabla\varphi+3\varphi)
  \,dy,
  \qquad
  Y^1(s):=\int_{\mathbb R^3}U\cdot(y\cdot\nabla\varphi+3\varphi)
  \,dy.
\]
For each \(s\), Cauchy--Schwarz gives
\[
  |Y_n^1(s)|
  \le
  \|V_n(s)\|_{L^2(K_\varphi)}
  \|y\cdot\nabla\varphi(s)+3\varphi(s)\|_{L^2(K_\varphi)}.
\]
The \(L^\infty(0,T;L^2(K_\varphi))\) bound in
Assumption~\ref{paper5:ass:selected-window-estimates} gives
\[
  \sup_n\|Y_n^1\|_{L^\infty(0,T)}\le C_1<\infty.
\]
Also,
\[
  \|Y_n^1-Y^1\|_{L^2(0,T)}
  \le
  \|V_n-U\|_{L^2(0,T;L^2(K_\varphi))}
  \|y\cdot\nabla\varphi+3\varphi\|_{L^\infty(0,T;L^2(K_\varphi))}
  \to0.
\]
Therefore
\[
  -\int_0^T\alpha_n(s)
  \int_{\mathbb R^3}(y\cdot\nabla V_n)\cdot\varphi
  \,dy\,ds
  \to
  a_\infty\int_0^T\!\int_{\mathbb R^3}U\cdot(y\cdot\nabla\varphi+3\varphi)
  \,dy\,ds.
\]
Integrating by parts once more in the limit gives
\[
  a_\infty\int_0^T\!\int_{\mathbb R^3}U\cdot(y\cdot\nabla\varphi+3\varphi)
  \,dy\,ds
  =
  -a_\infty\int_0^T\!\int_{\mathbb R^3}(y\cdot\nabla U)\cdot\varphi
  \,dy\,ds.
\]

\medskip
\noindent\textit{(vii) Term with \(b\).}
Since \(\beta_n(s)\) is independent of \(y\),
\[
  \int_{\mathbb R^3}(\beta_n\cdot\nabla V_n)\cdot\varphi
  =-\int_{\mathbb R^3}V_n\cdot(\beta_n\cdot\nabla\varphi)
  \,dy,
\]
so
\[
  -\int_0^T\!\int(\beta_n\cdot\nabla V_n)\cdot\varphi
  \,dy\,ds
  =
  \int_0^T\!\int V_n\cdot(\beta_n\cdot\nabla\varphi)
  \,dy\,ds.
\]
Let
\[
\begin{aligned}
Z_n(s)&:=\left(\int_{\mathbb R^3}V_n\cdot\partial_1\varphi\,dy,\,
\int_{\mathbb R^3}V_n\cdot\partial_2\varphi\,dy,\,
\int_{\mathbb R^3}V_n\cdot\partial_3\varphi\,dy\right),\\
Z(s)&:=\left(\int_{\mathbb R^3}U\cdot\partial_1\varphi\,dy,\,
\int_{\mathbb R^3}U\cdot\partial_2\varphi\,dy,\,
\int_{\mathbb R^3}U\cdot\partial_3\varphi\,dy\right).
\end{aligned}
\]
Then
\[
  \int_{\mathbb R^3}V_n\cdot(\beta_n\cdot\nabla\varphi)\,dy
  =\beta_n(s)\cdot Z_n(s),\qquad
  \int_{\mathbb R^3}U\cdot(b_\infty\cdot\nabla\varphi)\,dy
  =b_\infty\cdot Z(s).
\]
Consequently
\[
  \left|\int_0^T\!(\beta_n-b_\infty)\cdot Z_n\,ds\right|
  \le
  \|\beta_n-b_\infty\|_{L^1(0,T)}\sup_{s\in(0,T)}|Z_n(s)|\to0,
\]
because
\[
  |Z_n(s)|\le
  \|V_n(s)\|_{L^2(K_\varphi)}
  \left(\sum_{j=1}^3\|\partial_j\varphi(s)\|_{L^2(K_\varphi)}^2\right)^{1/2}
\]
and Assumption~\ref{paper5:ass:selected-window-estimates} bounds the right-hand side
uniformly in \(n\) and \(s\). Furthermore,
\[
  \|Z_n-Z\|_{L^2(0,T;\mathbb R^3)}
  \le
  \|V_n-U\|_{L^2(0,T;L^2(K_\varphi))}
  \left(\sum_{j=1}^3\|\partial_j\varphi\|_{L^\infty(0,T;L^2(K_\varphi))}^2\right)^{1/2}
  \to0.
\]
Therefore
\[
  \int_0^T b_\infty\cdot Z_n(s)\,ds
  \to
  \int_0^T b_\infty\cdot Z(s)\,ds.
\]
Hence
\[
  -\int_0^T\!\int (\beta_n\cdot\nabla V_n)\cdot\varphi
  \,dy\,ds
  \to
  \int_0^T\!\int U\cdot(b_\infty\cdot\nabla\varphi)
  \,dy\,ds.
\]
Since \(b_\infty\) is constant in \(y\), integration by parts yields
\[
  \int_0^T\!\int U\cdot(b_\infty\cdot\nabla\varphi)
  \,dy\,ds
  =
  -\int_0^T\!\int (b_\infty\cdot\nabla U)\cdot\varphi
  \,dy\,ds.
\]

\emph{Step 5: assemble the limit.}
Passing to the limit in \eqref{paper5:eq:weak-n} in each term (Steps 4(i)--4(vii)) gives
\[
  \int_0^T\!\int
  \bigl(
  -U\cdot\partial_s\varphi
  -(U\otimes U):\nabla\varphi
  +\nu\nabla U:\nabla\varphi
  -\Pi\nabla\cdot\varphi
  -a_\infty(U+y\cdot\nabla U)\cdot\varphi
  -(b_\infty\cdot\nabla U)\cdot\varphi
  \bigr)
  \,dy\,ds=0.
\]
Since this holds for every compactly supported vector field \(\varphi\), equation~\eqref{paper5:eq:autonomous-limit} holds in distributions on
\((0,T)\times\mathbb R^3\).
\end{proof}

\subsection{Decomposition of alternatives of the scale-collapse branch}
\label{paper5:sec:branch-closure}

The cost argument in \cref{paper5:sec:cost,paper5:sec:finite-cost} excludes the finite-cost scale-collapse branch.  The remaining
scale-collapse obstruction is treated by a decomposition of alternatives, following the local
estimate and autonomous-reduction argument of the preceding sections.  The conclusion separates the excluded branches from the residual alternatives; it does not assert a Liouville theorem for every autonomous negative-drift equation.

\begin{theorem}[Selected core retention]\label{paper5:thm:core-retention}
Let a represented Type II branch be obtained from the selected-core construction of the preceding sections.  Then there exist \(R>0\), a cutoff \(\chi\in C_c^\infty(B_R)\), \(\chi\ge0\), a
constant \(\eta>0\), and a terminal time \(\tau_1\) such that, along every selected window used
below,
\[
  \int_{\mathbb R^3}|V(y,\tau)|^2\chi(y)\,dy\ge\eta
  \quad\text{for a.e. }\tau\ge\tau_1.
\]
In particular, every translated sequence \(V_n(y,s)=V(y,\tau_n+s)\) satisfies the localized
nontriviality condition in Assumption~\ref{paper5:ass:selected-window-estimates}, after discarding
finitely many \(n\).
\end{theorem}

\begin{proof}
The selected-core construction chooses a compact renormalized core carrying a fixed positive
amount of localized kinetic mass.  Since the cutoff \(\chi\) is supported strictly inside that
core and equals a positive spatial weight there, the selected-core lower bound gives the stated
estimate.  Translating \(\tau=\tau_n+s\) gives the final assertion on every finite selected
window.
\end{proof}

\begin{theorem}[Autonomous modulation selection]\label{paper5:thm:autonomous-modulation-selection}
Let a represented scale-collapse branch satisfy the modulation bounds obtained in the repaired
gauge construction.  Suppose that a sequence of windows \([\tau_n,\tau_n+T]\) is selected in the
autonomous regime.  Then, after passing to a subsequence, there exist constants
\(a_\infty\in\mathbb R\) and \(b_\infty\in\mathbb R^3\) such that
\[
  a(\tau_n+\cdot)\to a_\infty
  \quad\text{in }L^1(0,T),
  \qquad
  b(\tau_n+\cdot)\to b_\infty
  \quad\text{in }L^1(0,T;\mathbb R^3).
\]
\end{theorem}

\begin{proof}
The autonomous-window selection in the repaired-gauge modulation analysis gives convergence
of the modulation functions to constant autonomous parameters on the fixed window length
\(T\).  Passing to the corresponding subsequence gives the displayed
\(L^1\)-convergences.
\end{proof}

\begin{theorem}[Scale-collapse decomposition of alternatives]\label{paper5:thm:scale-collapse-stratification}
Let a represented Type II branch have genuine scale collapse and nonvanishing selected core.
Then exactly one of the following alternatives occurs after passing to selected terminal
windows:
\begin{enumerate}[label=(\roman*)]
  \item the branch has finite scale-collapsing cost, which is excluded by
  Theorem~\ref{paper5:thm:cost-exclusion};
  \item the branch has finite absolute cost, which is excluded by
  Corollary~\ref{paper5:cor:absolute-cost-exclusion};
  \item the branch admits thick autonomous windows in the sense of
  Definition~\ref{paper5:def:thick-autonomous-windows}, satisfies the selected-window estimates of
  Assumption~\ref{paper5:ass:selected-window-estimates}, and has autonomous modulation convergence;
  \item the branch has a thin-drift alternative: no selected fixed-length window carries a uniform
  positive lower bound for the averaged weighted negative drift;
  \item the branch has an autonomous-modulation alternative: thick windows exist, but no subsequence
  has a constant-coefficient \(L^1\) modulation limit in the sense needed for
  Theorem~\ref{paper5:thm:autonomous-limit};
  \item the branch has a compactness alternative: the local estimates needed for
  Assumption~\ref{paper5:ass:selected-window-estimates} fail on the selected windows.
\end{enumerate}
\end{theorem}

\begin{proof}
The scale-collapse selection theorem decomposes terminal scale-collapse windows according to the
averaged negative-drift mass
\[
  \frac1T\int_{\tau_n}^{\tau_n+T}a_-(\tau)M_R(\tau)\,d\tau.
\]
If the accumulated cost is finite, alternatives (i) or (ii) apply according to the cost
functional under consideration.  If the averaged negative-drift mass stays bounded below on a
subsequence, one obtains thick windows; otherwise the branch lies in the thin-drift alternative.
On thick windows, either the local estimates and core retention hold, in which case
Proposition~\ref{paper5:prop:selected-window-derived} and Theorem~\ref{paper5:thm:core-retention} give
Assumption~\ref{paper5:ass:selected-window-estimates}, or the branch lies in the compactness alternative.  If
the estimates hold but no subsequence has the required constant-coefficient modulation limit,
the branch lies in the autonomous-modulation alternative.  The remaining case is the thick
autonomous alternative.
\end{proof}

\begin{definition}[Stationary omega-limit class]\label{paper5:def:stationary-omega-class}
Let \((U,\Pi)\) be a nonzero autonomous limit produced by
Theorem~\ref{paper5:thm:autonomous-limit}.  We say that it has a stationary omega-limit in the
admissible class if there exists a nonzero pair
\[
  (W,Q),\qquad W\in L^3(\mathbb R^3),
\]
with the local energy and pressure reconstruction properties inherited from the selected
windows, such that
\[
  (W\cdot\nabla)W+\nabla Q
  =
  \nu\Delta W+a_\infty(W+y\cdot\nabla W)+b_\infty\cdot\nabla W,
  \qquad
  \nabla\cdot W=0,
\]
in \(\mathcal D'(\mathbb R^3)\).
\end{definition}

\begin{definition}[Rigidity-covered stationary class]\label{paper5:def:rigidity-covered-class}
For a parameter pair \((a_\infty,b_\infty)\), the stationary class is rigidity-covered if every
admissible stationary solution in Definition~\ref{paper5:def:stationary-omega-class} is either zero or
is not a terminal Type II branch.  In the classical backward self-similar normalization this is
the role of the Nečas--Růžička--Šverák rigidity theorem, after the harmless translation drift is
removed \cite{NecasRuzickaSverak1996}.  Other parameter regimes require their own stated
Liouville or classification theorem.
\end{definition}

\begin{theorem}[Stationary-rigidity exclusion]\label{paper5:thm:stationary-rigidity-blocker}
Assume a thick autonomous branch produces a nonzero autonomous limit by
Theorem~\ref{paper5:thm:autonomous-limit}.  If that autonomous limit has a stationary omega-limit in
the admissible class and the corresponding stationary class is rigidity-covered, then the branch
cannot occur as a terminal Type II scale-collapse branch.
\end{theorem}

\begin{proof}
The stationary omega-limit gives a nonzero stationary profile \((W,Q)\) in the admissible class.
The rigidity-covered hypothesis says that no such nonzero profile can represent a terminal Type
II branch.  This contradicts membership in the terminal scale-collapse Type II class.
\end{proof}

\begin{theorem}[Scale-collapse alternatives]\label{paper5:thm:scale-collapse-alternatives}
Let a represented Type II branch have genuine scale collapse and a nonvanishing selected core.
Then one of the following mutually exclusive outcomes occurs:
\begin{enumerate}[label=(\roman*)]
  \item the finite scale-collapsing cost exclusion applies;
  \item the finite absolute cost exclusion applies;
  \item the branch has a thin-drift alternative;
  \item the branch has an autonomous-modulation alternative;
  \item the branch has a compactness alternative;
  \item the branch has a nonzero autonomous reduced limit, and either stationary rigidity
  excludes it or the remaining obstruction is the corresponding stationary/generalized
  self-similar Liouville problem.
\end{enumerate}
\end{theorem}

\begin{proof}
Apply \cref{paper5:thm:scale-collapse-stratification}.  The finite-cost alternatives are
items (i) and (ii).  The thin-drift, modulation, and compactness alternatives are items (iii)--(v).  In the
remaining thick autonomous case, Lemma~\ref{paper5:lem:negative-parameter} gives \(a_\infty<0\), and
Theorem~\ref{paper5:thm:autonomous-limit} produces a nonzero autonomous reduced limit.  If
Definitions~\ref{paper5:def:stationary-omega-class} and~\ref{paper5:def:rigidity-covered-class} apply, then
Theorem~\ref{paper5:thm:stationary-rigidity-blocker} excludes the branch.  If not, the remaining
obstruction is the stationary or generalized self-similar Liouville problem for the
parameter pair \((a_\infty,b_\infty)\).
\end{proof}

\subsection{Conclusion}

The exclusions in Sections~\ref{paper5:sec:cost}--\ref{paper5:sec:finite-cost} are proved unconditionally within
the renormalized setting:
no branch with finite scale-collapsing (hence finite absolute) cost can have
\(\lambda(\tau)\to0\) together with a nonvanishing localized core.

Lemma~\ref{paper5:lem:negative-parameter} shows that any thick autonomous window admits a negative
autonomous drift limit \(a_\infty\).  Lemmas~\ref{paper5:lem:window-time-regularity}
and~\ref{paper5:lem:autonomous-window-compactness}, together with Theorem~\ref{paper5:thm:autonomous-limit},
assemble the selected-window estimates from the preceding sections into a nontrivial autonomous
finite-window limit \((U,\Pi)\) satisfying the full distributional autonomous equation.

Section~\ref{paper5:sec:branch-closure} gives the corresponding decomposition of alternatives for the scale-collapse branch.  The autonomous reduction gives a contradiction only in parameter regimes where a
stationary omega-limit exists and the resulting stationary class is covered by a Liouville
or rigidity theorem, such as the classical Nečas--Růžička--Šverák self-similar class.  Outside
those regimes the conclusion is the thin-drift, modulation, compactness, or
Liouville alternative, rather than an unconditional exclusion.

\clearpage
\section{Corrected Monotonicity and Scale-Drift Barriers}
\label{sec:corrected-monotonicity-scale-drift}

This section records the remaining conditional analytic inputs needed to use a
barrier argument on represented Type II branches.  The statements are written
entirely in the variables of the repaired gauge:
\[
  \partial_\tau V+(V\cdot\nabla)V+\nabla P
  =
  \nu\Delta V+a(\tau)(V+y\cdot\nabla V)+b(\tau)\cdot\nabla V,
  \qquad \nabla\cdot V=0.
\]
All pressure terms are understood modulo functions of time.

\subsection{Application Data for a Conditional Type II Barrier}

\begin{definition}[Barrier-application data]\label{paper6a:def:barrier-application-data}
A represented Type II branch is said to have barrier-application data if the
following analytic facts hold.
\begin{enumerate}[label=\textup{(\roman*)}]
\item the branch is a scale-supercritical local Type II branch in the sense of
the local concentration route;
\item the physical Navier--Stokes energy inequality gives bounded energy along
the branch;
\item the branch is represented by the repaired-gauge variables
\((V,P,a,b)\) above, including the scale, center, pressure normalization, and
modulation data;
\item the active localized barrier cost is identified with
\[
  \tilde{\mathfrak D}_{R_0}(\tau)
  =
  \nu\int_{\mathbb R^3}|\nabla V(y,\tau)|^2\phi_{R_0}(y)\,dy
  +a_+(\tau)\int_{\mathbb R^3}|V(y,\tau)|^2\phi_{R_0}(y)\,dy,
\]
or with an explicitly stated tail-comparable nonnegative replacement;
\item a localized energy identity or suitable-solution local energy inequality
is available for the cutoff used in the cost;
\item a corrected localized monotonicity inequality, in the sense of
\cref{paper6a:def:corrected-monotonicity-input}, is available;
\item the conditional Type II barrier principle used in the final assembly has
the same conclusion as the local Type II exclusion target: once the branch is
scale-supercritical and its admissible nonnegative barrier cost is divergent,
the branch cannot remain in the represented Type II class.
\end{enumerate}
\end{definition}

\begin{lemma}[Scale route and energy]\label{paper6a:lem:scale-energy-data}
If a branch is routed through the local Type II concentration alternative and
satisfies the physical Navier--Stokes energy inequality, then it satisfies
\cref{paper6a:def:barrier-application-data}\textup{(i)} and
\textup{(ii)}.
\end{lemma}

\begin{proof}
The local Type II route asserts precisely that the selected branch is
scale-supercritical and carries a concentration profile.  The physical energy
inequality gives the bounded-energy part.  These are precisely items
\textup{(i)} and \textup{(ii)}.
\end{proof}

\begin{lemma}[Cost identification]\label{paper6a:lem:cost-identification}
Assume the Type II barrier cost is the localized Navier--Stokes cost
\(\tilde{\mathfrak D}_{R_0}\).  Then divergence of
\[
  \int_{\tau_0}^{\infty}\tilde{\mathfrak D}_{R_0}(\tau)\,d\tau
\]
is the barrier-cost divergence used by the conditional Type II barrier
principle.
\end{lemma}

\begin{proof}
The active cost is assumed to be equal to
\(\tilde{\mathfrak D}_{R_0}\).  Thus the cost appearing in the local
Navier--Stokes estimates and the cost consumed by the barrier principle are
the same nonnegative functional.  The divergence condition is therefore
identical in the two formulations.
\end{proof}

\begin{lemma}[Parabolic representation class]\label{paper6a:lem:parabolic-class}
If the repaired-gauge representation theorem applies to a Type II branch, then
the branch lies in the parabolic transport-diffusion class required by the
conditional Type II barrier principle.
\end{lemma}

\begin{proof}
The repaired-gauge theorem supplies the renormalized orbit \((V,P,a,b)\), the
scale and center, the pressure reconstruction modulo functions of time, and the
modulation coefficients.  The displayed equation is a parabolic
Navier--Stokes equation with transport, dilation, and translation drift terms.
Consequently the branch has the parabolic represented structure required by
the barrier principle, with the same scale-supercritical and barrier-cost
interpretations.
\end{proof}

\begin{definition}[Corrected localized monotonicity input]
\label{paper6a:def:corrected-monotonicity-input}
A represented branch has corrected localized monotonicity if there are a
localized renormalized energy \(\mathcal E_{R_0}\), a nonnegative barrier cost
\(\mathcal D_{R_0}\), a constant \(c_0>0\), and a nonnegative remainder
\(\mathcal R_{R_0}\) such that, on the tail of the branch,
\[
  \frac{d}{d\tau}\mathcal E_{R_0}(\tau)
  +c_0\mathcal D_{R_0}(\tau)
  +\mathcal R_{R_0}(\tau)\le0
\]
in distributions.  If the localized energy identity has an integrable error
\(B_{R_0}\), the corrected energy is
\[
  \mathcal E_{R_0}^{\mathrm{corr}}(\tau)
  =
  \mathcal E_{R_0}(\tau)
  +\int_\tau^\infty B_{R_0}(s)\,ds,
\]
and the same inequality is required for
\(\mathcal E_{R_0}^{\mathrm{corr}}\).  Multiplication of a nonnegative barrier
cost by \(c_0>0\) does not change whether its tail integral diverges.
\end{definition}

\begin{lemma}[Energy and Caccioppoli bounds alone are not monotonicity]
\label{paper6a:lem:caccioppoli-not-monotonicity}
The global physical energy inequality and compact-cylinder Caccioppoli
estimates do not, by themselves, imply the corrected localized monotonicity
input of \cref{paper6a:def:corrected-monotonicity-input}.
\end{lemma}

\begin{proof}
The global physical energy inequality controls total physical energy and
physical dissipation, but it does not identify a localized renormalized energy
at the active Type II scale whose derivative is monotone after pressure flux,
cutoff flux, transport, and modulation terms are included.  The Caccioppoli
estimate gives local spacetime \(H^1\) control from the local energy inequality.
It is an a priori estimate, not a corrected monotonicity formula.  Therefore an
additional absorption or tail-correction argument is needed.
\end{proof}

\begin{theorem}[Conditional barrier application]\label{paper6a:thm:conditional-barrier-application}
Let a represented Type II branch have the data of
\cref{paper6a:def:barrier-application-data}.  If its admissible nonnegative
barrier cost has divergent tail integral, then the conditional Type II barrier
principle applies to the branch, and the branch is excluded from the represented
Type II class covered by that principle.
\end{theorem}

\begin{proof}
\Cref{paper6a:lem:scale-energy-data} gives the scale-supercritical and
bounded-energy inputs.  \Cref{paper6a:lem:cost-identification} identifies the
localized Navier--Stokes cost with the cost used in the barrier argument.
\Cref{paper6a:lem:parabolic-class} gives the represented parabolic class.
The corrected monotonicity input is an explicit hypothesis of
\cref{paper6a:def:barrier-application-data}.  These are precisely the analytic
inputs listed in that definition, and the last item of the definition is the
conditional Type II barrier implication with the same local conclusion.
\end{proof}

\subsection{Fixed-Cutoff Monotonicity}

Let \(\phi=\phi_{R_0}\in C_c^\infty(\mathbb R^3)\), \(0\le\phi\le1\), be the
cutoff used in the localized cost.  Define
\[
  \mathcal E_\phi(\tau)=\frac12\int_{\mathbb R^3}|V(y,\tau)|^2\phi(y)\,dy,
\]
\[
  A_\phi(\tau)=\int_{\mathbb R^3}|\nabla V(y,\tau)|^2\phi(y)\,dy,\qquad
  M_\phi(\tau)=\int_{\mathbb R^3}|V(y,\tau)|^2\phi(y)\,dy,
\]
and
\[
  \tilde{\mathfrak D}_{R_0}(\tau)=\nu A_\phi(\tau)+a_+(\tau)M_\phi(\tau),
  \qquad a_+(\tau)=\max\{a(\tau),0\}.
\]
The localized energy identity used below is
\begin{align}
  \frac{d}{d\tau}\mathcal E_\phi
  ={}&-\nu A_\phi
  +\frac{\nu}{2}\int |V|^2\Delta\phi
  +\frac12\int |V|^2V\cdot\nabla\phi
  +\int P\,V\cdot\nabla\phi \notag\\
  &-\frac{a}{2}M_\phi
  -\frac{a}{2}\int |V|^2y\cdot\nabla\phi
  -\frac12 b\cdot\int |V|^2\nabla\phi .
  \label{paper6a:eq:fixed-local-energy}
\end{align}

\begin{definition}[Fixed-cutoff error density]\label{paper6a:def:fixed-error-density}
Set \(a_-(\tau)=\max\{-a(\tau),0\}\) and define
\begin{align*}
  B_{R_0}(\tau)
  :={}&
  \frac{\nu}{2}\left|\int |V|^2\Delta\phi\right|
  +\frac12\left|\int |V|^2V\cdot\nabla\phi\right|
  +\left|\int P\,V\cdot\nabla\phi\right|\\
  &+\frac12|b(\tau)|\left|\int |V|^2\nabla\phi\right|
  +\frac12a_-(\tau)M_\phi(\tau)
  +\frac12|a(\tau)|\left|\int |V|^2y\cdot\nabla\phi\right|.
\end{align*}
The fixed-cutoff finite-error condition is
\[
  \int_{\tau_0}^{\infty}B_{R_0}(\tau)\,d\tau<\infty .
\]
\end{definition}

\begin{remark}[Pressure normalization]
The pressure term in \(B_{R_0}\) is independent of the representative of the
pressure modulo functions of time.  Replacing \(P\) by \(P+c(\tau)\) changes
the pressure flux by
\[
  c(\tau)\int V\cdot\nabla\phi
  =
  -c(\tau)\int \phi\,\nabla\cdot V
  =0.
\]
\end{remark}

\begin{theorem}[Fixed-cutoff corrected monotonicity]
\label{paper6a:thm:fixed-corrected-monotonicity}
Assume \eqref{paper6a:eq:fixed-local-energy} holds on
\((\tau_0,\infty)\) and that \(B_{R_0}\in L^1(\tau_0,\infty)\).  Define
\[
  \mathcal E_\phi^{\mathrm{corr}}(\tau)
  =
  \mathcal E_\phi(\tau)+\int_\tau^\infty B_{R_0}(s)\,ds .
\]
Then, in distributions on \((\tau_0,\infty)\),
\[
  \frac{d}{d\tau}\mathcal E_\phi^{\mathrm{corr}}(\tau)
  +\frac12\tilde{\mathfrak D}_{R_0}(\tau)\le0 .
\]
\end{theorem}

\begin{proof}
Add
\[
  \frac12\tilde{\mathfrak D}_{R_0}
  =
  \frac{\nu}{2}A_\phi+\frac12a_+M_\phi
\]
to both sides of \eqref{paper6a:eq:fixed-local-energy}.  The gradient term
becomes
\[
  -\nu A_\phi+\frac{\nu}{2}A_\phi=-\frac{\nu}{2}A_\phi\le0.
\]
The core scale term satisfies
\[
  -\frac{a}{2}M_\phi+\frac12a_+M_\phi
  =
  \frac12a_-M_\phi
  \le B_{R_0}(\tau).
\]
Every remaining term in \eqref{paper6a:eq:fixed-local-energy} is bounded above
by its corresponding absolute-value contribution in \(B_{R_0}\).  Therefore
\[
  \frac{d}{d\tau}\mathcal E_\phi(\tau)
  +\frac12\tilde{\mathfrak D}_{R_0}(\tau)
  \le B_{R_0}(\tau)
\]
in distributions.  Since \(B_{R_0}\in L^1(\tau_0,\infty)\), the correction
tail is finite for every \(\tau\ge\tau_0\), absolutely continuous on compact
subintervals, and
\[
  \frac{d}{d\tau}\int_\tau^\infty B_{R_0}(s)\,ds=-B_{R_0}(\tau).
\]
Adding this identity gives the claimed inequality.
\end{proof}

\begin{corollary}[Fixed-cutoff monotonicity input]\label{paper6a:cor:fixed-input}
If the repaired-gauge representation, the cost identification, the localized
energy identity, and the fixed-cutoff finite-error condition all hold, then
\cref{paper6a:def:corrected-monotonicity-input} holds with
\(\mathcal D_{R_0}=\tilde{\mathfrak D}_{R_0}\) and \(c_0=1/2\).
\end{corollary}

\begin{proof}
The repaired-gauge representation supplies \((V,P,a,b)\).  The cost
identification supplies \(\tilde{\mathfrak D}_{R_0}\).  The localized energy
identity and finite-error condition allow
\cref{paper6a:thm:fixed-corrected-monotonicity} to be applied.
\end{proof}

\begin{remark}[Finite-tail subconditions]\label{paper6a:rem:fixed-tail-subconditions}
The fixed-cutoff finite-error condition is equivalent to the conjunction of tail
integrability for the viscous cutoff term, the convective flux, the pressure
flux, the center-drift cutoff term, the negative-scale term, and the
scale-cutoff term:
\[
\begin{gathered}
  \left|\int |V|^2\Delta\phi\right|,\qquad
  \left|\int |V|^2V\cdot\nabla\phi\right|,\qquad
  \left|\int P\,V\cdot\nabla\phi\right|,\\
  |b|\left|\int |V|^2\nabla\phi\right|,\qquad
  a_-M_\phi,\qquad
  |a|\left|\int |V|^2y\cdot\nabla\phi\right|.
\end{gathered}
\]
Uniform compact-window bounds do not by themselves give this global tail
integrability on \([\tau_0,\infty)\).
\end{remark}

\begin{definition}[Direct fixed-cutoff application data]\label{paper6a:def:direct-fixed-data}
A represented Type II branch has direct fixed-cutoff application data if it has
the scale-supercritical concentration route, bounded physical energy,
repaired-gauge representation, cost identification with
\(\tilde{\mathfrak D}_{R_0}\), divergent admissible barrier cost, localized
energy identity, and the finite-tail bound
\(\int_{\tau_0}^{\infty}B_{R_0}<\infty\).
\end{definition}

\begin{theorem}[Direct fixed-cutoff application]\label{paper6a:thm:direct-fixed-application}
Every branch with direct fixed-cutoff application data is covered by
\cref{paper6a:thm:conditional-barrier-application}.
\end{theorem}

\begin{proof}
By \cref{paper6a:cor:fixed-input}, the finite-tail condition supplies the
corrected monotonicity input.  The remaining assumptions are precisely the other
items in \cref{paper6a:def:barrier-application-data}, together with cost
divergence.
\end{proof}

\begin{definition}[Compact-barrier replacement]\label{paper6a:def:compact-barrier-replacement}
A compact-barrier replacement is available when the compact Type II hypotheses
of the preceding sections imply
\[
  \int_{\tau_0}^{\infty}\tilde{\mathfrak D}_{R_0}(\tau)\,d\tau=\infty,
\]
and this divergence is the admissible barrier-cost divergence used in
\cref{paper6a:thm:conditional-barrier-application}.  In particular, it may be
supplied by the combination of classification of local alternatives, uniform
critical tightness, local windowed \(H^1\) control, and the cost identification
above whenever those are the active route to the compact Type II theorem.
\end{definition}

\begin{theorem}[Expanded compact-barrier application]\label{paper6a:thm:expanded-compact-application}
If a represented branch has all direct fixed-cutoff application data except
that cost divergence is supplied through the compact-barrier replacement of
\cref{paper6a:def:compact-barrier-replacement}, then the conditional Type II
barrier principle applies.
\end{theorem}

\begin{proof}
The compact-barrier replacement supplies the missing divergent admissible cost.
Substituting this into the direct fixed-cutoff data gives the hypotheses of
\cref{paper6a:thm:direct-fixed-application}.
\end{proof}

\subsection{Moving-Cutoff Monotonicity}

Let \(\phi\in C_c^\infty(\mathbb R^3)\) be radial, \(0\le\phi\le1\),
\(\phi=1\) on \(B_1\), and \(\phi=0\) outside \(B_2\).  For
\(R(\tau)\ge R_0\), locally absolutely continuous and nondecreasing, set
\[
  \Phi(\tau,y)=\phi(y/R(\tau)),\qquad
  A_{R(\tau)}=\{R(\tau)\le |y|\le2R(\tau)\},
\]
\[
  \mathcal E_R(\tau)=\frac12\int |V(y,\tau)|^2\Phi(\tau,y)\,dy,
\]
and
\[
  \tilde{\mathfrak D}_{R(\tau)}(\tau)
  =
  \nu\int|\nabla V|^2\Phi
  +a_+(\tau)\int|V|^2\Phi .
\]

\begin{definition}[Moving-cutoff error density]\label{paper6a:def:moving-error-density}
Define
\begin{align*}
  B_{\mathrm{mov}}(\tau)
  :={}&
  \frac12\left|\int |V|^2\partial_\tau\Phi\right|
  +\frac{\nu}{2}\left|\int |V|^2\Delta\Phi\right|
  +\frac12\left|\int |V|^2V\cdot\nabla\Phi\right|
  +\left|\int P\,V\cdot\nabla\Phi\right|\\
  &+\frac12|b(\tau)|\left|\int |V|^2\nabla\Phi\right|
  +\frac12a_-(\tau)\int |V|^2\Phi
  +\frac12|a(\tau)|\left|\int |V|^2y\cdot\nabla\Phi\right|.
\end{align*}
The moving finite-error condition is
\[
  \int_{\tau_0}^{\infty}B_{\mathrm{mov}}(\tau)\,d\tau<\infty .
\]
\end{definition}

\begin{theorem}[Moving-cutoff corrected monotonicity]
\label{paper6a:thm:moving-corrected-monotonicity}
Assume the localized energy identity is valid with the time-dependent cutoff
\(\Phi\), and assume \(B_{\mathrm{mov}}\in L^1(\tau_0,\infty)\).  Then
\[
  \mathcal E_R^{\mathrm{corr}}(\tau)
  =
  \mathcal E_R(\tau)+\int_\tau^\infty B_{\mathrm{mov}}(s)\,ds
\]
satisfies
\[
  \frac{d}{d\tau}\mathcal E_R^{\mathrm{corr}}(\tau)
  +\frac12\tilde{\mathfrak D}_{R(\tau)}(\tau)\le0
\]
in distributions.
\end{theorem}

\begin{proof}
Use the localized energy identity with the time-dependent test function
\(\Phi(\tau,y)\).  Compared with the fixed-cutoff identity, the only new term
is
\[
  \frac12\int |V|^2\partial_\tau\Phi,
\]
which comes from differentiating the cutoff inside \(\mathcal E_R\).  All
other terms are the fixed-cutoff terms with \(\phi\) replaced by \(\Phi\).  Add
\(\frac12\tilde{\mathfrak D}_{R(\tau)}\) to both sides.  The gradient term
leaves
\[
  -\frac{\nu}{2}\int|\nabla V|^2\Phi\le0,
\]
and the core scale term contributes
\[
  \frac12a_-\int |V|^2\Phi .
\]
Each remaining term is bounded above by the corresponding absolute-value term
in \(B_{\mathrm{mov}}\).  Hence
\[
  \frac{d}{d\tau}\mathcal E_R(\tau)
  +\frac12\tilde{\mathfrak D}_{R(\tau)}(\tau)
  \le B_{\mathrm{mov}}(\tau).
\]
The moving finite-error condition makes the correction tail finite and
absolutely continuous on compact subintervals.  Differentiating that tail
subtracts \(B_{\mathrm{mov}}\), giving the stated inequality.
\end{proof}

\begin{definition}[Summable moving-annulus errors]\label{paper6a:def:summable-moving-errors}
The moving-annulus summability condition is the existence of a nondecreasing
locally absolutely continuous \(R(\tau)\to\infty\) such that each of the seven
quantities
\[
\begin{gathered}
  \left|\int |V|^2\partial_\tau\Phi\right|,\quad
  \left|\int |V|^2\Delta\Phi\right|,\quad
  \left|\int |V|^2V\cdot\nabla\Phi\right|,\quad
  \left|\int P\,V\cdot\nabla\Phi\right|,\\
  |b|\left|\int |V|^2\nabla\Phi\right|,\quad
  a_-\int |V|^2\Phi,\quad
  |a|\left|\int |V|^2y\cdot\nabla\Phi\right|
\end{gathered}
\]
is integrable on \([\tau_0,\infty)\).
\end{definition}

\begin{lemma}[Summable annuli imply moving finite error]
\label{paper6a:lem:summable-moving-errors}
If \cref{paper6a:def:summable-moving-errors} holds, then
\(B_{\mathrm{mov}}\in L^1(\tau_0,\infty)\).
\end{lemma}

\begin{proof}
The density \(B_{\mathrm{mov}}\) is a finite positive linear combination of
the seven quantities listed in
\cref{paper6a:def:summable-moving-errors}.  If each is integrable, then their
linear combination is integrable.
\end{proof}

\begin{definition}[Summable tightness schedule]\label{paper6a:def:summable-tightness}
A summable tightness schedule is a nondecreasing sequence \(R_n\to\infty\) such
that, on unit windows \(I_n=[\tau_0+n,\tau_0+n+1]\),
\[
  \sum_{n=0}^{\infty}
  \sup_{\tau\in I_n}
  \|V(\tau)\|_{L^3(A_{R_n})}^3<\infty,
\]
and the same schedule is compatible with the pressure-tail normalization used
for \(P\).  This is stronger than uniform critical tightness, which permits
large radii on each window but does not by itself give a monotone radius
schedule with a summable series.
\end{definition}

\begin{definition}[Moving-cost admissibility]\label{paper6a:def:moving-cost-admissibility}
The moving localized cost is admissible if
\[
  \tilde{\mathfrak D}_{R(\tau)}(\tau)
  =
  \nu\int|\nabla V|^2\phi_{R(\tau)}
  +a_+(\tau)\int|V|^2\phi_{R(\tau)}
\]
is accepted as the active Type II barrier cost, or is compared to the fixed
cost by a positive constant on the tail so that cost divergence is unchanged.
\end{definition}

\begin{theorem}[Moving-cutoff replacement]\label{paper6a:thm:moving-replacement}
Assume the repaired-gauge representation, localized energy identity,
moving-cost admissibility, and summable moving-annulus errors.  Then the moving
cutoff supplies the corrected monotonicity input needed in
\cref{paper6a:thm:conditional-barrier-application}.
\end{theorem}

\begin{proof}
\Cref{paper6a:lem:summable-moving-errors} gives
\(B_{\mathrm{mov}}\in L^1(\tau_0,\infty)\).  Applying
\cref{paper6a:thm:moving-corrected-monotonicity} gives corrected monotonicity
with cost \(\tilde{\mathfrak D}_{R(\tau)}\).  The moving-cost admissibility
assumption identifies this cost, or a tail-comparable version of it, with the
active barrier cost.
\end{proof}

\subsection{Negative Scale Drift}

Set
\[
  M_R(\tau)=\int |V(y,\tau)|^2\Phi(\tau,y)\,dy .
\]
The negative-scale contribution in the moving cutoff argument is
\[
  \int_{\tau_0}^{\infty}a_-(\tau)M_R(\tau)\,d\tau.
\]

\begin{definition}[Negative scale-drift alternatives]\label{paper6a:def:negative-drift}
We say that the branch has finite negative drift if
\[
  \int_{\tau_0}^{\infty}a_-(\tau)\,d\tau<\infty.
\]
It has bounded moving \(L^2\) core mass if
\[
  \sup_{\tau\ge\tau_0}M_R(\tau)<\infty.
\]
It has a moving \(L^2\) core floor if there are \(m_0>0\) and
\(\tau_1\ge\tau_0\) such that
\[
  M_R(\tau)\ge m_0
  \qquad\text{for a.e. }\tau\ge\tau_1.
\]
The scale-collapse drift obstruction is the divergent alternative
\[
  \int_{\tau_0}^{\infty}a_-(\tau)M_R(\tau)\,d\tau=\infty .
\]
\end{definition}

\begin{lemma}[Finite negative drift controls the scale term]
\label{paper6a:lem:finite-negative-drift}
If the branch has bounded moving \(L^2\) core mass and finite negative drift,
then
\[
  \int_{\tau_0}^{\infty}a_-(\tau)M_R(\tau)\,d\tau<\infty .
\]
\end{lemma}

\begin{proof}
Let \(M_*=\sup_{\tau\ge\tau_0}M_R(\tau)<\infty\).  Since \(a_-\ge0\),
\[
  \int_{\tau_0}^{\infty}a_-(\tau)M_R(\tau)\,d\tau
  \le
  M_*\int_{\tau_0}^{\infty}a_-(\tau)\,d\tau<\infty.
\]
\end{proof}

\begin{theorem}[Finite-or-infinite scale-negative alternative]
\label{paper6a:thm:finite-infinite-scale-negative}
For every represented orbit for which \(a_-\) and \(M_R\) are measurable and
nonnegative, precisely one of the following alternatives holds:
\[
  \int_{\tau_0}^{\infty}a_-(\tau)M_R(\tau)\,d\tau<\infty,
  \qquad
  \int_{\tau_0}^{\infty}a_-(\tau)M_R(\tau)\,d\tau=\infty.
\]
\end{theorem}

\begin{proof}
The integral
\[
  I_R=\int_{\tau_0}^{\infty}a_-(\tau)M_R(\tau)\,d\tau
\]
is an extended nonnegative number.  Hence either \(I_R<\infty\) or
\(I_R=\infty\), and the two cases are mutually exclusive.
\end{proof}

\begin{lemma}[Logarithmic scale identity]\label{paper6a:lem:log-scale}
For the repaired-gauge convention
\[
  \tau_t=\lambda(t)^{-2},\qquad a(\tau)=\lambda(t)\lambda_t(t),
\]
one has
\[
  \int_{\tau_1}^{\tau_2}a(\tau)\,d\tau
  =
  \log\lambda(\tau_2)-\log\lambda(\tau_1)
\]
whenever \(\tau_0\le\tau_1<\tau_2<\infty\).
\end{lemma}

\begin{proof}
The chain rule gives
\[
  \frac{d}{d\tau}\log\lambda
  =
  \frac{\lambda_t}{\lambda}\frac{dt}{d\tau}
  =
  \frac{\lambda_t}{\lambda}\lambda^2
  =
  \lambda\lambda_t=a.
\]
Integrating over \([\tau_1,\tau_2]\) proves the identity.
\end{proof}

\begin{theorem}[Collapse with core floor forces scale-drift obstruction]
\label{paper6a:thm:collapse-core-floor}
Assume
\[
  \lambda(\tau)\to0\qquad\text{as }\tau\to\infty,
\]
and assume the moving \(L^2\) core floor of
\cref{paper6a:def:negative-drift}.  Then
\[
  \int_{\tau_0}^{\infty}a_-(\tau)M_R(\tau)\,d\tau=\infty .
\]
\end{theorem}

\begin{proof}
By \cref{paper6a:lem:log-scale},
\[
  \int_{\tau_1}^{\tau_2}a(\tau)\,d\tau
  =
  \log\lambda(\tau_2)-\log\lambda(\tau_1).
\]
Since \(\lambda(\tau)\to0\), the right-hand side tends to \(-\infty\) as
\(\tau_2\to\infty\).  If
\(\int_{\tau_1}^{\infty}a_-(\tau)\,d\tau<\infty\), then
\[
  \int_{\tau_1}^{\tau_2}a(\tau)\,d\tau
  =
  \int_{\tau_1}^{\tau_2}a_+(\tau)\,d\tau
  -
  \int_{\tau_1}^{\tau_2}a_-(\tau)\,d\tau
  \ge
  -\int_{\tau_1}^{\infty}a_-(\tau)\,d\tau>-\infty,
\]
which contradicts the preceding limit.  Thus
\(\int_{\tau_1}^{\infty}a_-(\tau)\,d\tau=\infty\).

By the moving \(L^2\) core floor, \(M_R(\tau)\ge m_0\) for a.e.
\(\tau\ge\tau_1\).  Therefore
\[
  \int_{\tau_0}^{\infty}a_-(\tau)M_R(\tau)\,d\tau
  \ge
  m_0\int_{\tau_1}^{\infty}a_-(\tau)\,d\tau
  =
  \infty.
\]
\end{proof}

\begin{corollary}[Scale-drift routing]\label{paper6a:cor:scale-drift-routing}
In the moving-cutoff setup, the negative-scale term has the following
exhaustive routing.  If bounded moving \(L^2\) core mass and finite negative
drift hold, then the negative-scale term is integrable.  If the weighted
negative-drift integral is infinite, the branch lies in the scale-collapse
drift obstruction.  If genuine collapse and a moving \(L^2\) core floor also
hold, then the obstruction is forced.
\end{corollary}

\begin{proof}
The first case is \cref{paper6a:lem:finite-negative-drift}.  The exhaustive
split is \cref{paper6a:thm:finite-infinite-scale-negative}.  The final
assertion is \cref{paper6a:thm:collapse-core-floor}.
\end{proof}

\subsection{Scale-Drift Barrier Costs}

\begin{definition}[Scale-collapse drift cost]\label{paper6a:def:scale-collapse-cost}
For a represented branch define
\[
  \mathfrak C_{\mathrm{sc}}(\tau)
  =
  a_-(\tau)M_R(\tau),
  \qquad
  M_R(\tau)=\int |V(y,\tau)|^2\phi_{R(\tau)}(y)\,dy,
\]
and
\[
  \mathcal C_{\mathrm{sc}}[\tau_0,\infty)
  =
  \int_{\tau_0}^{\infty}\mathfrak C_{\mathrm{sc}}(\tau)\,d\tau .
\]
Thus the scale-collapse drift obstruction is precisely
\(\mathcal C_{\mathrm{sc}}[\tau_0,\infty)=\infty\).  The density is
nonnegative and local in the renormalized core.
\end{definition}

\begin{definition}[Scale-collapse cost admissibility]\label{paper6a:def:scale-collapse-cost-admissibility}
The scale-collapse drift cost is admissible when:
\begin{enumerate}[label=\textup{(\roman*)}]
\item \(\mathfrak C_{\mathrm{sc}}\) is measurable, nonnegative, and locally
integrable on finite \(\tau\)-intervals for represented branches;
\item divergence of \(\mathcal C_{\mathrm{sc}}[\tau_0,\infty)\) is an
accepted Type II barrier condition;
\item the resulting barrier conclusion is the same local Type II exclusion
conclusion used for the fixed localized cost.
\end{enumerate}
\end{definition}

\begin{theorem}[Scale-collapse drift is a barrier when admissible]
\label{paper6a:thm:scale-collapse-drift-barrier}
If the scale-collapse drift obstruction holds and the cost
\(\mathfrak C_{\mathrm{sc}}\) is admissible, then the divergent-cost hypothesis
of \cref{paper6a:thm:conditional-barrier-application} is satisfied for the
scale-collapse branch.
\end{theorem}

\begin{proof}
By \cref{paper6a:def:scale-collapse-cost}, the obstruction is precisely
\[
  \mathcal C_{\mathrm{sc}}[\tau_0,\infty)=\infty .
\]
By \cref{paper6a:def:scale-collapse-cost-admissibility}, this divergence is an
accepted nonnegative Type II barrier condition with the same local conclusion.
\end{proof}

\begin{corollary}[Scale-collapse suppression under the conditional principle]
\label{paper6a:cor:scale-collapse-suppression}
Assume a scale-supercritical represented Type II branch has scale-collapse
drift obstruction, admissible scale-collapse drift cost, and all other
application data required by the conditional Type II barrier principle.  Then
the branch is excluded from the represented Type II class covered by that
principle.
\end{corollary}

\begin{proof}
\Cref{paper6a:thm:scale-collapse-drift-barrier} supplies the divergent
admissible barrier cost.  Applying
\cref{paper6a:thm:conditional-barrier-application} gives the exclusion.
\end{proof}

\begin{definition}[Absolute scale-drift cost]\label{paper6a:def:absolute-scale-cost}
Define
\[
  \mathfrak C_{\mathrm{abs}}(\tau)
  =
  \nu\int|\nabla V|^2\phi_{R(\tau)}
  +|a(\tau)|M_R(\tau).
\]
The absolute scale-drift cost is admissible if divergence of
\(\int_{\tau_0}^{\infty}\mathfrak C_{\mathrm{abs}}(\tau)\,d\tau\) is accepted
as a Type II barrier condition with the same local conclusion.
\end{definition}

\begin{lemma}[Scale-collapse drift forces absolute-cost divergence]
\label{paper6a:lem:scale-drift-absolute}
If the scale-collapse drift obstruction holds, then
\[
  \int_{\tau_0}^{\infty}\mathfrak C_{\mathrm{abs}}(\tau)\,d\tau=\infty .
\]
\end{lemma}

\begin{proof}
Since
\[
  \mathfrak C_{\mathrm{abs}}(\tau)
  =
  \nu\int|\nabla V|^2\phi_{R(\tau)}
  +|a(\tau)|M_R(\tau)
  \ge a_-(\tau)M_R(\tau),
\]
divergence of \(\int a_-M_R\) forces divergence of
\(\int\mathfrak C_{\mathrm{abs}}\).
\end{proof}

\begin{corollary}[Absolute scale-cost barrier]\label{paper6a:cor:absolute-scale-cost-barrier}
If the scale-collapse drift obstruction holds and the absolute scale-drift
cost is admissible, then the divergent-cost hypothesis of
\cref{paper6a:thm:conditional-barrier-application} is satisfied.
\end{corollary}

\begin{proof}
\Cref{paper6a:lem:scale-drift-absolute} gives divergence of the absolute
scale-drift cost, and admissibility makes this divergence an accepted Type II
barrier condition.
\end{proof}

\begin{theorem}[Moving route with scale-collapse fallback]
\label{paper6a:thm:moving-route-scale-fallback}
Assume the moving-cutoff setup and a scale-supercritical represented Type II
branch.
\begin{enumerate}[label=\textup{(\roman*)}]
\item If the negative-scale term is integrable and the remaining moving-annulus
errors are summable, then the moving-cutoff monotonicity route is available.
\item If the scale-collapse drift obstruction holds and either the
scale-collapse drift cost or the absolute scale-drift cost is admissible, then
the divergent-cost hypothesis of the conditional barrier principle is
satisfied directly.
\item If, in addition, the remaining application data of
\cref{paper6a:def:barrier-application-data} hold, then the represented Type II
branch is excluded.
\end{enumerate}
\end{theorem}

\begin{proof}
The first branch is \cref{paper6a:thm:moving-replacement}.  The second branch
is \cref{paper6a:thm:scale-collapse-drift-barrier} or
\cref{paper6a:cor:absolute-scale-cost-barrier}, depending on which admissible
cost is used.  The third branch is
\cref{paper6a:thm:conditional-barrier-application}.
\end{proof}

\subsection{Terminal Conditional Assembly}

\begin{definition}[Abstract final Type II data]\label{paper6a:def:abstract-final-data}
The abstract final Type II data consist of:
\begin{enumerate}[label=\textup{(\roman*)}]
\item classification of represented Type II branches into the local alternatives
used in the final assembly;
\item exclusion of the radiative or noncompact alternative, equivalently the
uniform critical tightness conclusion on represented branches;
\item exclusion of the rough-core alternative, equivalently uniform local
windowed \(H^1\) control on represented branches;
\item the conditional Type II barrier principle of
\cref{paper6a:thm:conditional-barrier-application}.
\end{enumerate}
\end{definition}

\begin{theorem}[Abstract final conditional Type II theorem]
\label{paper6a:thm:abstract-final-typeII}
If the abstract final Type II data hold, then every represented Type II branch
in the covered class is excluded.
\end{theorem}

\begin{proof}
The classification data place every represented Type II branch in one of the
local alternatives.  The radiative and rough-core alternatives are excluded by
items \textup{(ii)} and \textup{(iii)} of
\cref{paper6a:def:abstract-final-data}.  Any branch carrying a divergent
admissible Type II barrier cost is excluded by item \textup{(iv)}.  The
remaining compact, multibubble, cascade, and scale-rigid alternatives are the
ones treated in the preceding local sections and in the final assembly below.
\end{proof}

\begin{definition}[Expanded terminal data]\label{paper6a:def:expanded-terminal-data}
Expanded terminal data are the following analytic inputs:
\begin{enumerate}[label=\textup{(\roman*)}]
\item classification of represented Type II branches into the ordered local
alternatives;
\item the aggregate technical Type II exclusions used for compact,
multibubble, scale-collapse, and rough-core branches, excluding the conditional
barrier principle itself;
\item the no-radiation route yielding uniform critical tightness;
\item an exhaustive terminal state-space partition and a local finite
critical-mass package for retained active terminal strata;
\item small-data stability in \(L^3\) and a critical Navier--Stokes profile
decomposition with nonlinear stability;
\item removal of the scattering branch and discard of physically exterior
regular profiles;
\item the repaired-gauge representation theorem;
\item the Caccioppoli/windowed regularity input;
\item the scale-collapse rigidity input used by the multibubble and
scale-collapse reductions;
\item the conditional Type II barrier principle.
\end{enumerate}
\end{definition}

\begin{lemma}[Terminal data give the radiative exclusion]
\label{paper6a:lem:terminal-radiative}
The expanded terminal data imply exclusion of the radiative or noncompact
alternative.
\end{lemma}

\begin{proof}
The repaired-gauge representation turns represented orbits into terminal
sequences.  The terminal state-space partition, the finite critical-mass input,
small-data \(L^3\) stability, and the critical profile decomposition give the
profile completeness needed for the terminal profile route.  The scattering
branch and exterior regular profiles are removed by hypothesis.  Together with
the aggregate technical Type II exclusions, the Caccioppoli input, and the
scale-collapse rigidity input, the multibubble alternatives are removed.  The
remaining non-bubble noncompact branch is the radiative alternative,
and the no-radiation route excludes it, equivalently giving uniform critical
tightness.
\end{proof}

\begin{lemma}[Terminal data give the rough-core exclusion]
\label{paper6a:lem:terminal-roughcore}
The expanded terminal data imply exclusion of the rough-core alternative.
\end{lemma}

\begin{proof}
The rough-core composition requires the represented branch, local pressure and
modulation control, compact-cylinder critical bounds, and Caccioppoli/windowed
regularity.  These are included in the expanded terminal data through the
aggregate technical exclusions, the repaired-gauge representation, and the
Caccioppoli input.  Hence the universal rough-core exclusion is available.
\end{proof}

\begin{lemma}[Expanded terminal data give abstract final data]
\label{paper6a:lem:expanded-to-abstract}
The expanded terminal data imply the abstract final Type II data.
\end{lemma}

\begin{proof}
The expanded terminal data include classification of represented Type II
branches and the conditional Type II barrier principle.  By
\cref{paper6a:lem:terminal-radiative} they give the radiative exclusion, and by
\cref{paper6a:lem:terminal-roughcore} they give the rough-core exclusion.
These are precisely the four entries in
\cref{paper6a:def:abstract-final-data}.
\end{proof}

\begin{theorem}[Expanded terminal Type II theorem]
\label{paper6a:thm:expanded-terminal-typeII}
If the expanded terminal data hold, then every represented Type II branch in
the covered class is excluded.
\end{theorem}

\begin{proof}
\Cref{paper6a:lem:expanded-to-abstract} gives the abstract final Type II data.
Applying \cref{paper6a:thm:abstract-final-typeII} excludes every represented
Type II branch in the covered class.
\end{proof}

\begin{remark}[Direct branchwise reading]
Equivalently, compact single-core branches are removed by the compact Type II
barrier; scale-collapse branches are removed by the scale-collapse rigidity
input, or by \cref{paper6a:thm:scale-collapse-drift-barrier} when the
scale-collapse drift cost is admissible; multibubble and radiative branches
are removed by the terminal profile and no-radiation inputs; rough-core
branches are removed by the rough-core estimate; and every divergent admissible
barrier cost is converted to the same local Type II exclusion conclusion by
the conditional barrier principle.
\end{remark}

\clearpage
\section{Assembly of the Local Type II Exclusion}\label{sec:typeII-assembly}
\subsection{Local Entry}

Let \((u,p)\) be a suitable weak solution of the three-dimensional
Navier--Stokes equations
\[
  \partial_tu+(u\cdot\nabla)u+\nabla p=\nu\Delta u,
  \qquad \nabla\cdot u=0,
\]
on \(\mathbb R^3\times(T-\delta,T)\).  For \(z_0=(x_0,T)\) and \(r>0\), set
\[
  Q_r(z_0)=B_r(x_0)\times(T-r^2,T),
\]
\[
  C(z_0,r)=r^{-2}\int_{Q_r(z_0)}|u|^3\,dx\,dt,
\]
and
\[
  D(z_0,r)=r^{-2}\int_{T-r^2}^{T}\int_{B_r(x_0)}
  |p(x,t)-(p)_{B_r(x_0)}(t)|^{3/2}\,dx\,dt.
\]
The pressure is always taken modulo functions of time, and the spatial mean
normalization is the one used in the Caffarelli--Kohn--Nirenberg criterion.

\begin{definition}[Singular set and local Type II point]
Let \(\Sigma(T)\) be the set of points \(x_0\) for which \(u\) is not locally
bounded in any backward cylinder \(Q_r(x_0,T)\).  A point \(x_0\in\Sigma(T)\)
is in the local Type II alternative if it carries a positive local
Caffarelli--Kohn--Nirenberg concentration sequence and no local Type I
scale-invariant bound holds at \((x_0,T)\).
\end{definition}

\begin{definition}[Positive local Type II concentration sequence]
\label{paper6:def:local-typeII-sequence}
A sequence \(r_n\downarrow0\) is a positive local Type II concentration sequence
at \((x_0,T)\) if there exists \(\eta>0\) such that
\[
  C((x_0,T),r_n)+D((x_0,T),r_n)\ge \eta
  \qquad\text{for all }n,
\]
and \((x_0,T)\) is in the local Type II alternative.  The associated parabolic
rescalings are
\[
  u_n(y,s)=r_nu(x_0+r_ny,T+r_n^2s),
  \qquad
  p_n(y,s)=r_n^2p(x_0+r_ny,T+r_n^2s).
\]
\end{definition}

\begin{theorem}[Local concentration and Type I/Type II dichotomy]
\label{paper6:thm:paper0}
If \(\Sigma(T)\ne\emptyset\), then there are
\(x_0\in\Sigma(T)\), \(\eta>0\), and \(r_n\downarrow0\) such that
\[
  C((x_0,T),r_n)+D((x_0,T),r_n)\ge\eta
  \qquad\text{for all }n.
\]
At the concentration point exactly one of the following alternatives holds:
there is a local Type I scale-invariant bound, or \((x_0,T)\) is in the local
Type II alternative.
\end{theorem}

\begin{proof}
The theorem is the standing local concentration dichotomy assumed in the assembly.
\end{proof}

\subsection{Local Theorems Used in the Assembly}

The local results needed for the final assembly are restated in the following form.

\begin{theorem}[Repaired-gauge representation]
\label{paper6:thm:representation}
Let a positive local Type II concentration sequence contain a nondegenerate
selected core.  Then, after passing to selected windows, there are local
scale-center variables \(\lambda(\tau)>0\), \(x_c(\tau)\), renormalized fields
\[
  V(y,\tau)=\lambda(\tau)u(x_c(\tau)+\lambda(\tau)y,t(\tau)),
  \qquad
  P(y,\tau)=\lambda(\tau)^2p(x_c(\tau)+\lambda(\tau)y,t(\tau))+\pi(\tau),
\]
and modulation coefficients \(a,b\) such that
\[
  \partial_\tau V+(V\cdot\nabla)V+\nabla P
  =
  \nu\Delta V+a(\tau)(V+y\cdot\nabla V)+b(\tau)\cdot\nabla V,
  \qquad \nabla\cdot V=0.
\]
On each compact renormalized window one has the local suitable energy
inequality, local pressure reconstruction modulo functions of time, stability
under subsequences, and the compact-window bounds proved in \cref{sec:repaired-gauge-construction}.  Failure of the nondegenerate scale-center selection is part of the
multibubble, cascade, or scale-rigidity alternatives.
\end{theorem}

\begin{proof}
\Cref{paper2:thm:C} is precisely the compact-window representation theorem.
\end{proof}

\begin{theorem}[Compact single-core exclusion]
\label{paper6:thm:single-core}
No compact single-core vanishing-dissipation branch can be a local Type II
sequence.  More explicitly, let a represented selected-window sequence have:
\begin{enumerate}[label=\textup{(\roman*)}]
\item positive local CKN concentration on a fixed compact cylinder;
\item uniformly bounded compact-cylinder suitable estimates \(A+E+C+D\);
\item nondegenerate local scale-center representation;
\item bounded modulation coefficients;
\item local pressure reconstruction and pressure stability modulo spatial
means;
\item vanishing localized dissipation on some retained compact cylinder.
\end{enumerate}
Then the sequence is incompatible with the local Type II condition.
\end{theorem}

\begin{proof}
\Cref{paper1:thm:single-core-criterion} gives the stated exclusion.
\end{proof}

\begin{theorem}[Multibubble and cascade exclusion]
\label{paper6:thm:multibubble}
Assume the terminal admissibility and scale-rigidity hypotheses stated in the
multibubble/cascade analysis.
Then no local multibubble, cascade, or gauge-degenerate Type II branch can occur
in the class produced by the local active-core extraction.  Equivalently, every
failure of the single nondegenerate compact-core alternative is excluded by the
local multibubble/cascade analysis under those stated hypotheses.
\end{theorem}

\begin{proof}
\Cref{paper3:thm:multi}, together with its removal corollary
\cref{paper3:cor:multi-removal}, gives the stated exclusion.
\end{proof}

\begin{theorem}[Rough-core reduction]
\label{paper6:thm:rough-core}
On any compact renormalized window satisfying the local energy inequality and
bounded modulation hypotheses, finite outer velocity-pressure CKN control
\(\mathcal C+\mathcal D\) implies finite inner windowed \(H^1\) control.  Hence
failure of local windowed \(H^1\) control on a retained inner cylinder forces
failure of the compact-cylinder \(\mathcal C+\mathcal D\) bounds on every
relevant enclosing cylinder.
\end{theorem}

\begin{proof}
\Cref{paper4:cor:rough-core} gives the stated reduction.
\end{proof}

\begin{theorem}[Finite-cost scale-collapse exclusion]
\label{paper6:thm:scale-cost}
No represented branch with finite scale-collapsing cost, and hence no branch
with finite absolute scale cost, can have genuine scale collapse together with a
nonvanishing selected core.
\end{theorem}

\begin{proof}
\Cref{paper5:thm:cost-exclusion,paper5:cor:absolute-cost-exclusion} give the stated exclusion.
\end{proof}

\begin{theorem}[Scale-collapse decomposition of alternatives]
\label{paper6:thm:scale-stratification}
Let a represented Type II branch have genuine scale collapse and a nonvanishing
selected core.  After passing to selected terminal windows, exactly one of the
alternatives isolated in \cref{sec:scale-collapse-cost} occurs:
\begin{enumerate}[label=\textup{(\roman*)}]
\item finite scale-collapsing cost;
\item finite absolute scale cost;
\item thin-drift alternative;
\item autonomous-modulation alternative;
\item compactness alternative;
\item a nonzero autonomous reduced limit in the scale-rigid state.
\end{enumerate}
In the local Type II class of alternatives used by the multibubble/cascade and
scale-collapse analyses, alternatives
\textup{(iii)}--\textup{(vi)} do not constitute additional compact single-core
mechanisms: the thin-drift, modulation, and compactness alternatives lie outside
the compact finite-cost scale-collapse branch, and the scale-rigid state is
covered by the scale-rigidity hypothesis in the multibubble/cascade theorem.
\end{theorem}

\begin{proof}
The assertion follows by combining the scale-collapse decomposition of
alternatives with the scale-rigidity alternative used in the
multibubble/cascade reduction.
\end{proof}

\begin{definition}[Virial-dissipative terminal state]
\label{paper6:def:virial-dissipative}
Let \(G_\nu(y)=e^{-|y|^2/(4\nu)}\).  A scale-rigid terminal state is called
velocity-virial dissipative if its selected terminal representative
\((V,P)\) satisfies, after the usual pressure normalization, the autonomous
self-similar equation
\[
  \partial_\tau V+(V\cdot\nabla)V+\nabla P
  =
  \nu\Delta V-\frac12(V+y\cdot\nabla V),
  \qquad \nabla\cdot V=0,
\]
and the Gaussian identity below is justified by cutoff approximation on the
selected renormalized windows.  It is called swirl-virial dissipative if, in
addition, the representative is axisymmetric, \(\Gamma=R V_\theta\) is the
scale-invariant circulation, and the corresponding weighted circulation
identity below is justified.
\end{definition}

\begin{theorem}[Virial exclusion for dissipative terminal states]
\label{paper6:thm:virial-exclusion}
Let a scale-rigid terminal state satisfy the hypotheses of
\cref{paper6:def:virial-dissipative} and be bounded for all negative renormalized
times in the sense needed for the Gaussian integrals below.

\begin{enumerate}[label=\textup{(\alph*)}]
\item Suppose the velocity virial defect is absorbable: for some
\(\kappa>0\) and a.e. \(\tau\),
\[
  -\frac{1}{2\nu}\int_{\mathbb R^3}|V|^2(V\cdot y)G_\nu\,dy
  -\frac1\nu\int_{\mathbb R^3}P(V\cdot y)G_\nu\,dy
  \le
  (1-\kappa)\int_{\mathbb R^3}|V|^2G_\nu\,dy
  +2\nu(1-\kappa)\int_{\mathbb R^3}|\nabla V|^2G_\nu\,dy .
\]
Then
\[
  \int_{\mathbb R^3}|V(y,\tau)|^2G_\nu(y)\,dy=0
  \qquad\text{for every }\tau.
\]
Consequently no retained terminal core with positive Gaussian velocity mass can
belong to this branch.

\item Suppose the terminal state is axisymmetric and the circulation defect is
absorbable: for some \(\kappa>0\) and a.e. \(\tau\),
\[
  -\frac{1}{2\nu}\int_{\mathbb R^3}\Gamma^2(V\cdot y)G_\nu\,dy
  \le
  (1-\kappa)\int_{\mathbb R^3}\Gamma^2G_\nu\,dy
  +2\nu(1-\kappa)\int_{\mathbb R^3}|\nabla \Gamma|^2G_\nu\,dy .
\]
Then
\[
  \int_{\mathbb R^3}\Gamma(y,\tau)^2G_\nu(y)\,dy=0
  \qquad\text{for every }\tau.
\]
Thus every retained axisymmetric terminal core whose remaining obstruction is a
positive Gaussian circulation mass is excluded by the weighted circulation
identity.
\end{enumerate}
\end{theorem}

\begin{proof}
For the velocity statement set
\[
  \mathcal E(\tau)=\int_{\mathbb R^3}|V(y,\tau)|^2G_\nu(y)\,dy .
\]
The Gaussian virial identity gives
\[
  \partial_\tau\mathcal E
  =
  -2\nu\int_{\mathbb R^3}|\nabla V|^2G_\nu\,dy
  -\int_{\mathbb R^3}|V|^2G_\nu\,dy
  -\frac{1}{2\nu}\int_{\mathbb R^3}|V|^2(V\cdot y)G_\nu\,dy
  -\frac1\nu\int_{\mathbb R^3}P(V\cdot y)G_\nu\,dy .
\]
The absorbability hypothesis therefore implies
\[
  \partial_\tau\mathcal E
  +\kappa\mathcal E
  +2\nu\kappa\int_{\mathbb R^3}|\nabla V|^2G_\nu\,dy
  \le0 .
\]
In particular \(\partial_\tau\mathcal E+\kappa\mathcal E\le0\).  Integrating
from \(s<\tau\) to \(\tau\) gives
\[
  \mathcal E(\tau)\le e^{-\kappa(\tau-s)}\mathcal E(s).
\]
The ancient boundedness of the terminal state makes \(\sup_{s<\tau}\mathcal
E(s)<\infty\).  Sending \(s\to-\infty\) yields \(\mathcal E(\tau)=0\).

For the axisymmetric statement set
\[
  \mathcal G(\tau)=\int_{\mathbb R^3}\Gamma(y,\tau)^2G_\nu(y)\,dy .
\]
The weighted circulation identity is
\[
  \partial_\tau\mathcal G
  +\mathcal G
  +2\nu\int_{\mathbb R^3}|\nabla\Gamma|^2G_\nu\,dy
  =
  -\frac{1}{2\nu}\int_{\mathbb R^3}\Gamma^2(V\cdot y)G_\nu\,dy .
\]
The assumed absorption gives
\[
  \partial_\tau\mathcal G+\kappa\mathcal G
  +2\nu\kappa\int_{\mathbb R^3}|\nabla\Gamma|^2G_\nu\,dy
  \le0 .
\]
The same backward Gronwall argument gives \(\mathcal G(\tau)=0\) for every
\(\tau\).
\end{proof}

\begin{remark}[Controlled swirl alone]
\label{paper6:rem:controlled-swirl}
If the axisymmetric representative satisfies \(\Gamma\in L^\infty\), the same
weighted circulation identity gives the a priori estimate
\[
  \mathcal G(\tau)^{1/2}
  \lesssim
  \|V\|_{L^\infty}\|\Gamma\|_{L^\infty}
\]
for bounded ancient terminal states.  This estimate enters the local
decomposition of alternatives but does not by itself exclude a terminal state.  The exclusion
requires the stated absorption, vanishing conclusion, or an equivalent
condition which makes the virial inequality dissipative.
\end{remark}

\subsection{Assembly of Alternatives}

\begin{theorem}[Local Type II decomposition into local alternatives]
\label{paper6:thm:state-decomposition}
Every positive local Type II concentration sequence in \cref{paper6:def:local-typeII-sequence}, after passing to subsequences and selected windows, enters one of the following alternatives:
\begin{enumerate}[label=\textup{(\roman*)}]
\item a compact single-core vanishing-dissipation branch;
\item a multibubble, cascade, or gauge-degenerate branch;
\item a rough-core branch, meaning loss of local windowed \(H^1\) control on a
retained compact cylinder;
\item a finite-cost scale-collapse branch with nonvanishing selected core;
\item a scale-rigid state covered by the scale-rigidity hypothesis in
\cref{paper6:thm:multibubble}.
\end{enumerate}
In the last alternative, any terminal subbranch satisfying
\cref{paper6:thm:virial-exclusion} is eliminated before the remaining scale-rigid
state is passed to the scale-rigidity hypothesis.
\end{theorem}

\begin{proof}
This decomposition is obtained by combining
\cref{paper6:thm:representation,paper6:thm:multibubble,paper6:thm:rough-core,paper6:thm:scale-stratification}.
The representation theorem gives the nondegenerate single-core chart whenever a
unique retained core is selected.  Failure of that selection, splitting of
positive CKN mass, compactness failure, or scale cascade is the
multibubble/cascade side.  On a represented retained core, rough-core loss is
equivalent by \cref{paper6:thm:rough-core} to failure of compact CKN control on an
enclosing cylinder, hence returns to the multibubble/cascade side.  Genuine
scale collapse is divided by \cref{paper6:thm:scale-stratification} into finite-cost
scale collapse or one of the scale-rigid terminal states used in
\cref{paper6:thm:multibubble}.  These cases exhaust the local alternatives.
\end{proof}

\begin{proposition}[Elimination of local Type II states]
\label{paper6:prop:state-elimination}
No positive local Type II concentration sequence can lie in any of the
alternatives \textup{(i)}--\textup{(v)} of
\cref{paper6:thm:state-decomposition}.
\end{proposition}

\begin{proof}
Alternative \textup{(i)} is excluded by the compact single-core criterion,
\cref{paper6:thm:single-core}.

Alternative \textup{(ii)} is excluded by the local multibubble and cascade
theorem, \cref{paper6:thm:multibubble}, under its terminal admissibility and
scale-rigidity hypotheses.

Alternative \textup{(iii)} is not an independent terminal state.  By
\cref{paper6:thm:rough-core}, rough-core loss of local \(H^1\) control forces
failure of the compact-cylinder \(\mathcal C+\mathcal D\) bounds on enclosing
compact cylinders.  Such failure yields an active compact-cylinder
concentration branch and therefore belongs to the multibubble/cascade
alternative, already excluded by \cref{paper6:thm:multibubble}.

Alternative \textup{(iv)} is excluded by the finite-cost scale-collapse theorem,
\cref{paper6:thm:scale-cost}.

Alternative \textup{(v)} is the scale-rigid terminal state in
\cref{paper6:thm:multibubble}.  If this terminal state is velocity-virial dissipative,
\cref{paper6:thm:virial-exclusion} forces its Gaussian velocity mass to vanish, which
contradicts retention of a positive terminal core.  If it is axisymmetric and
swirl-virial dissipative with positive retained circulation mass, the same
theorem forces \(\Gamma\equiv0\), contradicting that positive circulation mass.
The remaining scale-rigid terminal states are covered by the
scale-rigidity hypothesis in \cref{paper6:thm:multibubble}.  Under that hypothesis, the
state either has a nonzero bounded selected-window limit, which is incompatible
with the local Type II condition, or it reduces to the compact single-core
branch, which is excluded by \cref{paper6:thm:single-core}.  Hence it cannot be a local
Type II branch.
\end{proof}

\subsection{Assembly Theorem}

\begin{theorem}[Local Type II exclusion]
\label{paper6:thm:local-typeII-exclusion}
Let \((u,p)\) be a suitable weak solution on
\(\mathbb R^3\times(T-\delta,T)\).  Then no point of \(\Sigma(T)\) lies in the
local Type II alternative, within the class of local alternatives covered by
\cref{paper6:thm:state-decomposition}.
\end{theorem}

\begin{proof}
Assume, for contradiction, that \(x_0\in\Sigma(T)\) lies in the local Type II
alternative.  By \cref{paper6:thm:paper0}, there is a positive local Type II
concentration sequence \(r_n\downarrow0\) at \((x_0,T)\).  By the decomposition into local alternatives, \cref{paper6:thm:state-decomposition}, this sequence lies,
after passing to selected windows, in one of the five alternatives listed
there.  Each alternative is impossible by \cref{paper6:prop:state-elimination}.  This
contradiction proves that no such \(x_0\) exists in the covered local
class of alternatives.
\end{proof}

\begin{corollary}[Final local dichotomy after Type I separation]
\label{paper6:cor:final-local-dichotomy}
Every singular point at time \(T\) in the covered class of local alternatives
belongs to the local Type I alternative.  Equivalently, once the Type I case is
covered by the Type I part of the theory, the local Type II class of alternatives contains no remaining singular branch.
\end{corollary}

\begin{proof}
Let \(x_0\in\Sigma(T)\).  By \cref{paper6:thm:paper0}, a concentration point is
either in the local Type I alternative or in the local Type II alternative.
\Cref{paper6:thm:local-typeII-exclusion} excludes the latter.
\end{proof}

\subsection{Conclusion}

The proof uses only local PDE statements.  The local concentration theorem gives
positive local CKN concentration and the Type I/Type II dichotomy.  The
repaired-gauge theorem gives the local renormalized representation for
nondegenerate selected cores.  The single-core theorem excludes compact
single-core vanishing-dissipation branches.  The multibubble/cascade theorem
excludes splitting, cascades, gauge degeneracy, and the scale-rigid terminal
state under its stated local hypotheses.  The virial identities eliminate the
velocity-virial dissipative terminal states and the axisymmetric
swirl-virial dissipative terminal states before the remaining scale-rigid
alternatives are covered by that hypothesis.  The rough-core estimate reduces
loss of windowed \(H^1\) control to failure of compact CKN control.  The
scale-collapse theorem excludes finite-cost scale collapse and decomposes the remaining scale-collapse alternatives locally.  Thus the local alternatives are exhausted without adding any further classification theorem.

\clearpage
\section{PDE Translation of the C1--C5 Classification Inputs}
\label{sec:c1-c5-pde-translation}

This section records, in PDE notation, the Type II classification inputs encoded
in the C1--C5 cluster notes.  The statements are conditional classification
statements for declared Type II branches.  They do not assert the existence of
such branches and do not add any regularity theorem beyond the analytic
hypotheses explicitly stated below.

\subsection{Declared Type II branches and exhaustion}

\begin{definition}[Declared local Type II class]
\label{def:c1-declared-typeII-class}
Let \(\mathcal U_{\mathrm{II}}^{NS}\) be the class of candidate terminal
branches \((u,p,T^*)\) such that:
\begin{enumerate}[label=\textup{(\roman*)}]
\item \((u,p)\) is a three-dimensional incompressible Navier--Stokes solution
on a terminal interval \([0,T^*)\) in the analytic class under consideration;
\item \(T^*<\infty\) is the terminal time of the branch;
\item the branch is not in the local Type I class associated with the
scale-invariant Type I criterion;
\item the branch is a Type II candidate rather than a continuation, Type I,
boundary, or non-Navier--Stokes-domain alternative.
\end{enumerate}
Membership in \(\mathcal U_{\mathrm{II}}^{NS}\) is only a declaration of the
classification problem; no nonemptiness assertion is made.
\end{definition}

\begin{definition}[Type II route data]
\label{def:c1-typeII-route-data}
A branch \(\omega=(u,p,T^*)\in\mathcal U_{\mathrm{II}}^{NS}\) has the Type II
route data if the following three analytic outputs are present:
\begin{enumerate}[label=\textup{(\roman*)}]
\item a concentration profile or blow-up germ modulo the Navier--Stokes
symmetry group;
\item a supercritical scale route, meaning that the branch enters the Type II
scale alternative rather than the Type I scale alternative;
\item a complete local profile output in the Navier--Stokes profile class used
for the subsequent compactness and representation arguments.
\end{enumerate}
The route is exhaustive if every declared Type II branch either has these three
outputs or falls into one of the ordered failures in
\cref{def:c1-exhaustion-failures}.
\end{definition}

\begin{definition}[Ordered exhaustion failures]
\label{def:c1-exhaustion-failures}
For a declared Type II branch, the ordered failures of the Type II route are:
\begin{enumerate}[label=\textup{(\roman*)}]
\item concentration extraction fails;
\item concentration extraction succeeds but the supercritical scale route fails;
\item concentration extraction and scale routing succeed but the local profile
output is not complete;
\item the declared Type II branch lies outside the certified
concentration/scale/profile route.
\end{enumerate}
The first applicable failure in this order is the emitted exhaustion
alternative.
\end{definition}

\begin{theorem}[Type II branch exhaustion]
\label{thm:c1-typeII-branch-exhaustion}
Assume the Type II route data of \cref{def:c1-typeII-route-data} are exhaustive
on \(\mathcal U_{\mathrm{II}}^{NS}\).  Then every
\(\omega\in\mathcal U_{\mathrm{II}}^{NS}\) has the concentration output, the
supercritical scale route, and the complete local profile output.
\end{theorem}

\begin{proof}
Fix \(\omega=(u,p,T^*)\in\mathcal U_{\mathrm{II}}^{NS}\).  By the
concentration-extraction component of the route data, the branch has a
concentration profile or blow-up germ.  By the scale-routing component, the
same branch lies in the supercritical Type II scale route.  By profile
completeness, the extracted profile belongs to the local Navier--Stokes profile
class used in the subsequent arguments.  Since \(\omega\) was arbitrary, the
three outputs hold for every declared Type II branch.
\end{proof}

\begin{corollary}[Ordered C1 classification]
\label{cor:c1-ordered-exhaustion}
For each \(\omega\in\mathcal U_{\mathrm{II}}^{NS}\), exactly one ordered output
occurs: the Type II route data hold, or one of the four failures in
\cref{def:c1-exhaustion-failures} occurs.  The Type II route is exhaustive on
\(\mathcal U_{\mathrm{II}}^{NS}\) if and only if the first output holds for
every \(\omega\in\mathcal U_{\mathrm{II}}^{NS}\).
\end{corollary}

\begin{proof}
Test the four route requirements in the stated order.  If concentration
extraction, scale routing, and profile completeness all hold, the branch has
the Type II route data.  If a test fails, the first failed test gives the
corresponding ordered failure.  The alternatives are disjoint by the first
failure convention and exhaustive by construction.  The final equivalence is
the definition of exhaustion on \(\mathcal U_{\mathrm{II}}^{NS}\).
\end{proof}

\begin{theorem}[Entry into the downstream Type II analysis]
\label{thm:c1-downstream-entry}
Assume the Type II route is exhaustive.  Then each declared Type II branch
enters the downstream analysis in exactly one of the following ways:
\begin{enumerate}[label=\textup{(\roman*)}]
\item it reaches the compact Type II barrier hypotheses after the
representation, critical-mass, tightness, and windowed-\(H^1\) tests;
\item it falls into an ordered representation failure, critical-mass failure,
loss of critical tightness, or loss of local windowed \(H^1\) control.
\end{enumerate}
\end{theorem}

\begin{proof}
By \cref{thm:c1-typeII-branch-exhaustion}, every declared Type II branch has
the concentration, scale, and profile route data.  The representation step then
either gives a repaired-gauge orbit or one of the ordered representation
failures in \cref{def:c2-representation-failures}.  If a repaired-gauge orbit
exists, the critical \(L^3\) test gives either the positive finite annulus or
one of the ordered critical-mass alternatives in
\cref{def:c3-critical-mass-alternatives}.  If the positive finite annulus
holds, the remaining compact-barrier inputs are critical tightness and local
windowed \(H^1\) control.  Their failures are exactly the radiative/noncompact
and rough-core alternatives of \cref{def:c5-compact-tests}.  If both hold, the
compact Type II barrier applies.
\end{proof}

\subsection{Repaired-gauge representation from a Type II route}

\begin{definition}[Raw Type II chart]
\label{def:c2-raw-chart}
A raw Type II chart consists of
\[
  (u,p,T^*,x_c,\lambda,\tau,V,P,\mathcal I),
  \qquad \mathcal I=(t_0,T^*),
\]
where \(u,p\) solve
\[
  \partial_tu+(u\cdot\nabla)u+\nabla p=\nu\Delta u,\qquad
  \nabla\cdot u=0,
\]
\(\lambda(t)>0\), \(x_c(t)\in\mathbb R^3\), \(x_c,\lambda\in AC_{\rm loc}\),
\(\lambda(t)\to0\) as \(t\uparrow T^*\), and
\[
  \frac{d\tau}{dt}=\lambda(t)^{-2},\qquad
  \tau(t)\to\infty\quad\text{as }t\uparrow T^*.
\]
The renormalized variables are
\[
  y=\frac{x-x_c(t)}{\lambda(t)},\qquad
  u(x,t)=\lambda(t)^{-1}V(y,\tau(t)),\qquad
  p(x,t)=\lambda(t)^{-2}P(y,\tau(t)),
\]
with enough regularity to justify the distributional chain rule on compact
\(y\)-sets.
\end{definition}

\begin{lemma}[Raw renormalized equation]
\label{lem:c2-raw-renormalized-equation}
For a raw Type II chart, define
\[
  a_{\rm raw}(\tau(t))=\lambda(t)\lambda_t(t),\qquad
  b_{\rm raw}(\tau(t))=\lambda(t)x_c'(t).
\]
Then \(V,P\) satisfy, distributionally on compact \(y\)-sets,
\[
  \partial_\tau V+(V\cdot\nabla)V+\nabla P
  =
  \nu\Delta V
  +a_{\rm raw}(\tau)(V+y\cdot\nabla V)
  +b_{\rm raw}(\tau)\cdot\nabla V,
  \qquad \nabla\cdot V=0 .
\]
\end{lemma}

\begin{proof}
Write
\[
u(x,t)=\lambda(t)^{-1}V(y,\tau(t)),\qquad
y=\frac{x-x_c(t)}{\lambda(t)},\qquad \tau_t=\lambda^{-2}.
\]
Then
\[
  \nabla_xu=\lambda^{-2}\nabla_yV,\qquad
  \Delta_xu=\lambda^{-3}\Delta_yV,\qquad
  (u\cdot\nabla_x)u=\lambda^{-3}(V\cdot\nabla_y)V.
\]
Since
\[
  y_t=-\frac{x_c'(t)}{\lambda(t)}
      -\frac{\lambda_t(t)}{\lambda(t)}y,
\]
one has
\[
  \partial_tu
  =
  \lambda^{-3}\partial_\tau V
  -\lambda_t\lambda^{-2}(V+y\cdot\nabla V)
  -\lambda^{-2}x_c'(t)\cdot\nabla V .
\]
Moreover \(\nabla_xp=\lambda^{-3}\nabla_yP\).  Substitution in
Navier--Stokes and multiplication by \(\lambda^3\) give
\[
  \partial_\tau V-\lambda\lambda_t(V+y\cdot\nabla V)
  -\lambda x_c'(t)\cdot\nabla V
  +(V\cdot\nabla)V+\nabla P
  =\nu\Delta V .
\]
Moving the modulation terms to the right-hand side gives the equation with
\(a_{\rm raw}=\lambda\lambda_t\) and \(b_{\rm raw}=\lambda x_c'\).  Finally,
\(\nabla_x\cdot u=\lambda^{-2}\nabla_y\cdot V\), so incompressibility of \(u\)
implies \(\nabla_y\cdot V=0\).
\end{proof}

\begin{definition}[Repaired-gauge realization]
\label{def:c2-repaired-gauge-realization}
Let \(0<p<3\), \(\Theta_0>0\), and let \(\psi_R\) be a smooth compactly
supported centering cutoff.  A raw Type II chart admits a repaired-gauge
realization if, after an admissible time-dependent scale and translation
correction, the profile satisfies
\[
  G_{\rm sc}(V)=\int_{\mathbb R^3}|y|^{-p}|V(y)|^3\,dy-\Theta_0=0,
\]
and
\[
  G_j(V)=\int_{\mathbb R^3}y_j|V(y)|^2\psi_R(y)\,dy=0,
  \qquad j=1,2,3,
\]
for a.e. renormalized time.  The corrected profile belongs to the admissible
gauge class \(\mathcal A_p\), has the regularity needed to differentiate these
functionals, and all time-dependent symmetry corrections have been absorbed
into final modulation coefficients \(a,b\).
\end{definition}

\begin{definition}[Repaired-gauge representation]
\label{def:c2-representation}
A Type II branch has a repaired-gauge representation if it has:
\begin{enumerate}[label=\textup{(\roman*)}]
\item a raw Type II chart;
\item a repaired-gauge realization;
\item pressure reconstruction modulo functions of time;
\item final modulation coefficients \(a,b\) appearing in the repaired-gauge
equation.
\end{enumerate}
The output is the tuple
\[
  (V,P,a,b,G_{\rm sc},G_1,G_2,G_3,\tau_0).
\]
\end{definition}

\begin{theorem}[C2 representation implication]
\label{thm:c2-representation-implication}
Assume a Type II profile branch has a raw Type II chart, a repaired-gauge
realization, pressure reconstruction modulo functions of time, and final
modulation coefficients.  Then the branch has a repaired-gauge representation
in the sense of \cref{def:c2-representation}.  In particular,
\((V,P,a,b)\) satisfies the PDE class required by the compact Type II barrier.
\end{theorem}

\begin{proof}
The raw chart gives \(V,P\) and, by \cref{lem:c2-raw-renormalized-equation},
the renormalized Navier--Stokes equation with raw modulation coefficients.  The
repaired-gauge realization places the profile on the scale and centering gauge
surface and absorbs the time-dependent symmetry correction into updated
coefficients.  The pressure reconstruction identifies the pressure associated
with \(V\), modulo the usual functions of time.  The modulation data identify
the final coefficients \(a,b\).  These are the components of
\cref{def:c2-representation}.
\end{proof}

\begin{definition}[Ordered representation failures]
\label{def:c2-representation-failures}
If a repaired-gauge representation is not obtained, the ordered failures are:
\begin{enumerate}[label=\textup{(\roman*)}]
\item no raw orbit is supplied by the concentration/profile output;
\item the raw orbit exists but the repaired gauge cannot be realized;
\item pressure reconstruction modulo functions of time is unavailable;
\item the final modulation coefficients are unavailable.
\end{enumerate}
The first applicable failure is the emitted representation alternative.
\end{definition}

\begin{corollary}[Ordered C2 representation classification]
\label{cor:c2-ordered-representation}
Every Type II profile branch in the declared class has exactly one ordered
output: a repaired-gauge representation, or one of the four failures in
\cref{def:c2-representation-failures}.
\end{corollary}

\begin{proof}
Verify the four representation inputs in order.  If all are present,
\cref{thm:c2-representation-implication} gives the repaired-gauge
representation.  If not, the first missing input gives the corresponding
ordered failure.  This is exhaustive and disjoint by first-failure ordering.
\end{proof}

\subsection{Minimal chart and gauge realization}

\begin{definition}[Minimal representation data]
\label{def:c2r-minimal-representation-data}
The minimal data for a routed Type II branch consist of:
\begin{enumerate}[label=\textup{(\roman*)}]
\item an absolutely continuous concentration chart
\[
  (u,p,T^*,x_c,\lambda,\rho,U,Q,\mathcal I),
\]
where \(\mathcal I=(t_0,T^*)\), \(\lambda(t)>0\), \(\lambda(t)\to0\),
\(\rho_t=\lambda^{-2}\), and
\[
  u(x,t)=\lambda(t)^{-1}U(y,\rho(t)),\qquad
  p(x,t)=\lambda(t)^{-2}Q(y,\rho(t))+c(t),\qquad
  y=\frac{x-x_c(t)}{\lambda(t)};
\]
\item an absolutely continuous repaired-gauge correction consisting of
\(\rho=\rho(\tau)\), \(\mu(\tau)>0\), and \(q(\tau)\in\mathbb R^3\), with
\[
  \rho_\tau=\mu(\tau)^2,
\]
\[
  V(Y,\tau)=\mu(\tau)U(\mu(\tau)Y+q(\tau),\rho(\tau)),\qquad
  P(Y,\tau)=\mu(\tau)^2Q(\mu(\tau)Y+q(\tau),\rho(\tau)),
\]
and
\[
  G_{\rm sc}(V(\tau))=0,\qquad
  G_j(V(\tau))=0,\quad j=1,2,3.
\]
The condition \(\rho_\tau=\mu^2\) preserves the viscosity coefficient after the
time-dependent scale correction.
\end{enumerate}
\end{definition}

\begin{lemma}[Chart data imply the raw orbit]
\label{lem:c2r-chart-raw-orbit}
Under the chart data in \cref{def:c2r-minimal-representation-data}, \(U,Q\)
satisfy
\[
  \partial_\rho U+(U\cdot\nabla)U+\nabla Q
  =
  \nu\Delta U+A(\rho)(U+y\cdot\nabla U)+B(\rho)\cdot\nabla U,
  \qquad \nabla\cdot U=0,
\]
where
\[
  A(\rho(t))=\lambda(t)\lambda_t(t),\qquad
  B(\rho(t))=\lambda(t)x_c'(t).
\]
\end{lemma}

\begin{proof}
The chart data are the raw orbit data with raw time \(\rho\).  The additive
pressure function \(c(t)\) has zero spatial gradient.  The chain-rule
computation is the one in \cref{lem:c2-raw-renormalized-equation}, with
\(\tau\) replaced by \(\rho\).  Thus spatial derivatives scale as
\(\nabla_xu=\lambda^{-2}\nabla_yU\), \(\Delta_xu=\lambda^{-3}\Delta_yU\), the
convection term scales as \(\lambda^{-3}(U\cdot\nabla)U\), and
\(\rho_t=\lambda^{-2}\).  The terms generated by \(\lambda_t\) and \(x_c'\)
give \(A=\lambda\lambda_t\) and \(B=\lambda x_c'\).  Incompressibility pulls
back as \(\nabla_x\cdot u=\lambda^{-2}\nabla_y\cdot U\).
\end{proof}

\begin{lemma}[Pressure pullback]
\label{lem:c2r-pressure-pullback}
Under the chart data,
\[
  -\Delta Q=\partial_i\partial_j(U_iU_j)
\]
in distributions, modulo functions of \(\rho\).  After any admissible
repaired-gauge correction, the transformed pressure satisfies
\[
  -\Delta P=\partial_i\partial_j(V_iV_j)
\]
in distributions, modulo functions of \(\tau\).
\end{lemma}

\begin{proof}
Taking divergence of the physical Navier--Stokes equation gives
\[
  -\Delta p=\partial_i\partial_j(u_iu_j),
\]
with pressure determined up to a function of time.  Pulling this identity back
through \(p=\lambda^{-2}Q+c(t)\), \(u=\lambda^{-1}U\), and
\(y=(x-x_c)/\lambda\) gives the pressure equation for \(Q\).  If
\[
  V(Y)=\mu U(\mu Y+q),\qquad P(Y)=\mu^2Q(\mu Y+q),
\]
then
\[
  -\Delta_YP
  =\mu^4(-\Delta Q)(\mu Y+q)
  =\mu^4\partial_i\partial_j(U_iU_j)(\mu Y+q)
  =\partial_i\partial_j(V_iV_j).
\]
Adding a function of \(\rho\) or \(\tau\) to the pressure does not affect the
identity.
\end{proof}

\begin{lemma}[Forced modulation coefficients]
\label{lem:c2r-forced-modulation}
Assume the minimal data of \cref{def:c2r-minimal-representation-data}.  Then
\((V,P)\) satisfies
\[
  \partial_\tau V+(V\cdot\nabla)V+\nabla P
  =
  \nu\Delta V
  +a(\tau)(V+Y\cdot\nabla V)
  +b(\tau)\cdot\nabla V,
\]
where
\[
  a(\tau)=\mu(\tau)^2A(\rho(\tau))+\frac{\mu_\tau(\tau)}{\mu(\tau)}
\]
and
\[
  b(\tau)=\mu(\tau)B(\rho(\tau))
      +\mu(\tau)A(\rho(\tau))q(\tau)
      +\frac{q_\tau(\tau)}{\mu(\tau)}.
\]
\end{lemma}

\begin{proof}
Let
\[
  z=\mu(\tau)Y+q(\tau),\qquad
  V(Y,\tau)=\mu(\tau)U(z,\rho(\tau)),\qquad
  \rho_\tau=\mu^2.
\]
The raw equation in \((z,\rho)\) is
\[
  U_\rho+(U\cdot\nabla_z)U+\nabla_zQ
  =
  \nu\Delta_zU+A(U+z\cdot\nabla_zU)+B\cdot\nabla_zU.
\]
The identities
\[
  \nabla_YV=\mu^2\nabla_zU,\qquad
  \Delta_YV=\mu^3\Delta_zU,\qquad
  (V\cdot\nabla_Y)V=\mu^3(U\cdot\nabla_z)U,\qquad
  \nabla_YP=\mu^3\nabla_zQ
\]
together with differentiation in \(\tau\) give
\[
  V_\tau+(V\cdot\nabla_Y)V+\nabla_YP-\nu\Delta_YV
  =
  \left(\frac{\mu_\tau}{\mu}+\mu^2A\right)(V+Y\cdot\nabla_YV)
  +
  \left(\mu B+\mu Aq+\frac{q_\tau}{\mu}\right)\cdot\nabla_YV.
\]
This is the displayed equation.  Since \(\mu,q,\rho\) are absolutely continuous
and \(A,B\) are the raw chart coefficients, \(a,b\) are measurable and locally
integrable on every finite renormalized window.
\end{proof}

\begin{theorem}[Minimal representation theorem]
\label{thm:c2r-minimal-representation}
If a routed Type II branch has the minimal data of
\cref{def:c2r-minimal-representation-data}, then it has a repaired-gauge
representation in the sense of \cref{def:c2-representation}.
\end{theorem}

\begin{proof}
\Cref{lem:c2r-chart-raw-orbit} gives the raw orbit.  The gauge correction
places the profile on the repaired scale and centering surface, hence gives the
repaired-gauge realization.  \Cref{lem:c2r-pressure-pullback} gives pressure
reconstruction modulo functions of time.  \Cref{lem:c2r-forced-modulation}
gives the final modulation coefficients.  These are the four inputs in
\cref{def:c2-representation}.
\end{proof}

\begin{definition}[Minimal representation failures]
\label{def:c2r-minimal-failures}
Outside the declared repaired-gauge representation class, the ordered failures
are:
\begin{enumerate}[label=\textup{(\roman*)}]
\item no absolutely continuous concentration chart is available;
\item the chart exists but the repaired scale and centering gauges cannot be
solved admissibly and absolutely continuously.
\end{enumerate}
Pressure-pullback and modulation-coefficient failures are not independent
alternatives when the chart and the absolutely continuous gauge solve exist.
\end{definition}

\begin{corollary}[Ordered minimal representation classification]
\label{cor:c2r-ordered-classification}
Every routed Type II branch has exactly one ordered output: a repaired-gauge
representation, failure of the raw chart, or failure of the repaired-gauge
solve.
\end{corollary}

\begin{proof}
If the chart is absent, the first failure occurs.  If the chart exists but the
admissible absolutely continuous gauge solve is absent, the second failure
occurs.  If both are present, \cref{thm:c2r-minimal-representation} gives the
repaired-gauge representation.  The alternatives are disjoint by their order
and exhaustive for the representation row.
\end{proof}

\subsection{\texorpdfstring{Critical \(L^3\) annulus}{Critical L3 annulus}}

\begin{definition}[Critical mass and positive finite annulus]
\label{def:c3-critical-mass}
For a repaired-gauge renormalized Type II candidate \((V,P,a,b)\) on
\([\tau_0,\infty)\), define
\[
  N_3(\tau)=\|V(\tau)\|_{L^3(\mathbb R^3)}.
\]
The orbit has positive finite critical mass if there are constants
\(0<\eta\le M<\infty\) such that
\[
  \eta\le N_3(\tau)\le M,\qquad \tau\ge\tau_0.
\]
This condition replaces the exact normalization
\(\|V(\tau)\|_{L^3}=1\); no amplitude rescaling is performed.
\end{definition}

\begin{definition}[Critical-mass alternatives]
\label{def:c3-critical-mass-alternatives}
For a represented orbit, the ordered critical-mass alternatives are:
\begin{enumerate}[label=\textup{(\roman*)}]
\item \(N_3\) is not a well-defined extended measurable critical norm on the
represented branch;
\item \(N_3(\tau)=\infty\) for some \(\tau\), or
\(\sup_{\tau\ge\tau_0}N_3(\tau)=\infty\);
\item \(\sup_{\tau\ge\tau_0}N_3(\tau)<\infty\) and
\(\liminf_{\tau\to\infty}N_3(\tau)=0\);
\item after discarding at most a finite initial time interval, the positive
finite critical annulus of \cref{def:c3-critical-mass} holds.
\end{enumerate}
The first applicable alternative is emitted.
\end{definition}

\begin{theorem}[Nonzero-mass compact Type II barrier]
\label{thm:c3-nonzero-mass-barrier}
Let \((V,P,a,b)\) be a repaired-gauge renormalized Type II candidate.  Assume:
\begin{enumerate}[label=\textup{(\roman*)}]
\item the positive finite critical annulus holds:
\[
  0<\eta\le \|V(\tau)\|_{L^3(\mathbb R^3)}\le M<\infty;
\]
\item uniform critical tightness holds:
\[
  \forall\varepsilon>0\ \exists R_\varepsilon:\quad
  \sup_{\tau\ge\tau_0}
  \int_{|y|>R_\varepsilon}|V(y,\tau)|^3\,dy<\varepsilon;
\]
\item local windowed \(H^1\) control holds:
\[
  \forall m\ge1:\quad
  \sup_{n\in\mathbb N}
  \int_{\tau_0+n}^{\tau_0+n+1}
  \|V(\tau)\|_{H^1(B_m)}^2\,d\tau<\infty.
\]
\end{enumerate}
Then finite total localized renormalization cost is impossible:
\[
  \int_{\tau_0}^{\infty}
  \tilde{\mathfrak D}_{R_0}(\tau)\,d\tau=\infty.
\]
\end{theorem}

\begin{proof}
The proof is the compact Type II barrier proof with the exact normalization
\(\|V(\tau)\|_{L^3}=1\) replaced by the lower bound
\(N_3(\tau)\ge\eta\).  Assume for contradiction that the total localized
renormalization cost is finite.  The good-window selection and compactness
portion of the compact barrier gives times \(\sigma_j\to\infty\) such that,
after passing to a subsequence,
\[
  V(\sigma_j)\to V_\infty\quad\text{strongly in }L^3_{\rm loc},
\]
the selected sequence has vanishing local dissipation on the exhaustion used by
the compact barrier, and consequently
\[
  \nabla V_\infty=0
\]
distributionally on every fixed ball.  These conclusions use only tightness,
local windowed \(H^1\) control, and finite total cost.

The low-dissipation limit is spatially constant on each fixed ball.  The
constants are compatible as the balls increase, so \(V_\infty\equiv c\) on
\(\mathbb R^3\).  The upper critical-mass bound gives, for every \(R\),
\[
  |c|^3|B_R|
  =
  \lim_{j\to\infty}\int_{B_R}|V(y,\sigma_j)|^3\,dy
  \le M^3 .
\]
Letting \(R\to\infty\) forces \(c=0\).  Thus
\[
  V(\sigma_j)\to0\quad\text{strongly in }L^3_{\rm loc}.
\]
Choose \(R\) so that the critical \(L^3\)-tail outside \(B_R\) is
\(<\eta^3/4\) uniformly in \(\tau\).  Strong convergence on \(B_R\) gives
\[
  \int_{B_R}|V(y,\sigma_j)|^3\,dy<\eta^3/4
\]
for \(j\) large.  Therefore
\[
  \|V(\sigma_j)\|_{L^3(\mathbb R^3)}^3<\eta^3/2,
\]
contradicting \(N_3(\sigma_j)\ge\eta\).  Hence the total localized
renormalization cost is infinite.
\end{proof}

\begin{theorem}[\texorpdfstring{Critical \(L^3\)-annulus classification}{Critical L3-annulus classification}]
\label{thm:c3-l3-annulus-classification}
For a represented repaired-gauge Type II candidate, exactly one of the ordered
alternatives in \cref{def:c3-critical-mass-alternatives} occurs.  In the fourth
case the compact Type II barrier remains valid with the positive finite
critical annulus in place of exact unit \(L^3\)-normalization.
\end{theorem}

\begin{proof}
If \(N_3\) is not a well-defined measurable extended norm, the first
alternative holds.  Otherwise \(N_3\) is an extended nonnegative measurable
function.  If it is infinite at some time, or has infinite supremum on the
tail, the second alternative holds.  If the supremum is finite but
\(\liminf_{\tau\to\infty}N_3(\tau)=0\), the third alternative holds.  In the
remaining case,
\[
  0<\liminf_{\tau\to\infty}N_3(\tau)
  \quad\text{and}\quad
  \sup_{\tau\ge\tau_0}N_3(\tau)<\infty.
\]
After discarding a finite initial time interval, there exist
\(0<\eta\le M<\infty\) such that
\[
  \eta\le N_3(\tau)\le M
\]
for all remaining \(\tau\).  Relabeling this tail initial time as \(\tau_0\)
gives the fourth alternative.  The barrier statement is
\cref{thm:c3-nonzero-mass-barrier}.
\end{proof}

\subsection{Localized cost and the Type II barrier condition}

\begin{definition}[Localized Type II cost]
\label{def:c4-localized-cost}
For a repaired-gauge Type II candidate, let
\[
  \tilde{\mathfrak D}_{R_0}(\tau)
  =
  \nu\int_{\mathbb R^3}|\nabla V(y,\tau)|^2\phi_{R_0}(y)\,dy
  +
  a_+(\tau)\int_{\mathbb R^3}|V(y,\tau)|^2\phi_{R_0}(y)\,dy,
\]
where \(a_+=\max(a,0)\), \(\phi_{R_0}\) is the compact Type II cutoff, and
\(\nu>0\).  A Type II barrier cost is a measurable nonnegative function
\(\mathfrak C_{\mathrm{II}}\) that is locally integrable on finite
\(\tau\)-intervals and whose divergence on \([\tau_0,\infty)\) is accepted as
the blocked Type II scale condition.
\end{definition}

\begin{definition}[Cost comparison]
\label{def:c4-cost-comparison}
A cost comparison for a repaired-gauge Type II candidate consists of
\(R_0,\phi_{R_0},\mathfrak C_{\mathrm{II}},c_0,\tau_1\), with \(c_0>0\) and
\(\tau_1\ge\tau_0\), such that
\[
  \mathfrak C_{\mathrm{II}}(\tau)
  \ge c_0\,\tilde{\mathfrak D}_{R_0}(\tau)
  \qquad\text{for a.e. }\tau\ge\tau_1.
\]
The canonical Navier--Stokes choice is
\[
  \mathfrak C_{\mathrm{II}}(\tau)=\tilde{\mathfrak D}_{R_0}(\tau),
  \qquad c_0=1,\qquad \tau_1=\tau_0.
\]
\end{definition}

\begin{lemma}[Identity-cost bridge]
\label{lem:c4-identity-cost}
If \(\tilde{\mathfrak D}_{R_0}\) is measurable, nonnegative, and locally
integrable on finite \(\tau\)-intervals, then the canonical choice
\(\mathfrak C_{\mathrm{II}}=\tilde{\mathfrak D}_{R_0}\) defines a valid Type II
barrier cost and satisfies the cost comparison with \(c_0=1\) and
\(\tau_1=\tau_0\).
\end{lemma}

\begin{proof}
The assumed measurability, nonnegativity, and local finite-interval
integrability give the cost properties in \cref{def:c4-localized-cost}.  The
identity
\[
  \mathfrak C_{\mathrm{II}}=\tilde{\mathfrak D}_{R_0}
\]
gives the comparison in \cref{def:c4-cost-comparison} with \(c_0=1\) and
\(\tau_1=\tau_0\).
\end{proof}

\begin{theorem}[Cost bridge]
\label{thm:c4-cost-bridge}
Let \((V,P,a,b)\) be a repaired-gauge Type II candidate.  Assume:
\begin{enumerate}[label=\textup{(\roman*)}]
\item \(\tilde{\mathfrak D}_{R_0}\) is measurable, nonnegative, and locally
integrable on finite \(\tau\)-intervals;
\item
\[
  \int_{\tau_0}^{\infty}
  \tilde{\mathfrak D}_{R_0}(\tau)\,d\tau=\infty;
\]
\item either the canonical cost
\(\mathfrak C_{\mathrm{II}}=\tilde{\mathfrak D}_{R_0}\) is used, or a cost
comparison as in \cref{def:c4-cost-comparison} is available.
\end{enumerate}
Then the blocked Type II scale condition associated with
\(\mathfrak C_{\mathrm{II}}\) holds.
\end{theorem}

\begin{proof}
In the canonical case, \cref{lem:c4-identity-cost} supplies the cost comparison.
In the general case, the comparison is assumed.  Hence there are
\(c_0>0\) and \(\tau_1\ge\tau_0\) such that
\[
  \mathfrak C_{\mathrm{II}}(\tau)
  \ge c_0\tilde{\mathfrak D}_{R_0}(\tau)
\]
for a.e. \(\tau\ge\tau_1\).  Since
\(\tilde{\mathfrak D}_{R_0}\ge0\) and is locally integrable on finite
intervals, divergence on \([\tau_0,\infty)\) implies
\[
  \int_{\tau_1}^{\infty}
  \tilde{\mathfrak D}_{R_0}(\tau)\,d\tau=\infty.
\]
Therefore
\[
  \int_{\tau_1}^{\infty}
  \mathfrak C_{\mathrm{II}}(\tau)\,d\tau
  \ge
  c_0\int_{\tau_1}^{\infty}
  \tilde{\mathfrak D}_{R_0}(\tau)\,d\tau
  =\infty.
\]
Since \(\mathfrak C_{\mathrm{II}}\) is nonnegative and locally integrable on
finite intervals,
\[
  \int_{\tau_0}^{\infty}
  \mathfrak C_{\mathrm{II}}(\tau)\,d\tau
  =
  \int_{\tau_0}^{\tau_1}\mathfrak C_{\mathrm{II}}(\tau)\,d\tau
  +
  \int_{\tau_1}^{\infty}\mathfrak C_{\mathrm{II}}(\tau)\,d\tau
  =\infty.
\]
By definition of the Type II barrier cost, this divergence is the blocked Type
II scale condition.
\end{proof}

\begin{corollary}[Compact barrier implies blocked scale]
\label{cor:c4-compact-barrier-blocks-scale}
Assume the hypotheses of \cref{thm:c3-nonzero-mass-barrier} and use the canonical
Navier--Stokes cost \(\mathfrak C_{\mathrm{II}}=\tilde{\mathfrak D}_{R_0}\).
Then the blocked Type II scale condition holds.  If, in addition, the
supercritical scale route is present and the corresponding Navier--Stokes
promotion hypothesis is available, then the promoted Type II suppression
condition holds.
\end{corollary}

\begin{proof}
\Cref{thm:c3-nonzero-mass-barrier} gives
\[
  \int_{\tau_0}^{\infty}
  \tilde{\mathfrak D}_{R_0}(\tau)\,d\tau=\infty.
\]
By \cref{lem:c4-identity-cost}, the canonical cost satisfies the cost comparison.
\Cref{thm:c4-cost-bridge} converts the infinite PDE cost into the blocked Type
II scale condition.  Combining the blocked condition with the supercritical
scale route gives the promoted suppression condition only under the stated
Navier--Stokes promotion hypothesis.
\end{proof}

\subsection{Two-bucket finite-cost classification}

\begin{definition}[Remaining compact-barrier tests]
\label{def:c5-compact-tests}
For a represented repaired-gauge Type II orbit, define the critical-tightness
condition by
\[
  \forall\varepsilon>0\ \exists R_\varepsilon:\quad
  \sup_{\tau\ge\tau_0}
  \int_{|y|>R_\varepsilon}|V(y,\tau)|^3\,dy<\varepsilon.
\]
Its failure is the radiative/noncompact alternative:
\[
  \exists\varepsilon_0>0\ \forall R>0\ \exists\tau_R\ge\tau_0:\quad
  \int_{|y|>R}|V(y,\tau_R)|^3\,dy\ge\varepsilon_0.
\]
Define local windowed \(H^1\) control by
\[
  \forall m\ge1:\quad
  \sup_{n\in\mathbb N}
  \int_{\tau_0+n}^{\tau_0+n+1}
  \|V(\tau)\|_{H^1(B_m)}^2\,d\tau<\infty.
\]
Its failure is the rough-core alternative:
\[
  \exists m\ge1:\quad
  \sup_{n\in\mathbb N}
  \int_{\tau_0+n}^{\tau_0+n+1}
  \|V(\tau)\|_{H^1(B_m)}^2\,d\tau=\infty.
\]
\end{definition}

\begin{definition}[Classification-ready compact route]
\label{def:c5-class-complete}
A declared Type II branch reaches the classification-ready compact route if
it has:
\begin{enumerate}[label=\textup{(\roman*)}]
\item the Type II route data of \cref{def:c1-typeII-route-data};
\item the repaired-gauge representation of \cref{def:c2-representation};
\item the positive finite critical annulus of \cref{def:c3-critical-mass};
\item the locally integrable nonnegative cost and cost comparison of
\cref{def:c4-localized-cost,def:c4-cost-comparison}.
\end{enumerate}
This condition does not include critical tightness or windowed \(H^1\) control;
those are the two remaining compact-barrier tests.
\end{definition}

\begin{lemma}[Compact positive branch is suppressed]
\label{lem:c5-compact-positive-suppressed}
Assume a declared Type II candidate reaches the classification-ready compact
route and satisfies both compact-barrier tests in \cref{def:c5-compact-tests}.
If the Navier--Stokes promotion hypothesis for the blocked Type II scale
condition is available, then the candidate satisfies the promoted Type II
suppression condition.
\end{lemma}

\begin{proof}
The classification-ready compact route gives the supercritical scale route,
a repaired-gauge orbit \((V,P,a,b)\), the positive finite critical annulus, and
the cost comparison.  Together with critical tightness and local windowed
\(H^1\) control, the hypotheses of \cref{thm:c3-nonzero-mass-barrier} hold.
Thus
\[
  \int_{\tau_0}^{\infty}
  \tilde{\mathfrak D}_{R_0}(\tau)\,d\tau=\infty.
\]
\Cref{thm:c4-cost-bridge} converts this divergence into the blocked Type II
scale condition.  Combining the blocked condition with the supercritical scale
route gives the promoted Type II suppression condition under the stated
Navier--Stokes promotion hypothesis.
\end{proof}

\begin{theorem}[Suppressed-or-two-bucket classification]
\label{thm:c5-suppressed-or-two-bucket}
Assume every declared Type II candidate reaches the classification-ready
compact route and that the Navier--Stokes promotion hypothesis is available.
Then each declared Type II candidate has exactly one ordered output:
\begin{enumerate}[label=\textup{(\roman*)}]
\item the promoted compact Type II suppression condition;
\item the radiative/noncompact alternative, i.e. failure of critical tightness;
\item the rough-core alternative, i.e. failure of local windowed \(H^1\) control
after critical tightness has been verified.
\end{enumerate}
\end{theorem}

\begin{proof}
Fix a declared Type II candidate.  Since it reaches the
classification-ready compact route, only the two compact-barrier tests
remain.  First evaluate critical tightness.  If tightness fails, the candidate
is in the radiative/noncompact alternative.  If tightness holds, evaluate
local windowed \(H^1\) control.  If the windowed \(H^1\) condition fails, the
candidate is in the rough-core alternative.  If it holds, then
\cref{lem:c5-compact-positive-suppressed} gives the promoted compact Type II
suppression condition.  The outputs are exhaustive and disjoint by this ordered
evaluation.
\end{proof}

\begin{definition}[Finite-cost non-suppressed branch]
\label{def:c5-finite-cost-nonsuppressed}
A declared Type II branch is finite-cost if
\[
  \int_{\tau_0}^{\infty}\mathfrak C_{\mathrm{II}}(\tau)\,d\tau<\infty.
\]
In the canonical Navier--Stokes cost this is equivalent to finite localized PDE
cost, because
\(\mathfrak C_{\mathrm{II}}=\tilde{\mathfrak D}_{R_0}\).  It is
non-suppressed if the promoted compact Type II suppression condition is not
emitted for that branch.
\end{definition}

\begin{corollary}[Two-bucket finite-cost classification]
\label{cor:c5-two-bucket-finite-cost}
Under the hypotheses of \cref{thm:c5-suppressed-or-two-bucket}, every
finite-cost non-suppressed declared Type II branch is either
radiative/noncompact or rough-core.
\end{corollary}

\begin{proof}
Apply \cref{thm:c5-suppressed-or-two-bucket}.  Since the branch is
non-suppressed, the first output is unavailable.  Therefore the ordered output
is the radiative/noncompact alternative or the rough-core alternative.  The
finite-cost assumption records that the resulting branch remains finite-cost;
the exclusion of the suppressed case is the explicit
non-suppression hypothesis.
\end{proof}

\begin{remark}[Content of the C1--C5 translation]
The preceding statements translate the C1--C5 cluster proof obligations into
PDE form.  Exhaustion failure, missing repaired-gauge representation,
ambiguous critical \(L^3\) normalization, and cost-comparison failure are all
replaced by ordered alternatives.  Once those alternatives are excluded or
assumed absent, the only finite-cost non-suppressed Type II alternatives left
by C5 are the radiative/noncompact alternative and the rough-core alternative.
\end{remark}

\clearpage
\bibliographystyle{alpha}
\bibliography{type_II_regularity}

\end{document}
